% Produire le plan détaillé d’une dissertation — Français 1re Tech — Cours signé OnBuch+
% Programme officiel MINESEC. Compiler avec : tectonic N02-produire-plan-detaille-dissertation.tex
\newcommand{\DOCMATIERE}{Français}
\newcommand{\DOCNIVEAU}{1re Tech}
\newcommand{\DOCMODULE}{Français – classe de Première technique}
\newcommand{\DOCLECON}{N02}
\newcommand{\DOCTITRE}{Produire le plan détaillé d’une dissertation}
% OnBuch+ — préambule « cours signés » ENRICHI (compilé avec tectonic / XeLaTeX)
% Système visuel « Pop » d'OnBuch+ : fond crème (jamais blanc), bordures encre
% épaisses, ombres « dures » décalées, titres Archivo Black en capitales.
% Version MATHÉMATIQUES : graphes (pgfplots/tikz), boîtes pédagogiques
% (définition, propriété, méthode, attention, le savais-tu, exercice résolu,
% exercice, TP, bilan), page de garde et signature OnBuch+ ancrée en bas.
%
% Macros attendues dans main.tex AVANT \input{preamble} :
%   \DOCMATIERE  \DOCNIVEAU  \DOCMODULE  \DOCLECON (numéro)  \DOCTITRE
\documentclass[11pt]{article}

\usepackage{fontspec}
\usepackage[french]{babel}
\usepackage[a4paper,margin=2cm,top=3.4cm,bottom=2.8cm,headheight=1.4cm,footskip=1.3cm]{geometry}
\usepackage{amsmath,amssymb,amsfonts}
\usepackage{mathtools} % \xmapsto, \coloneqq... souvent utilisés par les modèles
\usepackage{cancel} % simplifications barrées (bilans de forces)
% \abs{z} (module, valeur absolue) et \norm{u} (norme), utilisés par les modèles
\providecommand{\abs}{}\let\abs\relax\DeclarePairedDelimiter{\abs}{\lvert}{\rvert}
\providecommand{\norm}{}\let\norm\relax\DeclarePairedDelimiter{\norm}{\lVert}{\rVert}
\usepackage[table]{xcolor} % [table] : fournit aussi \rowcolors (alternance de lignes)
\usepackage{pagecolor}
\usepackage{tikz}
\usetikzlibrary{arrows.meta,positioning,calc,shapes.geometric,shapes.misc,shapes.arrows,decorations.pathmorphing,decorations.pathreplacing,decorations.markings,patterns,shadows,mindmap,trees,fit,backgrounds,matrix,intersections,angles,quotes}
\usepackage{pgfplots}
\usepackage[version=4]{mhchem} % équations nucléaires \ce{^{235}_{92}U}
\usepackage[european,siunitx]{circuitikz} % schémas électriques
\pgfplotsset{compat=1.18}
\usepgfplotslibrary{fillbetween,groupplots} % aires entre courbes (calcul intégral)
\usepackage{enumitem}
\usepackage{fancyhdr}
% [most] charge toutes les bibliothèques tcolorbox (listings, minted, raster,
% documentation...) et gonfle inutilement la mémoire TeX sur un document long
% (~300 boîtes) : on ne charge que ce qui est réellement utilisé.
\usepackage[skins,breakable]{tcolorbox}
\usepackage{graphicx}
\usepackage{colortbl}
\usepackage{booktabs}
\usepackage{tabularx}
\usepackage{multirow}
\usepackage{diagbox} % case d'en-tête diagonale des tableaux à double entrée
% Colonnes extensibles centrée / gauche / droite (C, L, R), souvent utilisées par les modèles
\newcolumntype{C}{>{\centering\arraybackslash}X}
\newcolumntype{L}{>{\raggedright\arraybackslash}X}
\newcolumntype{R}{>{\raggedleft\arraybackslash}X}
\usepackage{array}
\usepackage{multicol}
\usepackage{siunitx}
% parse-numbers=false : accepte aussi \SI{2,5 \times 10^{-3}}{...}
\sisetup{locale=FR,output-decimal-marker={,},group-minimum-digits=4,parse-numbers=false}
\DeclareSIUnit\ohm{\ensuremath{\symup{\Omega}}} % U+2126 absent de la police
\DeclareSIUnit\division{div} % réglages d'oscilloscope (ms/div, V/div)
\DeclareSIUnit\tr{tr} % tours (vitesse de rotation, ex. \SI{1500}{\tr\per\minute})
\DeclareSIUnit\molar{M} % concentration molaire (ex. \SI{3}{\molar}, \SI{1}{\micro\molar})
\DeclareSIUnit\an{an} % année (le modèle confond parfois avec \year, un compteur TeX) — ex. \SI{5}{\kilo\meter\per\an}
\DeclareSIUnit\FCFA{FCFA} % franc CFA (ex. \SI{500}{\FCFA\per\kilo\gram})
\DeclareSIUnit\degreeBrix{\textdegree Brix} % teneur en sucre d'un sirop/jus (réfractométrie)
\DeclareSIUnit\month{mois} % le modèle utilise parfois \month, un compteur TeX, comme unité de durée
\DeclareSIUnit\microliter{\micro\liter} % le modèle écrit parfois \microliter en un seul token
\DeclareSIUnit\million{millions} % dénombrement (ex. \SI{15}{\million\per\mL} spermatozoïdes)
\usepackage{qrcode}
% \degree : commande fréquemment utilisée par les modèles pour un angle ou une
% température en dehors de \SI{}{} (non fournie par les paquets chargés).
\providecommand{\degree}{^{\circ}}
% \qed (fin de démonstration), utilisé par les modèles sans amsthm
\providecommand{\qed}{\hfill$\square$}
% \barre : double barre des valeurs interdites dans un tableau de variations (tkz-tab)
\providecommand{\barre}{\ensuremath{\|}}
% environnement proof (amsthm non chargé) utilisé par les modèles
\ifdefined\proof\else\newenvironment{proof}[1][Démonstration]{\par\noindent\textit{#1.}\ }{\qed\par}\fi
\usepackage{lastpage}
\usepackage{hyperref}
\hypersetup{hidelinks}

% ============================================================
% Palette « Pop » OnBuch+ — mêmes hex que la classe P (plus_ui.dart)
% ============================================================
\definecolor{poporange}{HTML}{F59321}
\definecolor{popdark}{HTML}{E07A0C}
\definecolor{popcream}{HTML}{FFF6E8}
\definecolor{popink}{HTML}{1C1714}
\definecolor{poppurple}{HTML}{7A5AE0}
\definecolor{popgreen}{HTML}{2E9E6A}
\definecolor{poppink}{HTML}{E0457B}
\definecolor{popblue}{HTML}{2F7DE1}
\definecolor{popgold}{HTML}{C98A12}
\definecolor{popmuted}{HTML}{8A7F76}
\definecolor{popline}{HTML}{EADFD2}
\definecolor{popgray}{HTML}{8A7F76}
\colorlet{popgrey}{popgray}
% Alias tolérants (le modèle nomme parfois les couleurs différemment)
\colorlet{popwhite}{white}
\colorlet{popblack}{popink}
\colorlet{popred}{poppink}
\colorlet{popyellow}{popgold}
% Teintes claires (fonds de tableaux/encadrés) — le modèle nomme parfois
% les couleurs avec un suffixe L (ex. popblueL) sans les avoir définies.
\colorlet{poporangeL}{poporange!20}
\colorlet{popdarkL}{popdark!20}
\colorlet{poppurpleL}{poppurple!20}
\colorlet{popgreenL}{popgreen!20}
\colorlet{poppinkL}{poppink!20}
\colorlet{popblueL}{popblue!20}
\colorlet{popgoldL}{popgold!20}
\colorlet{popredL}{popred!20}
\colorlet{popgrayL}{popgray!20}
% Teintes claires pour fonds de tableaux / zones
\colorlet{popblueL}{popblue!10!white}
\colorlet{poporangeL}{poporange!14!white}
\colorlet{popgreenL}{popgreen!12!white}
\colorlet{poppurpleL}{poppurple!12!white}
\colorlet{poppinkL}{poppink!12!white}
\colorlet{popgoldL}{popgold!14!white}

\pagecolor{popcream}

% ============================================================
% Typographie
% ============================================================
\defaultfontfeatures{Ligatures=TeX}
\setmainfont{PlusJakartaSans-Medium.ttf}[Path=../../../fonts/,BoldFont=PlusJakartaSans-Bold.ttf]
% Maths en sans-serif (Fira Math) avec chiffres et lettres droites de Plus
% Jakarta, pour que les formules se fondent dans le texte.
\usepackage[math-style=ISO,bold-style=ISO]{unicode-math}
\setmathfont{FiraMath-Regular.otf}[Path=../../../fonts/]
\setmathfont{PlusJakartaSans-Medium.ttf}[Path=../../../fonts/,range={up/num,up/{Latin,latin},\mathup}]
\usepackage{icomma} % virgule décimale sans espace en mode math
% Caractères mathématiques Unicode absents des polices de texte (Plus Jakarta,
% Archivo Black) : rendus par leur équivalent mathématique, en texte comme en
% formule — sinon ils s'affichent en carré vide (ex. « Arithmétique dans ℤ »).
\usepackage{newunicodechar}
\newunicodechar{ℝ}{\ensuremath{\mathbb{R}}}
\newunicodechar{ℤ}{\ensuremath{\mathbb{Z}}}
\newunicodechar{ℂ}{\ensuremath{\mathbb{C}}}
\newunicodechar{ℕ}{\ensuremath{\mathbb{N}}}
\newunicodechar{ℚ}{\ensuremath{\mathbb{Q}}}
\newunicodechar{𝔻}{\ensuremath{\mathbb{D}}}
\newunicodechar{α}{\ensuremath{\alpha}}
\newunicodechar{β}{\ensuremath{\beta}}
\newunicodechar{γ}{\ensuremath{\gamma}}
\newunicodechar{δ}{\ensuremath{\delta}}
\newunicodechar{ε}{\ensuremath{\varepsilon}}
\newunicodechar{θ}{\ensuremath{\theta}}
\newunicodechar{λ}{\ensuremath{\lambda}}
\newunicodechar{μ}{\ensuremath{\mu}}
\newunicodechar{σ}{\ensuremath{\sigma}}
\newunicodechar{τ}{\ensuremath{\tau}}
\newunicodechar{φ}{\ensuremath{\varphi}}
\newunicodechar{ψ}{\ensuremath{\psi}}
\newunicodechar{ω}{\ensuremath{\omega}}
\newunicodechar{Σ}{\ensuremath{\Sigma}}
% majuscules produites par \MakeUppercase dans les titres (α-aminés -> Α)
\newunicodechar{Α}{\ensuremath{\alpha}}
\newunicodechar{Β}{\ensuremath{\beta}}
\newunicodechar{Γ}{\ensuremath{\Gamma}}
\newunicodechar{Δ}{\ensuremath{\Delta}}
\newunicodechar{Ε}{\ensuremath{\varepsilon}}
\newunicodechar{Θ}{\ensuremath{\Theta}}
\newunicodechar{Λ}{\ensuremath{\Lambda}}
\newunicodechar{Μ}{\ensuremath{\mu}}
\newunicodechar{Τ}{\ensuremath{\tau}}
\newunicodechar{Φ}{\ensuremath{\Phi}}
\newunicodechar{Ψ}{\ensuremath{\Psi}}
\newunicodechar{Ω}{\ensuremath{\Omega}}
\newunicodechar{↦}{\ensuremath{\mapsto}}
\newunicodechar{∈}{\ensuremath{\in}}
\newunicodechar{∉}{\ensuremath{\notin}}
\newunicodechar{∧}{\ensuremath{\wedge}}
\newunicodechar{⇔}{\ensuremath{\Leftrightarrow}}
\newunicodechar{⟺}{\ensuremath{\Longleftrightarrow}}
\newunicodechar{⇒}{\ensuremath{\Rightarrow}}
\newunicodechar{⟹}{\ensuremath{\Longrightarrow}}
\newunicodechar{∘}{\ensuremath{\circ}}
\newunicodechar{∪}{\ensuremath{\cup}}
\newunicodechar{∩}{\ensuremath{\cap}}
\newunicodechar{≡}{\ensuremath{\equiv}}
\newunicodechar{⊂}{\ensuremath{\subset}}
\newunicodechar{⊥}{\ensuremath{\perp}}
\newunicodechar{∃}{\ensuremath{\exists}}
\newunicodechar{∀}{\ensuremath{\forall}}
\newunicodechar{∛}{\ensuremath{\sqrt[3]{\,}}}
\newunicodechar{∜}{\ensuremath{\sqrt[4]{\,}}}
\newunicodechar{⃗}{\ensuremath{{}^{\to}}}
\newunicodechar{ⁿ}{\ifmmode^{n}\else\textsuperscript{\ensuremath{n}}\fi}
\newunicodechar{ˣ}{\ifmmode^{x}\else\textsuperscript{\ensuremath{x}}\fi}
\newunicodechar{ᵇ}{\ifmmode^{b}\else\textsuperscript{\ensuremath{b}}\fi}
\newunicodechar{ʳ}{\ifmmode^{r}\else\textsuperscript{\ensuremath{r}}\fi}
\newunicodechar{ᵘ}{\ifmmode^{u}\else\textsuperscript{\ensuremath{u}}\fi}
\newunicodechar{ʸ}{\ifmmode^{y}\else\textsuperscript{\ensuremath{y}}\fi}
\newunicodechar{ᵃ}{\ifmmode^{a}\else\textsuperscript{\ensuremath{a}}\fi}
\newunicodechar{ᵉ}{\ifmmode^{e}\else\textsuperscript{\ensuremath{e}}\fi}
\newunicodechar{ᵏ}{\ifmmode^{k}\else\textsuperscript{\ensuremath{k}}\fi}
\newunicodechar{ᵖ}{\ifmmode^{p}\else\textsuperscript{\ensuremath{p}}\fi}
\newunicodechar{ᴺ}{\ifmmode^{N}\else\textsuperscript{\ensuremath{N}}\fi}
\newunicodechar{ᶜ}{\ifmmode^{c}\else\textsuperscript{\ensuremath{c}}\fi}
\newunicodechar{ʰ}{\ifmmode^{h}\else\textsuperscript{\ensuremath{h}}\fi}
\newunicodechar{ᵠ}{\ifmmode^{q}\else\textsuperscript{\ensuremath{q}}\fi}
\newunicodechar{ₙ}{\ifmmode_{n}\else\textsubscript{\ensuremath{n}}\fi}
\newunicodechar{ₐ}{\ifmmode_{a}\else\textsubscript{\ensuremath{a}}\fi}
\newunicodechar{ᵢ}{\ifmmode_{i}\else\textsubscript{\ensuremath{i}}\fi}
\newunicodechar{ᵣ}{\ifmmode_{r}\else\textsubscript{\ensuremath{r}}\fi}
\newunicodechar{ₖ}{\ifmmode_{k}\else\textsubscript{\ensuremath{k}}\fi}
\newunicodechar{ᵦ}{\ifmmode_{\beta}\else\textsubscript{\ensuremath{\beta}}\fi}
\newunicodechar{ₓ}{\ifmmode_{x}\else\textsubscript{\ensuremath{x}}\fi}
\newfontfamily\popdisplay{ArchivoBlack-Regular.ttf}[Path=../../../fonts/]
\newfontfamily\poplogo{SpaceGrotesk-Bold.ttf}[Path=../../../fonts/]
\newfontfamily\popsemi{PlusJakartaSans-SemiBold.ttf}[Path=../../../fonts/]
\newfontfamily\popxbold{PlusJakartaSans-ExtraBold.ttf}[Path=../../../fonts/]
\newfontfamily\popxboldls{PlusJakartaSans-ExtraBold.ttf}[Path=../../../fonts/,LetterSpace=3.0]

\setlist[enumerate]{leftmargin=1.7em,itemsep=5pt,topsep=4pt}
\setlist[itemize]{leftmargin=1.7em,itemsep=4pt,topsep=4pt,label={\color{poporange}\textbullet}}
\setlength{\parindent}{0pt}
\setlength{\parskip}{5pt}
\renewcommand{\arraystretch}{1.25}

% pgfplots : style Pop par défaut
\pgfplotsset{
  every axis/.append style={axis line style={popink,line width=1pt},tick style={popink},
    grid style={popline,line width=0.5pt},label style={font=\small},tick label style={font=\footnotesize}},
  popcurve/.style={poporange,line width=1.8pt,no markers,smooth},
}
% Même style disponible dans un \draw TikZ simple (hors pgfplots)
\tikzset{popcurve/.style={poporange,line width=1.8pt}}
\tikzset{popgrid/.style={popline,very thin}}
\colorlet{popcurve}{poporange} % le nom du style est parfois utilisé comme couleur

% ============================================================
% Pill / squiggle
% ============================================================
\newcommand{\poppill}[3]{%
  \tikz[baseline=-0.55ex]\node[draw=popink,line width=1.2pt,rounded corners=2.6pt,%
    fill=#2,inner xsep=6.5pt,inner ysep=3pt]{%
    \popxboldls\fontsize{7}{7}\selectfont\color{#3}\MakeUppercase{#1}};%
}
\newcommand{\popsquiggle}[1]{%
  \tikz[baseline=-0.4ex]{
    \draw[#1,line width=2.6pt,line cap=round]
      plot [smooth] coordinates {(0,0.09) (0.34,0.01) (0.68,0.21) (1.02,0.01) (1.36,0.10)};
  }%
}
% Mot-clé surligné (marqueur orange sous le texte)
\newcommand{\cle}[1]{{\popxbold\color{popdark}#1}}

% ============================================================
% En-tête / pied
% ============================================================
\pagestyle{fancy}
\fancyhf{}
\renewcommand{\headrulewidth}{0pt}
\renewcommand{\footrulewidth}{0pt}
\lhead{\raisebox{-0.15ex}{\poplogo\fontsize{13}{13}\selectfont\color{popink}OnBuch\color{poporange}+}}
\rhead{\poppill{\DOCMATIERE\ \textbullet\ \DOCNIVEAU\ \textbullet\ Leçon \DOCLECON}{white}{popink}}
\lfoot{\popxbold\fontsize{7.6}{7.6}\selectfont\color{popmuted}\MakeUppercase{Cours signé OnBuch+}\ \textbullet\ {\popsemi\DOCTITRE}}
\rfoot{\tikz[baseline=-0.6ex]\node[rounded corners=8.5pt,fill=popink,minimum height=17pt,minimum width=17pt,inner xsep=5pt,inner sep=0]{\popxbold\fontsize{8}{8}\selectfont\color{popcream}\thepage\,/\,\pageref*{LastPage}};}

% ============================================================
% Titres
% ============================================================
\newcommand{\chaptitre}[2]{%
  \begin{tcolorbox}[enhanced,colback=poporange,colframe=popink,boxrule=2.6pt,arc=11pt,
      drop shadow=popink,left=15pt,right=15pt,top=12pt,bottom=11pt,before skip=3pt,after skip=18pt]
    \poppill{#2}{popink}{popcream}\par\vspace{9pt}
    {\raggedright\hyphenpenalty=10000\exhyphenpenalty=10000\popdisplay\fontsize{22}{24}\selectfont\color{popcream}\MakeUppercase{#1}\par}
  \end{tcolorbox}
}
% Section numérotée : pastille orange avec le numéro + titre Archivo
\newcounter{popsec}
\newcommand{\coursec}[1]{%
  \stepcounter{popsec}\setcounter{popsub}{0}%
  \par\vspace{16pt}\noindent
  \begin{minipage}{\linewidth}
  \tikz[baseline=-0.7ex]\node[draw=popink,line width=1.6pt,fill=poporange,rounded corners=4pt,minimum size=22pt,inner sep=0]{\popdisplay\fontsize{12}{12}\selectfont\color{popcream}\thepopsec};%
  \hspace{8pt}{\popdisplay\fontsize{15}{17}\selectfont\color{popink}\MakeUppercase{#1}}\par
  \vspace{4pt}\noindent{\color{poporange}\rule{36pt}{3.6pt}}
  \end{minipage}\par\nopagebreak\vspace{8pt}
}
\newcounter{popsub}
\newcommand{\courssub}[1]{%
  \stepcounter{popsub}%
  \par\vspace{9pt}\noindent{\popxbold\fontsize{11.5}{13}\selectfont\color{popink}{\color{poporange}\thepopsec.\thepopsub}\hspace{6pt}#1}\par\nopagebreak\vspace{4pt}
}

% ============================================================
% Boîtes pédagogiques — même recette : fond blanc, bordure épaisse colorée,
% ombre dure encre, badge plein. Titre optionnel [#1].
% ============================================================
\tcbset{popbase/.style={breakable,enhanced,colback=white,boxrule=2.3pt,arc=9pt,
  shadow={3pt}{-3pt}{0pt}{popink},left=14pt,right=14pt,top=11pt,bottom=11pt,before skip=11pt,after skip=15pt,
  fontupper=\popsemi\fontsize{10.6}{14.5}\selectfont\color{popink},
  fontlower=\popsemi\fontsize{10.6}{14.5}\selectfont\color{popink}}}
% Coche dessinée (le glyphe ✓ n'existe pas dans les polices du document)
\newcommand{\popcheck}[1]{\tikz[baseline=-0.5ex]\draw[#1,line width=1.6pt,line cap=round,line join=round](-0.13,0.0)--(-0.04,-0.09)--(0.13,0.1);}
\providecommand{\coche}{\popcheck{popgreen}}
\providecommand{\resultat}[1]{\par\smallskip\noindent\fcolorbox{popgreen}{popgreenL}{\parbox{\dimexpr\linewidth-2\fboxsep-2\fboxrule\relax}{\textbf{Résultat.} #1}}\par\smallskip}
\newcommand{\popboxhead}[3]{\poppill{#1}{#2}{white}\ifx\relax#3\relax\else\hspace{6pt}{\popxbold\fontsize{10.6}{12}\selectfont\color{popink}#3}\fi\par\vspace{7pt}}

\newtcolorbox{aretenir}[1][]{popbase,colframe=popgreen,before upper={\popboxhead{À retenir}{popgreen}{#1}}}
\newtcolorbox{exemplebox}[1][]{popbase,colframe=poporange,before upper={\popboxhead{Exemple}{poporange}{#1}}}
\newtcolorbox{definition}[1][]{popbase,colframe=popblue,before upper={\popboxhead{Définition}{popblue}{#1}}}
\newtcolorbox{propriete}[1][]{popbase,colframe=poppurple,before upper={\popboxhead{Propriété}{poppurple}{#1}}}
\newtcolorbox{methode}[1][]{popbase,colframe=popgold,before upper={\popboxhead{Méthode}{popgold}{#1}}}
\newtcolorbox{attention}[1][]{popbase,colframe=poppink,before upper={\popboxhead{Attention}{poppink}{#1}}}
\newtcolorbox{experience}[1][]{popbase,colframe=popblue,colback=popblueL,before upper={\popboxhead{Expérience}{popblue}{#1}}}
% « Le savais-tu ? » : carte violette pleine (contexte, culture, Cameroun)
\newtcolorbox{savaistu}[1][]{popbase,colback=poppurple,colframe=popink,
  fontupper=\popsemi\fontsize{10.4}{14.2}\selectfont\color{white},
  before upper={\poppill{Le savais-tu\,?}{white}{poppurple}\ifx\relax#1\relax\else\hspace{6pt}{\popxbold\color{white}#1}\fi\par\vspace{7pt}}}
% Exercice résolu : énoncé en haut, \tcblower puis la solution sur fond crème
\newcounter{popexo}\newcounter{popexr}
\newtcolorbox{exoresolu}[1][]{popbase,colframe=poporange,colbacklower=popcream,
  segmentation style={popink,line width=1.2pt,dashed},
  before upper={\stepcounter{popexr}\popboxhead{Exercice résolu \thepopexr}{poporange}{#1}},
  before lower={\poppill{Solution}{popink}{popcream}\par\vspace{6pt}}}
% Exercice (non résolu en place) — la correction est dans la partie Corrigés
\newtcolorbox{exercice}[1][]{popbase,colframe=popink,
  before upper={\stepcounter{popexo}\popboxhead{Exercice \thepopexo}{popink}{#1}}}
% Corrigé
\newtcolorbox{corrige}[1][]{popbase,colframe=popgreen,colback=popgreenL,
  before upper={\popboxhead{Corrigé}{popgreen}{#1}}}
% Objectifs / prérequis / bilan
\newtcolorbox{objectifs}{popbase,colframe=popink,colback=poporangeL,
  before upper={\popboxhead{Objectifs de la leçon}{popink}{}}}
\newtcolorbox{prerequis}{popbase,colframe=popink,colback=popblueL,
  before upper={\popboxhead{Prérequis}{popblue}{}}}
\newtcolorbox{bilan}{popbase,colframe=popink,colback=white,boxrule=2.8pt,
  before upper={\popboxhead{Fiche bilan}{poporange}{L'essentiel à connaître pour l'examen}}}
% Figure encadrée avec légende
\newtcolorbox{popfigure}[1][]{enhanced,breakable=false,colback=white,colframe=popline,boxrule=1.4pt,arc=8pt,
  left=10pt,right=10pt,top=10pt,bottom=8pt,before skip=10pt,after skip=12pt,halign=center,
  lower separated=false,halign lower=center,
  fontlower=\popsemi\fontsize{9}{11.5}\selectfont\color{popmuted},
  #1}
\newcommand{\legende}[1]{\par\vspace{6pt}{\popsemi\fontsize{9}{11.5}\selectfont\color{popmuted}#1}}

% ============================================================
% Signature OnBuch+
% ============================================================
\newcommand{\poplogoinline}[1]{{\poplogo\fontsize{#1}{#1}\selectfont\color{popink}OnBuch\color{poporange}+}}
\newcommand{\signatureonbuch}{%
  \begin{tcolorbox}[enhanced,colback=popink,colframe=popink,boxrule=2.2pt,arc=10pt,
      drop shadow=poporange,left=14pt,right=14pt,top=10pt,bottom=10pt,before skip=0pt,after skip=0pt]
    \tikz[baseline=-0.7ex]\node[circle,fill=popgreen,draw=popcream,line width=1.3pt,minimum size=22pt,inner sep=0]{\popcheck{white}};%
    \hspace{9pt}\begin{minipage}[c]{0.62\linewidth}
      {\popxbold\fontsize{12}{13}\selectfont\color{popcream}Signé\ \textbullet\ {\poplogo OnBuch\color{poporange}+}}\par\vspace{2pt}
      {\raggedright\popsemi\fontsize{8}{10.5}\selectfont\color{popline}\MakeUppercase{Équipe pédagogique OnBuch+}\\\MakeUppercase{Conforme au programme officiel MINESEC}\par}
    \end{minipage}\hfill
    {\poplogo\fontsize{22}{22}\selectfont\color{popcream}OnBuch\color{poporange}+}
  \end{tcolorbox}%
}

% ============================================================
% Page de garde
% #1 titre · #2 matière · #3 niveau · #4 module · #5 n° leçon · #6 durée ·
% #7 description / objectifs (texte ou itemize)
% ============================================================
\newcommand{\pagegarde}[7]{%
  \thispagestyle{empty}
  \newgeometry{a4paper,margin=1.7cm,top=1.6cm,bottom=1.4cm}%
  \begin{tikzpicture}[remember picture,overlay]
    \fill[poporange!18] ([xshift=-3.2cm,yshift=-3.6cm]current page.north east) circle (4.6cm);
    \fill[poppurple!14] ([xshift=2.4cm,yshift=7.6cm]current page.south west) circle (3.8cm);
  \end{tikzpicture}%
  {\poplogo\fontsize{30}{30}\selectfont\color{popink}OnBuch\color{poporange}+}\hfill
  \raisebox{0.6ex}{\poppill{Cours signé \textbullet\ Premium}{popink}{popcream}}\par
  \vspace{1pt}\popsquiggle{poppurple}\par
  \vspace{8pt}
  \begin{tcolorbox}[enhanced,colback=poporange,colframe=popink,boxrule=2.8pt,arc=13pt,
      drop shadow=popink,left=18pt,right=18pt,top=12pt,bottom=13pt,after skip=10pt]
    \poppill{#2 \textbullet\ #3}{popink}{popcream}\hspace{5pt}\poppill{#4}{white}{popink}\par\vspace{8pt}
    {\popdisplay\fontsize{14}{14}\selectfont\color{popink}LEÇON #5}\par\vspace{4pt}
    {\raggedright\hyphenpenalty=10000\exhyphenpenalty=10000\popdisplay\fontsize{25}{27}\selectfont\color{popcream}\MakeUppercase{#1}\par}
    \vspace{8pt}
    {\popxbold\fontsize{9.5}{9.5}\selectfont\color{popink}\MakeUppercase{Durée conseillée : #6}}
  \end{tcolorbox}
  \begin{tcolorbox}[enhanced,colback=white,colframe=popink,boxrule=2.2pt,arc=10pt,
      drop shadow=popink,left=16pt,right=16pt,top=10pt,bottom=10pt,after skip=8pt]
    \poppill{Dans cette leçon}{poppurple}{white}\par\vspace{5pt}
    \popsemi\fontsize{9.3}{12.6}\selectfont\color{popink}#7
  \end{tcolorbox}
  \noindent\begin{minipage}[t]{0.487\linewidth}
    \begin{tcolorbox}[enhanced,colback=popgreen,colframe=popink,boxrule=2.2pt,arc=10pt,
        drop shadow=popink,left=11pt,right=11pt,top=10pt,bottom=10pt,height=3.1cm,valign=center]
      \tcbox[colback=white,colframe=popink,boxrule=1.3pt,arc=5pt,boxsep=0pt,nobeforeafter,
        left=3pt,right=3pt,top=3pt,bottom=3pt,valign=center]{\qrcode[height=1.9cm]{https://play.google.com/store/apps/details?id=com.luvvix.onbuch&hl=fr}}%
      \hspace{8pt}\begin{minipage}[c]{0.5\linewidth}
        {\popdisplay\fontsize{11}{12.5}\selectfont\color{white}\MakeUppercase{Télécharge OnBuch}}\par\vspace{4pt}
        {\raggedright\popsemi\fontsize{8.4}{11}\selectfont\color{white}Gratuit \textbullet\ Google Play\\Tuteur IA, annales, concours, résultats\par}
      \end{minipage}
    \end{tcolorbox}
  \end{minipage}\hfill
  \begin{minipage}[t]{0.487\linewidth}
    \begin{tcolorbox}[enhanced,colback=poppurple,colframe=popink,boxrule=2.2pt,arc=10pt,
        drop shadow=popink,left=11pt,right=11pt,top=10pt,bottom=10pt,height=3.1cm,valign=center]
      \tikz[baseline=-0.7ex]\node[circle,fill=white,draw=popink,line width=1.3pt,minimum size=1.6cm,inner sep=0]{\popdisplay\fontsize{24}{24}\selectfont\color{poppurple}+};%
      \hspace{8pt}\begin{minipage}[c]{0.6\linewidth}
        {\popdisplay\fontsize{11}{12.5}\selectfont\color{white}\MakeUppercase{Passe à OnBuch+}}\par\vspace{4pt}
        {\raggedright\popsemi\fontsize{8.4}{11}\selectfont\color{white}Tous les cours signés, Tuteur IA illimité, fiches PDF — onglet OnBuch+ de l'app\par}
      \end{minipage}
    \end{tcolorbox}
  \end{minipage}
  \vspace{8pt}
  \popcontenu
  \vfill
  \signatureonbuch
  \restoregeometry
  \newpage
}

% Rangée « Ce cours contient » (page de garde)
\newcommand{\poptile}[3]{% #1 couleur #2 gros texte #3 légende
  \begin{tcolorbox}[enhanced,colback=#1,colframe=popink,boxrule=2pt,arc=9pt,shadow={2.5pt}{-2.5pt}{0pt}{popink},
      left=6pt,right=6pt,top=9pt,bottom=9pt,halign=center,valign=center,height=1.75cm,width=\linewidth,nobeforeafter]
    {\popdisplay\fontsize{12.5}{13}\selectfont\color{white}\MakeUppercase{#2}}\par\vspace{3pt}
    {\centering\popsemi\fontsize{7.8}{9.5}\selectfont\color{white}#3\par}
  \end{tcolorbox}}
\newcommand{\popcontenu}{%
  \par\noindent{\popxboldls\fontsize{7.5}{7.5}\selectfont\color{popmuted}CE COURS SIGNÉ CONTIENT}\par\vspace{7pt}
  \noindent\begin{minipage}[t]{0.235\linewidth}\poptile{popblue}{Cours}{complet, illustré, pas à pas}\end{minipage}\hfill
  \begin{minipage}[t]{0.235\linewidth}\poptile{poporange}{Méthodes}{et exercices résolus}\end{minipage}\hfill
  \begin{minipage}[t]{0.235\linewidth}\poptile{poppink}{Exercices}{type examen + corrigés}\end{minipage}\hfill
  \begin{minipage}[t]{0.235\linewidth}\poptile{popgreen}{Fiche bilan}{l'essentiel à retenir}\end{minipage}\par}

% Sommaire
\newtcolorbox{sommaire}{enhanced,colback=white,colframe=popink,boxrule=2.2pt,arc=10pt,
  drop shadow=popink,left=18pt,right=18pt,top=13pt,bottom=13pt,before skip=6pt,after skip=20pt,
  before upper={\poppill{Sommaire}{poppurple}{white}\par\vspace{9pt}\popsemi\fontsize{10.6}{14.5}\selectfont\color{popink}}}

% Signature de fin de cours — poussée tout en bas de la dernière page
\newcommand{\findecours}{%
  \par\vfill\nopagebreak
  \begin{center}\popsquiggle{poporange}\end{center}\vspace{4pt}
  \signatureonbuch
}
\AtBeginDocument{\def\llbracket{\mathopen{[\mkern-3.5mu[}}\def\rrbracket{\mathclose{]\mkern-3.5mu]}}} % intervalles d'entiers (glyphe ⟦ absent de la police)
% Glyphes absents de Fira Math : redéfinis à partir de glyphes disponibles
\AtBeginDocument{%
  \def\setminus{\mathbin{\backslash}}\let\smallsetminus\setminus
  \def\vdots{\mathord{\vbox{\baselineskip3.2pt\lineskiplimit0pt\kern4pt\hbox{.}\hbox{.}\hbox{.}}}}%
  \def\ddots{\mathinner{\mkern1mu\raise6.5pt\hbox{.}\mkern2mu\raise3.6pt\hbox{.}\mkern2mu\raise0.7pt\hbox{.}\mkern1mu}}%
  \def\triangle{\mathord{\scalebox{0.9}{$\symup{\Delta}$}}}%
  \def\nabla{\mathord{\raisebox{\depth}{\scalebox{1}[-1]{$\symup{\Delta}$}}}}%
  \def\checkmark{\popcheck{popdark}}%
  \def\star{\mathbin{\ast}}\def\bigstar{\mathord{\ast}}%
  \def\diamond{\mathbin{\scalebox{0.6}{$\Diamond$}}}%
  \def\bigcup{\mathop{\mathchoice{\raisebox{-0.3em}{\scalebox{1.7}{$\cup$}}}{\scalebox{1.25}{$\cup$}}{\cup}{\cup}}\displaylimits}%
  \def\bigcap{\mathop{\mathchoice{\raisebox{-0.3em}{\scalebox{1.7}{$\cap$}}}{\scalebox{1.25}{$\cap$}}{\cap}{\cap}}\displaylimits}%
  \def\lesseqqgtr{\mathrel{\lessgtr}}\def\bigoplus{\mathop{\oplus}\displaylimits}%
}
\providecommand{\ssi}{si et seulement si} % abréviation fréquente
\AtBeginDocument{\providecommand{\wideparen}{\overparen}} % arc de cercle

%% FIGFIX-BEGIN v5 — correctifs globaux des figures (docs/onbuch-generation/tools/figfix)
\usepackage{adjustbox}\usepackage{etoolbox}\tcbuselibrary{raster}
% toute figure TikZ/pgfplots plus large que la ligne est réduite à la largeur disponible
\BeforeBeginEnvironment{tikzpicture}{\begin{adjustbox}{max width=\linewidth}}
\AfterEndEnvironment{tikzpicture}{\end{adjustbox}}
% légendes pgfplots lisibles par-dessus les courbes
\pgfplotsset{every axis/.append style={legend style={fill=white,fill opacity=0.92,text opacity=1,draw=black!15,inner sep=3pt}}}
% étiquettes d'arêtes (syntaxe edge["..."]) avec fond blanc
\tikzset{every edge quotes/.append style={fill=white,inner sep=1.2pt}}
%% FIGFIX-END
\begin{document}

\pagegarde{Produire le plan détaillé d’une dissertation}{Français}{1re Tech}{Français – classe de Première technique}{N02}{25 séances (fiche de progression harmonisée)}{Ce module premium couvre l'intégralité de la fiche de progression harmonisée MINESEC 2026-27 pour la 1re Tech (25 séances). Il articule l'étude de l'œuvre *Le Lion et la Perle*, la maîtrise de la langue (ponctuation, syntaxe), la rhétorique argumentative et la méthodologie complète de la dissertation jusqu'à l'élaboration du plan détaillé.
\vspace{4pt}
\begin{multicols}{2}\small
\begin{itemize}
\item À la fin de ce module, je sais analyser les paratextes et lire méthodiquement *Le Lion et la Perle* pour en rendre un compte rendu structuré.
\item Je sais identifier les situations et conditions de communication pour adapter mon discours.
\item Je maîtrise l'usage de la ponctuation forte et faible ainsi que la distinction phrase simple / phrase composée pour écrire correctement.
\item Je reconnais les caractéristiques et les stratégies du texte argumentatif (thèses, arguments, connecteurs, figures de style).
\item J'applique la méthode en 4 étapes pour analyser un sujet de dissertation et formuler une problématique pertinente.
\item Je mobilise les techniques de recherche d'idées (brainstorming, QQOQCP, analogies) pour alimenter ma réflexion.
\end{itemize}\end{multicols}}

\begin{sommaire}
\begin{multicols}{2}
\begin{enumerate}
\item 1. Entrée dans l'œuvre : *Le Lion et la Perle* de Wole Soyinka
\item 2. Communication \& Langue : Situations, Ponctuation \& Syntaxe
\item 3. Rhétorique : Le Texte Argumentatif - Caractéristiques \& Stratégies
\item 4. Méthodologie de la Dissertation : Du Sujet à la Problématique
\item 5. Recherche d'Idées \& Élaboration du Plan Détaillé
\item 6. Intégration, Évaluation \& Remédiation
\item Activité d'intégration
\item Méthodes et exercices résolus
\item Exercices
\item Corrigés des exercices
\item Fiche bilan
\end{enumerate}
\end{multicols}
\end{sommaire}

\begin{objectifs}
\begin{itemize}
\item À la fin de ce module, je sais analyser les paratextes et lire méthodiquement *Le Lion et la Perle* pour en rendre un compte rendu structuré.
\item Je sais identifier les situations et conditions de communication pour adapter mon discours.
\item Je maîtrise l'usage de la ponctuation forte et faible ainsi que la distinction phrase simple / phrase composée pour écrire correctement.
\item Je reconnais les caractéristiques et les stratégies du texte argumentatif (thèses, arguments, connecteurs, figures de style).
\item J'applique la méthode en 4 étapes pour analyser un sujet de dissertation et formuler une problématique pertinente.
\item Je mobilise les techniques de recherche d'idées (brainstorming, QQOQCP, analogies) pour alimenter ma réflexion.
\item Je construis un plan détaillé (dialectique, thématique, analytique) équilibré, cohérent et conforme aux attendus du Baccalauréat technique.
\item Je réalise une tâche d'intégration complexe : du sujet brut au plan détaillé abouti, en autonomie.
\end{itemize}
\end{objectifs}

\begin{prerequis}
\begin{itemize}
\item Notion de paratexte et lecture cursive (programme de 4ème / 3ème).
\item Les fonctions de la phrase (sujet, verbe, COD, COI, CC) et classes grammaticales de base.
\item Les connecteurs logiques de base (addition, opposition, cause, conséquence).
\item Structure de base d'un paragraphe argumentatif (idée directrice, arguments, exemple).
\item Culture générale sur le Nigeria colonial et l'œuvre de Wole Soyinka (contexte d'étude).
\end{itemize}
\end{prerequis}

\newpage

\coursec{Entrée dans l'œuvre : \emph{Le Lion et la Perle} de Wole Soyinka}

Au carrefour de Mvog-Ada à Yaoundé, un élève de Première technique hésite devant un sujet de dissertation au Probatoire : « La tradition est-elle un frein au progrès ? ». Il possède des idées mais ne sait ni les ordonner ni les articuler en un plan rigoureux. Comme Lakunle et Baroka dans \emph{Le Lion et la Perle}, il oscille entre modernité mal maîtrisée et tradition mal comprise. Ce module lui donne les outils pour transformer le chaos des idées en architecture logique, condition de la réussite à l'examen et de la citoyenneté active.

\textbf{Problématique de la section :} comment entrer dans une œuvre théâtrale africaine, comprendre son univers (personnages, conflit, culture yoruba) et en rendre compte par écrit, afin de disposer d'un corpus solide pour nourrir ensuite une dissertation ?

\courssub{Lire les paratextes : poser des hypothèses de lecture}

Avant même d'ouvrir la pièce, le lecteur attentif commence par observer tout ce qui \emph{entoure} le texte. C'est une démarche d'enquête : comme le technicien qui examine la plaque signalétique d'une machine avant de la démonter, le lecteur recueille des indices pour anticiper le sens de l'œuvre.

\begin{definition}[Le paratexte]
Le \cle{paratexte} est l'ensemble des éléments qui entourent le texte principal et guident la lecture : titre, sous-titre, nom de l'auteur, couverture (image, couleurs), quatrième de couverture, préface, notes éditoriales, dédicace. La notion a été théorisée par Gérard Genette (\emph{Seuils}, 1987). On distingue le \textbf{péritexte} (éléments \emph{dans} le livre : titre, préface, notes) et l'\textbf{épitexte} (éléments \emph{hors} du livre : interviews, articles, biographie de l'auteur).
\end{definition}

Appliquons cette grille à \emph{Le Lion et la Perle} (\emph{The Lion and the Jewel}, 1959) :

\begin{center}
\begin{tabularx}{\linewidth}{|l|X|X|}
\hline
\rowcolor{popblueL} \textbf{Paratexte} & \textbf{Observation} & \textbf{Hypothèse de lecture} \\
\hline
Titre & « Le Lion \emph{et} la Perle » : deux êtres unis par « et » & Un duel ou une union entre deux personnages symboliques : la force (lion) contre la beauté (perle) \\
\hline
Auteur & Wole Soyinka, dramaturge nigérian & Une pièce enracinée dans la culture yoruba, probablement critique envers la colonisation \\
\hline
Couverture & Souvent un masque africain ou une danseuse & Théâtre, danse, oralité : une esthétique africaine assumée \\
\hline
Genre & « Pièce en trois parties : Matinée, Midi, Soirée » & Action concentrée sur une seule journée (unité de temps) \\
\hline
Préface / notes & Vocabulaire yoruba, chansons traduites & Le texte exigera un effort de compréhension culturelle \\
\hline
\end{tabularx}
\end{center}

\begin{savaistu}[Wole Soyinka, Nobel de littérature]
Wole Soyinka, né en 1934 au Nigeria, est le premier Africain à recevoir le prix Nobel de littérature (1986). Il a écrit \emph{Le Lion et la Perle} en anglais en 1959, à la veille de l'indépendance du Nigeria (1960). La pièce mêle l'héritage du théâtre classique (unité de temps, unité d'action, rôle collectif du village proche du chœur antique) et la tradition yoruba (proverbes, chants, masques, mime dansé). Soyinka fut aussi un opposant politique : emprisonné pendant la guerre du Biafra (1967-1970), il incarne l'écrivain engagé.
\end{savaistu}

\begin{popfigure}
\begin{center}
\begin{tikzpicture}[mindmap, grow cyclic, every node/.style=concept, concept color=popblue!40, level 1/.style={level distance=4.6cm, sibling angle=90}, level 2/.style={level distance=2.8cm, sibling angle=45}]
\node {\emph{Le Lion et la Perle}}
child { node {Paratextes} child { node {Titre} } child { node {Auteur} } }
child { node {Contexte} child { node {Nigeria 1960} } child { node {Colonisation} } }
child { node {Personnages} child { node {Sidi (Perle)} } child { node {Baroka (Lion)} } child { node {Lakunle (Instituteur)} } }
child { node {Thèmes} child { node {Tradition/Modernité} } child { node {Pouvoir/Séduction} } child { node {Oralité/Théâtre} } };
\end{tikzpicture}
\end{center}
\legende{Carte mentale d'entrée dans l'œuvre : quatre portes d'entrée (paratextes, contexte, personnages, thèmes) pour organiser la lecture.}
\end{popfigure}

\courssub{Lecture ouvroir : découvrir l'univers de la pièce}

La \cle{lecture ouvroir} (lecture « porte d'entrée », première lecture globale et rapide) vise à identifier trois repères : le cadre spatio-temporel, les personnages, le conflit central.

\textbf{Le cadre.} L'action se déroule à \textbf{Ilujinle}, village imaginaire du Nigeria, à l'époque coloniale (vers la fin des années 1950). Le village vit au rythme de la tradition yoruba, mais les signes de la modernité pénètrent : l'école, le magazine photographique, le projet de route, le timbre postal.

\textbf{Les personnages principaux.}

\begin{center}
\begin{tabularx}{\linewidth}{|l|X|X|}
\hline
\rowcolor{popblueL} \textbf{Personnage} & \textbf{Identité} & \textbf{Ce qu'il représente} \\
\hline
Sidi & La plus belle jeune fille du village, la « Perle », rendue vaniteuse par ses photos publiées dans un magazine & La beauté, l'orgueil, la féminité convoitée \\
\hline
Lakunle & Jeune instituteur, habillé « à l'européenne », qui refuse de payer le prix de la fiancée & Une modernité superficielle, mimétique et ridicule \\
\hline
Baroka & Le Bale (chef) du village, 62 ans, le « Lion », rusé et séducteur & La tradition vivante, le pouvoir, la ruse — mais aussi une vraie capacité d'adaptation \\
\hline
Sadiku & Première épouse de Baroka, messagère et complice & Les femmes du palais, gardiennes des rites \\
\hline
\end{tabularx}
\end{center}

\textbf{Le conflit central.} Sidi est courtisée par deux hommes : Lakunle veut l'épouser « à la moderne » (sans dot, avec des baisers « comme à Londres »), Baroka veut en faire sa nouvelle épouse selon la coutume. Derrière ce triangle amoureux se joue le vrai débat de la pièce : \textbf{tradition contre modernité}. Mais Soyinka brouille les cartes : le « moderne » Lakunle est un perroquet ridicule, tandis que le « traditionnel » Baroka comprend le monde moderne (il a saboté le projet de route qui menaçait son village, mais sait utiliser le timbre postal comme symbole de son pouvoir).

\begin{attention}[Piège de lecture n°1]
Erreur fréquente : confondre les deux prétendants. Lakunle incarne une modernité \emph{superficielle} (vocabulaire anglais mal digéré, mépris des coutumes qu'il ne comprend pas) ; Baroka incarne une tradition \emph{vivante}, rusée, et paradoxalement plus « moderne » dans l'âme car il sait évoluer. Ne réduisez jamais Baroka à un « vieux polygame » : c'est un stratège qui feint l'impuissance pour piéger Sidi. À l'inverse, ne voyez pas en Lakunle un héros du progrès : Soyinka le traite en pantin comique.
\end{attention}

\courssub{Lecture méthodique n°1 : la « Matinée »}

La première partie, « Matinée » (\emph{Morning}), installe tout le dispositif dramatique. Observons-la comme un artisan observe une pièce mécanique : pièce par pièce.

\textbf{Premier moment : Sidi et Lakunle.} Sidi, un seau d'eau sur la tête, croise Lakunle qui la sermonne : une femme ne devrait pas porter de charges lourdes, devrait s'habiller « décemment », et il refuse de payer la dot car c'est « une coutume barbare ». Sidi le remet à sa place avec esprit : sans dot payée, elle passerait pour une femme de peu de valeur. On voit ici le \textbf{langage de Soyinka} : dialogues vifs, humour, et surtout les \textbf{proverbes} yoruba qui donnent à la parole une sagesse collective.

\textbf{Deuxième moment : la célébrité de Sidi et l'apparition de Baroka.} Arrive la nouvelle : un étranger photographe a publié les photos de Sidi dans un magazine ; le Bale lui-même n'y occupe qu'une petite place, à côté des latrines du village. Sidi devient une célébrité locale. Puis le village célèbre l'événement par un \textbf{chant et une danse} : c'est le \textbf{théâtre dans le théâtre} — les villageois miment la venue du photographe et rejouent sa maladresse. Baroka apparaît, feint l'indifférence, mais décide déjà de conquérir Sidi.

\begin{aretenir}[Ce que la « Matinée » met en place]
\begin{itemize}
\item Le \textbf{conflit des valeurs} : la dot (tradition) contre le mariage « moderne » ;
\item Le \textbf{ressort de la vanité} : les photos transforment Sidi, la perle devient orgueilleuse ;
\item Le \textbf{style de Soyinka} : proverbes, chants, danse, mime (théâtre dans le théâtre), ironie dramatique ;
\item La \textbf{position de Baroka} : observateur silencieux qui prépare sa ruse.
\end{itemize}
\end{aretenir}

\begin{popfigure}
\begin{center}
\begin{tikzpicture}[scale=1, every node/.style={font=\small}]
\draw[-{Stealth}, thick, popdark] (0,0) -- (12.5,0) node[below left]{Une seule journée};
\foreach \x/\t/\d in {1.5/{Matinée}/{Sidi devient célèbre ; Baroka la remarque}, 6/{Midi}/{Sadiku porte la demande ; Baroka feint l'impuissance}, 10.5/{Soirée}/{Sidi tombe dans le piège et épouse Baroka}}{
\fill[poporange] (\x,0) circle (0.12);
\draw[popdark] (\x,0) -- (\x,0.5);
\node[above, align=center, text width=3.6cm] at (\x,0.6) {\textbf{\t}\\ \d};
}
\end{tikzpicture}
\end{center}
\legende{Structure tripartite de \emph{Le Lion et la Perle} : trois parties correspondant aux trois moments d'une même journée (unité de temps), au service d'un seul enjeu : la conquête de Sidi (unité d'action).}
\end{popfigure}

\begin{exemplebox}[Un exemple de construction dramatique resserrée]
La pièce compte \textbf{3 parties} (Matinée, Midi, Soirée), soit l'\textbf{unité de temps} classique : toute l'action tient en \emph{une seule journée}. Cette concentration temporelle sert l'\textbf{unité d'action} : un seul enjeu, la conquête de Sidi. En une journée : le matin, Sidi devient célèbre et Baroka la remarque ; à midi, Baroka envoie Sadiku demander Sidi en mariage, puis répand le mensonge de sa propre impuissance ; le soir, Sidi, venue se moquer du « vieux lion impuissant », tombe dans le piège et devient son épouse. La ruse du Lion triomphe de l'orgueil de la Perle.
\end{exemplebox}

\courssub{Contrôle de lecture et compte rendu critique}

Après la lecture, l'élève doit vérifier sa compréhension (contrôle de lecture : questions sur les personnages, l'ordre des événements, le sens des proverbes) puis produire un \textbf{compte rendu critique}. Ce n'est ni un simple résumé, ni une opinion vague : c'est un texte structuré qui informe, analyse et juge.

\begin{methode}[Rédiger un compte rendu critique en 4 étapes]
\begin{enumerate}
\item \textbf{Fiche d'identité de l'œuvre} : auteur, titre, date de publication, genre littéraire, contexte (Nigeria, fin de la colonisation).
\item \textbf{Résumé synthétique} : raconter le début (les photos, la célébrité de Sidi), le milieu (la demande en mariage, le mensonge de Baroka), la fin (le piège, le mariage de Sidi et Baroka) en 8 à 12 lignes, au présent de narration.
\item \textbf{Analyse} : dégager les personnages (leurs caractères, leurs symboles), les thèmes (féminité, pouvoir, ruse, tradition et modernité) et le style (proverbes, chants, humour, ironie dramatique).
\item \textbf{Appréciation critique justifiée} : donner son jugement personnel en l'appuyant sur des arguments précis et des références au texte (jamais « j'ai aimé » sans dire \emph{pourquoi}).
\end{enumerate}
\end{methode}

\begin{exoresolu}[Question de contrôle de lecture]
Énoncé : Pourquoi Sidi refuse-t-elle d'épouser Lakunle sans le paiement de la dot ? Que révèle ce refus de sa conception de la valeur de la femme ?
\tcblower
\textbf{Solution détaillée.} Sidi refuse le mariage sans dot parce que, dans la coutume d'Ilujinle, le prix de la fiancée atteste publiquement la valeur de la femme et l'engagement sérieux de l'époux. Épouser Lakunle « gratuitement » la ferait passer, aux yeux du village, pour une femme sans valeur, voire déshonorée. Ce refus révèle que Sidi, loin d'être une simple beauté vaniteuse, intériorise les codes de sa communauté : sa dignité passe par la reconnaissance sociale. Soyinka montre ainsi que la « modernité » de Lakunle, qui méprise la dot sans proposer de dignité de remplacement, est une fausse libération : elle rabaisse la femme au lieu de l'élever. La réponse mobilise deux thèmes majeurs du compte rendu : la féminité et le conflit tradition/modernité.
\end{exoresolu}

\courssub{Négocier le projet d'étude de l'œuvre}

En classe de Première technique, l'étude d'une œuvre se fait en \textbf{projet négocié} : l'enseignant et les élèves choisissent ensemble les axes d'approfondissement, en fonction des intérêts de la classe et des compétences à construire (ici : alimenter la dissertation). Trois axes s'imposent pour \emph{Le Lion et la Perle} :

\begin{enumerate}
\item \textbf{La condition féminine} : Sidi est-elle un trophée ou un sujet libre ? Que disent Sadiku et les épouses du palais de la place des femmes dans la société yoruba ?
\item \textbf{La satire de la modernité factice} : Lakunle, avec son costume trop chaud et ses mots anglais mal employés, incarne le décalquage ridicule de l'Occident. Cet axe prépare directement les sujets de dissertation sur la tradition et le progrès.
\item \textbf{L'oralité et le théâtre} : proverbes, chants, danses, mime, théâtre dans le théâtre — comment Soyinka fait-il entrer la tradition orale africaine dans une pièce écrite en anglais ?
\end{enumerate}

La négociation aboutit à un contrat de classe : répartition des lectures, calendrier des comptes rendus, choix des extraits pour les lectures méthodiques suivantes.

\courssub{Remédiation ciblée : surmonter les difficultés de compréhension}

La première lecture révèle des obstacles récurrents. La remédiation les traite un par un :

\begin{center}
\begin{tabularx}{\linewidth}{|l|X|X|}
\hline
\rowcolor{popblueL} \textbf{Difficulté} & \textbf{Exemple} & \textbf{Stratégie de remédiation} \\
\hline
Vocabulaire yoruba & \emph{Bale} (chef de village), noms des rites et des masques & Construire un glossaire collectif de la classe, enrichi à chaque lecture \\
\hline
Références culturelles & La dot, la polygamie, le statut du Bale & Replacer chaque coutume dans son contexte sans jugement hâtif ; comparer avec des réalités camerounaises proches \\
\hline
Ironie dramatique & Le spectateur sait que Baroka n'est pas impuissant ; Sidi, elle, l'ignore & Repérer qui sait quoi sur scène : l'écart entre le savoir du public et celui du personnage crée le comique et le suspense \\
\hline
Langue et proverbes & Les images proverbiales semblent obscures & Traduire le proverbe en idée abstraite, puis le rapprocher d'un proverbe français ou camerounais équivalent \\
\hline
\end{tabularx}
\end{center}

\begin{attention}[Piège de lecture n°2]
Ne lisez pas la pièce comme un roman : un texte théâtral est fait pour être \emph{vu} et \emph{entendu}. Les didascalies (indications scéniques : gestes, chants, danses) font partie de l'œuvre au même titre que les répliques. Les ignorer, c'est perdre la moitié du sens — en particulier les scènes de mime, où le « théâtre dans le théâtre » commente l'action principale.
\end{attention}

\begin{exercice}[Pour vérifier l'acquisition de la section]
\begin{enumerate}
\item Définissez le paratexte et citez trois éléments paratextuels de \emph{Le Lion et la Perle} avec l'hypothèse de lecture que chacun permet de formuler.
\item Complétez : le conflit central oppose \ldots\ à \ldots\ à travers les personnages de \ldots\ et \ldots, autour de \ldots.
\item Expliquez en cinq lignes pourquoi la structure en trois parties (Matinée, Midi, Soirée) illustre l'unité de temps.
\item Rédigez la fiche d'identité complète de l'œuvre (auteur, titre, date, genre, contexte).
\end{enumerate}
\end{exercice}

\begin{corrige}[Exercice « Pour vérifier l'acquisition de la section »]
\begin{enumerate}
\item Le paratexte est l'ensemble des éléments qui entourent le texte et guident la lecture (Genette). Exemples : le titre (duel ou union entre la force et la beauté), le nom de l'auteur (Soyinka : théâtre yoruba, regard critique sur la colonisation), la mention du genre et des trois parties (action concentrée sur une journée).
\item Le conflit central oppose la \textbf{tradition} à la \textbf{modernité} à travers les personnages de \textbf{Baroka} et \textbf{Lakunle}, autour de \textbf{Sidi} (la conquête de la Perle).
\item Les trois parties correspondent aux trois moments d'une même journée (matin, midi, soir) : toute l'action se déroule en moins de vingt-quatre heures, ce qui réalise l'unité de temps héritée du théâtre classique et renforce l'unité d'action (la conquête de Sidi).
\item Auteur : Wole Soyinka (Nigeria, Nobel 1986) ; titre : \emph{Le Lion et la Perle} (\emph{The Lion and the Jewel}) ; date : 1959 ; genre : comédie dramatique (théâtre) ; contexte : village yoruba d'Ilujinle, Nigeria colonial, à la veille des indépendances.
\end{enumerate}
\end{corrige}

\begin{aretenir}[Bilan de la section]
Entrer dans \emph{Le Lion et la Perle}, c'est : (1) interroger les paratextes pour formuler des hypothèses ; (2) identifier le cadre (Ilujinle, Nigeria colonial), les personnages (Sidi, Lakunle, Baroka, Sadiku) et le conflit (tradition contre modernité) ; (3) lire méthodiquement la « Matinée » en prêtant attention aux proverbes, chants, danses et au théâtre dans le théâtre ; (4) rendre compte par un compte rendu critique en quatre étapes ; (5) négocier les axes du projet d'étude ; (6) lever les obstacles (vocabulaire yoruba, références culturelles, ironie dramatique). Ces acquis fourniront les exemples et les arguments de la dissertation, objet des sections suivantes.
\end{aretenir}

\coursec{Communication \& Langue : Situations, Ponctuation \& Syntaxe}

Dans la section précédente, nous avons franchi le seuil de \emph{Le Lion et la Perle} de Wole Soyinka : nous avons observé ses paratextes et entamé la lecture du texte ouvroir. Or, une pièce de théâtre est avant tout un immense acte de \cle{communication} : des personnages parlent, s'écoutent, se trompent, se disputent. Pour comprendre comment Baroka le Lion séduit Sidi la Perle, il faut d'abord comprendre comment fonctionne toute communication humaine. C'est l'objet de cette section : démonter le mécanisme de la communication, puis maîtriser les deux outils de la langue qui rendent nos messages lisibles — la structure de la phrase et la ponctuation.

\courssub{Les situations de communication : le modèle de Jakobson}

\textbf{Intuition d'abord.} Imaginez deux scènes de la vie quotidienne au Cameroun. Scène A : un élève de Première fait un exposé oral devant sa classe à Yaoundé. Scène B : une mère envoie un SMS à son fils étudiant à Douala : « Envoie le reçu des frais ce soir ». Dans les deux cas, quelqu'un parle, quelqu'un écoute, un contenu circule. Pourtant, tout est différent : la voix contre le téléphone, le français soutenu contre le style abrégé, la salle de classe contre la distance. Le linguiste Roman Jakobson a montré que toute communication, quelle qu'elle soit, repose sur les mêmes six éléments.

\begin{definition}[Les six éléments de la communication (Jakobson)]
Toute situation de communication met en jeu six éléments :
\begin{itemize}
\item l'\cle{émetteur} : celui qui produit le message (le locuteur, l'expéditeur) ;
\item le \cle{récepteur} : celui qui reçoit le message (l'auditeur, le destinataire) ;
\item le \cle{message} : le contenu transmis (paroles, texte, gestes) ;
\item le \cle{canal} : le moyen physique de transmission (voix, ondes radio, téléphone, papier) ;
\item le \cle{code} : le système de signes partagé (français, ewondo, anglais, langage SMS) ;
\item le \cle{contexte} : la situation, le lieu, le moment, les connaissances communes auxquels le message renvoie.
\end{itemize}
\end{definition}

\begin{popfigure}
\begin{tikzpicture}[
  >=stealth, shorten >=1pt, auto, node distance=2.5cm, thick,
  main node/.style={circle, draw, font=\sffamily\bfseries, minimum size=1.5cm}
]
\node[main node] (emetteur) {Émetteur};
\node[main node] (message) [right=of emetteur] {Message};
\node[main node] (recepteur) [right=of message] {Récepteur};
\node[main node] (contexte) [below=of message] {Contexte};
\node[main node] (code) [below left=of message] {Code};
\node[main node] (canal) [below right=of message] {Canal};
\path[->, font=\sffamily\small]
  (emetteur) edge (message)
  (message) edge (recepteur)
  (contexte) edge [bend left] (message)
  (code) edge [bend right] (message)
  (canal) edge [bend left] (message);
\end{tikzpicture}
\legende{Schéma de la communication selon Jakobson : l'émetteur adresse un message au récepteur ; le message dépend du contexte, du code et du canal.}
\end{popfigure}

\begin{exemplebox}[Identifier les éléments dans des situations camerounaises]
\begin{center}
\begin{tabularx}{\linewidth}{|l|X|X|X|X|}
\hline
\rowcolor{popblueL} \textbf{Situation} & \textbf{Émetteur / Récepteur} & \textbf{Message} & \textbf{Canal} & \textbf{Code / Contexte} \\
\hline
Exposé oral en classe & Élève / professeur et camarades & Contenu de l'exposé & Voix, tableau & Français soutenu / salle de classe, sujet imposé \\
\hline
Lettre administrative au Principal & Parent d'élève / Principal & Demande d'exonération & Papier, courrier & Français formel, formules de politesse / cadre scolaire \\
\hline
SMS à un ami & Abonné / correspondant & « wé j'arrive 2m1 » & Réseau téléphonique & Langage SMS abrégé / rendez-vous convenu \\
\hline
Débat radio (CRTV) & Journaliste et invités / auditeurs & Arguments sur un sujet d'actualité & Ondes hertziennes & Français ou pidgin / débat public en direct \\
\hline
\end{tabularx}
\end{center}
\end{exemplebox}

\begin{savaistu}[Et au théâtre ?]
Dans \emph{Le Lion et la Perle}, quand Baroka fait parvenir à Sidi sa fausse nouvelle d'impuissance, il manipule tous les éléments : il choisit un émetteur intermédiaire (Sadiku), un canal privé (la confidence), un code adapté (le langage de l'intimité féminine). Comprendre Jakobson, c'est déjà comprendre la ruse du Lion.
\end{savaistu}

\courssub{Les conditions et les obstacles de la communication}

\textbf{Observation.} Pourquoi certains messages « passent » et d'autres non ? Pourquoi un exposé peut-il échouer alors que l'élève connaît sa leçon ? L'analyse des échecs de communication révèle quatre conditions nécessaires.

\begin{aretenir}[Les quatre conditions d'une communication réussie]
\begin{enumerate}
\item \textbf{Intentionnalité} : l'émetteur doit vouloir communiquer quelque chose (informer, convaincre, émouvoir, ordonner).
\item \textbf{Partage du code} : émetteur et récepteur doivent connaître le même code. Un discours en français soutenu devant un public qui ne parle que le fulfuldé échoue.
\item \textbf{Canal fiable} : le moyen de transmission doit fonctionner. Une ligne téléphonique brouillée, une voix trop faible dans une salle bruyante détruisent le message.
\item \textbf{Contexte partagé} : le récepteur doit pouvoir rattacher le message à une situation connue. « Il est arrivé » ne veut rien dire si l'on ne sait pas de qui l'on parle.
\end{enumerate}
\end{aretenir}

\begin{definition}[Les obstacles à la communication]
\begin{itemize}
\item le \cle{bruit} : toute perturbation du canal (vacarme, mauvaise réception, rature dans une copie) ;
\item l'\cle{ambiguïté} : un message qui peut recevoir plusieurs interprétations (« J'ai vu l'homme avec les jumelles » : qui a les jumelles ?) ;
\item la \cle{barrière linguistique} : absence de code commun ou maîtrise insuffisante du code (fautes de grammaire, vocabulaire inadapté).
\end{itemize}
\end{definition}

\begin{attention}[Piège fréquent à l'examen]
Ne confondez pas \textbf{canal} et \textbf{code}. Le canal est matériel (la voix, le papier, les ondes) ; le code est le système de signes (la langue, le style SMS). Une panne de réseau est un problème de \emph{canal} ; un mot mal orthographié est un problème de \emph{code}.
\end{attention}

\courssub{Morphosyntaxe : phrase simple et phrase composée}

\textbf{Intuition.} Pourquoi dit-on qu'une phrase est « simple » ou « composée » ? Le critère n'est ni la longueur ni la difficulté : c'est le nombre de \cle{verbes conjugués}. Comptez les verbes conjugués, et vous avez la réponse.

\begin{definition}[Phrase simple]
La \cle{phrase simple} ne contient qu'\textbf{une seule forme verbale conjuguée} (un seul noyau prédicatif). Exemple : \emph{Sidi danse.} — Un seul verbe conjugué (\emph{danse}), donc une seule proposition.
\end{definition}

\begin{definition}[Phrase composée]
La \cle{phrase composée} contient \textbf{plusieurs verbes conjugués}, donc plusieurs propositions. On distingue :
\begin{itemize}
\item la \textbf{juxtaposition} : les propositions sont placées côte à côte, séparées par une virgule ou un point-virgule, sans mot de liaison. \emph{Le lion rugit, la perle brille.}
\item la \textbf{coordination} : les propositions sont reliées par un \textbf{conjonction de coordination} (mais, ou, et, donc, or, ni, car). \emph{Baroka est vieux, mais il est rusé.}
\item la \textbf{subordination} : une proposition dépend d'une autre, introduite par un \textbf{pronom relatif} (qui, que, où, dont...) ou une \textbf{conjonction de subordination} (quand, parce que, si, bien que...). \emph{Sidi refuse le mariage parce qu'elle est fière de sa beauté.}
\end{itemize}
\end{definition}

\begin{popfigure}
\begin{tikzpicture}[
  level distance=1.5cm,
  sibling distance=3.5cm,
  edge from parent/.style={draw, -stealth, thick},
  every node/.style={font=\sffamily\small, align=center, inner sep=2pt}
]
\node {Phrase Composée}
  child { node {Proposition Principale} child { node {Le lion rugit} } }
  child { node {Coordination (et)} }
  child { node {Proposition Juxtaposée} child { node {la perle brille ;} } }
  child { node {Subordonnée Circonstancielle (temps)}
    child { node {Conjonction : quand} }
    child { node {Sujet : le soleil} }
    child { node {Verbe : se couche} }
  };
\end{tikzpicture}
\legende{Arbre syntaxique d'une phrase composée : une proposition principale, une proposition coordonnée et une subordonnée circonstancielle de temps.}
\end{popfigure}

\begin{propriete}[Types de propositions subordonnées]
Dans une phrase avec subordination, on identifie :
\begin{itemize}
\item la \textbf{proposition principale} : celle dont dépendent les autres ;
\item la \textbf{subordonnée relative} : introduite par un pronom relatif (qui, que, dont, où), elle complète un nom. \emph{La photo \textbf{qui circule dans le village} a rendu Sidi célèbre.}
\item la \textbf{subordonnée conjonctive} : introduite par \emph{que} ou une conjonction de subordination, elle complète un verbe ou exprime une circonstance. \emph{Baroka prétend \textbf{qu'il est devenu impuissant}.} (complétive) ; \emph{\textbf{Quand le soleil se couche}, le village s'anime.} (circonstancielle de temps).
\end{itemize}
\end{propriete}

\begin{methode}[Analyser une phrase complexe en 4 étapes]
\begin{enumerate}
\item Souligner tous les verbes conjugués : leur nombre donne le nombre de propositions.
\item Repérer les mots de liaison (et, mais, que, qui, quand...) : ils indiquent le type de relation.
\item Identifier la proposition principale : c'est celle qui peut exister seule.
\item Classer chaque subordonnée : relative (complète un nom), conjonctive complétive (complète un verbe) ou circonstancielle (temps, cause, but, condition...).
\end{enumerate}
\end{methode}

\begin{exoresolu}[Analyse syntaxique guidée]
Énoncé : analysez la phrase suivante (nombre et nature des propositions) : \emph{Lakunle veut épouser Sidi, mais elle refuse parce qu'il ne veut pas payer la dot.}
\tcblower
Solution détaillée :
\begin{enumerate}
\item Verbes conjugués : \emph{veut}, \emph{refuse}, \emph{veut} $\rightarrow$ \textbf{3 propositions}.
\item Mots de liaison : \emph{mais} (coordination), \emph{parce que} (subordination de cause).
\item Proposition principale : \emph{Lakunle veut épouser Sidi}.
\item \emph{elle refuse} : proposition coordonnée par \emph{mais} (opposition). \emph{il ne veut pas payer la dot} : subordonnée conjonctive circonstancielle de cause, introduite par \emph{parce que}.
\end{enumerate}
Conclusion : phrase composée associant coordination et subordination.
\end{exoresolu}

\courssub{La ponctuation faible : virgule, point-virgule, deux-points}

\textbf{Intuition.} La ponctuation est la respiration du texte. La ponctuation \emph{faible} organise l'intérieur de la phrase : elle marque des pauses sans conclure.

\begin{aretenir}[Règles d'usage de la ponctuation faible]
\begin{itemize}
\item La \textbf{virgule} (,) sépare les éléments d'une énumération (\emph{Sidi est belle, fière et coquette}), isole une apposition (\emph{Baroka, le chef du village, accueille l'étranger}) ou détache une subordonnée circonstancielle en tête de phrase (\emph{Quand la photo paraît, Sidi devient célèbre}).
\item Le \textbf{point-virgule} (;) relie deux propositions liées par le sens (cause, conséquence, opposition) ou sépare les éléments d'une énumération complexe. \emph{Le Lion est vieux ; sa ruse reste redoutable.}
\item Les \textbf{deux-points} (:) annoncent une explication, une citation ou une énumération. \emph{Sidi comprend la vérité : Baroka l'a trompée.}
\end{itemize}
\end{aretenir}

\begin{attention}[Erreur fréquente : l'abus du point-virgule]
Le point-virgule ne sépare \textbf{pas} deux phrases indépendantes sans lien logique : c'est le rôle du \textbf{point}. Il relie deux propositions liées sémantiquement (cause, conséquence, opposition) ou des éléments d'énumération complexes. Écrire \emph{Il pleut ; mon frère aime le football} est fautif, car les deux faits n'ont aucun rapport.
\end{attention}

\courssub{La ponctuation forte : délimiter et faire sens}

\begin{propriete}[La ponctuation forte]
La \cle{ponctuation forte} marque la fin d'une unité syntaxique et intonative majeure : le \textbf{point} (.), le \textbf{point d'interrogation} (?), le \textbf{point d'exclamation} (!) et les \textbf{points de suspension} (...). Elle structure le discours global : chaque signe fort clôt une phrase et prépare la suivante.
\end{propriete}

Au-delà de la délimitation, la ponctuation forte produit des \cle{effets de sens} :
\begin{itemize}
\item les \textbf{points de suspension} créent le suspens, l'hésitation ou la menace implicite : \emph{Baroka sourit et dit : « Viens donc voir mes timbres, ce soir... »} — le lecteur devine le piège ;
\item le \textbf{point d'exclamation} exprime l'emphase, l'ordre ou l'indignation : \emph{Quelle beauté que la Perle d'Ilujinle !} ;
\item le \textbf{point d'interrogation} peut marquer l'ironie dans une question rhétorique : \emph{Crois-tu vraiment que le Lion soit trop vieux ?}
\end{itemize}

\begin{exemplebox}[Pourquoi cela compte au Probatoire et au Baccalauréat technique]
Dans un sujet de dissertation type, l'énoncé compte en moyenne 3 à 5 phrases, et votre copie en contiendra des centaines. Chaque phrase doit être correctement ponctuée pour garantir la lisibilité : le critère « Langue » représente 4 à 5 points sur 20. Une phrase sans ponctuation forte sur dix lignes est sanctionnée comme une faute de construction.
\end{exemplebox}

\begin{methode}[Ponctuer un texte brut en 3 étapes]
\begin{enumerate}
\item \textbf{Identifier les verbes conjugués} : ils signalent les unités prédicatives.
\item \textbf{Délimiter les phrases} : poser la ponctuation forte (point, ?, !) là où une unité de sens s'achève.
\item \textbf{À l'intérieur de chaque phrase}, repérer les groupes syntaxiques (sujets, compléments, appositions, subordonnées) pour poser la ponctuation faible (virgule, point-virgule, deux-points).
\end{enumerate}
\end{methode}

\begin{exoresolu}[Ponctuation d'un texte brut inspiré du \emph{Lion et la Perle}]
Énoncé : ponctuez le passage suivant : \emph{sidi la perle du village regarde sa photo elle sourit baroka le lion l'observe en silence il prépare son piège la jeune fille ne se doute de rien}
\tcblower
Solution détaillée :
\begin{enumerate}
\item Verbes conjugués : \emph{regarde}, \emph{sourit}, \emph{observe}, \emph{prépare}, \emph{se doute} $\rightarrow$ 5 unités prédicatives.
\item Ponctuation forte : on obtient trois phrases de sens complet.
\item Ponctuation faible : appositions à isoler par des virgules ; deux propositions liées par le sens peuvent être réunies par un point-virgule.
\end{enumerate}
Texte ponctué : \emph{Sidi, la perle du village, regarde sa photo. Elle sourit. Baroka, le lion, l'observe en silence ; il prépare son piège. La jeune fille ne se doute de rien...}
Remarque : les points de suspension finaux créent un effet de suspens, annonciateur du drame.
\end{exoresolu}

\courssub{Exercices d'application}

\begin{exercice}[Exercice 1 : situations de communication]
Pour chaque situation, identifiez les six éléments de Jakobson : (a) un chef de village annonce une réunion au tam-tam ; (b) un élève écrit une lettre de motivation pour un stage à la SABC ; (c) deux amis discutent en pidgin au marché de Bafoussam.
\end{exercice}

\begin{exercice}[Exercice 2 : transformation de phrases]
Transformez les phrases simples suivantes en phrases composées, selon la consigne entre parenthèses :
\begin{enumerate}
\item \emph{Sidi danse.} (ajoutez une proposition coordonnée avec \emph{et})
\item \emph{Baroka invite Sidi.} (ajoutez une subordonnée circonstancielle de but avec \emph{pour que})
\item \emph{La photo paraît dans le magazine.} (ajoutez une subordonnée relative portant sur \emph{photo})
\end{enumerate}
\end{exercice}

\begin{exercice}[Exercice 3 : ponctuation]
Ponctuez correctement : \emph{lakunle le maître d'école crie sidi tu es folle elle rit il ajoute je ne paierai jamais la dot jamais}
\end{exercice}

\begin{corrige}[Exercice 1]
(a) Émetteur : le chef ; récepteur : les villageois ; message : l'annonce de la réunion ; canal : le son du tam-tam ; code : le rythme codé du tam-tam et la langue locale ; contexte : la vie du village. (b) Émetteur : l'élève ; récepteur : le responsable des stages ; message : la demande de stage ; canal : le courrier écrit ; code : le français formel administratif ; contexte : la recherche d'emploi. (c) Émetteurs/récepteurs : les deux amis (rôles alternés) ; message : la conversation ; canal : la voix ; code : le pidgin ; contexte : le marché, lieu informel.
\end{corrige}

\begin{corrige}[Exercice 2]
1. \emph{Sidi danse et le village l'admire.} (coordination) — 2. \emph{Baroka invite Sidi pour qu'elle devienne sa dernière épouse.} (subordonnée de but) — 3. \emph{La photo qui paraît dans le magazine rend Sidi célèbre.} (subordonnée relative).
\end{corrige}

\begin{corrige}[Exercice 3]
\emph{Lakunle, le maître d'école, crie : « Sidi, tu es folle ! » Elle rit. Il ajoute : « Je ne paierai jamais la dot, jamais ! »} — Appositions entre virgules, deux-points avant la citation, point d'exclamation pour l'emphase, virgule avant la répétition insistante de \emph{jamais}.
\end{corrige}

\begin{aretenir}[L'essentiel de la section]
Toute communication repose sur six éléments (émetteur, récepteur, message, canal, code, contexte) et quatre conditions (intention, code partagé, canal fiable, contexte commun). Côté langue : une phrase simple n'a qu'un verbe conjugué ; une phrase composée en a plusieurs, reliés par juxtaposition, coordination ou subordination. La ponctuation faible organise l'intérieur de la phrase ; la ponctuation forte délimite les phrases et crée des effets de sens. Ces compétences seront indispensables pour rédiger, dans les sections suivantes, une dissertation claire et correcte.
\end{aretenir}

\coursec{Rhétorique : Le Texte Argumentatif - Caractéristiques \& Stratégies}

Dans la section précédente, nous avons appris à construire des phrases correctes (phrase simple, phrase composée) et à les ponctuer avec soin. Mais une phrase bien écrite ne suffit pas : il faut encore qu'elle serve un \emph{but}. Or, dans la vie quotidienne comme à l'examen, nous écrivons et parlons le plus souvent pour \cle{convaincre} : convaincre un jury, un lecteur, un père de famille, un électeur. C'est précisément l'objet de la rhétorique, l'art de bien parler pour persuader, que nous étudions maintenant. Cette section est le cœur de la leçon : sans elle, impossible de bâtir le plan détaillé d'une dissertation.

\courssub{Définition et finalités du texte argumentatif}

\begin{definition}[Texte argumentatif]
Un \cle{texte argumentatif} est un texte dont le but principal est de faire adhérer le destinataire (lecteur ou auditeur) à une \cle{thèse}, c'est-à-dire à un point de vue, une opinion, une position sur une question. Il s'oppose au texte narratif (raconter), au texte descriptif (dépeindre) et au texte explicatif (faire comprendre).
\end{definition}

Depuis Aristote, on distingue trois grandes finalités, trois « leviers » de l'argumentation :

\begin{itemize}
\item \cle{Convaincre} : s'adresser à la \textbf{raison} du destinataire. On utilise des arguments logiques, des faits, des démonstrations. C'est le registre du \emph{logos}. Exemple : « L'école technique forme des techniciens immédiatement employables ; les statistiques d'insertion professionnelle le prouvent. »
\item \cle{Persuader} : s'adresser aux \textbf{émotions} (peur, pitié, fierté, amour). C'est le registre du \emph{pathos}. Exemple : « Imaginez vos enfants grandissant sans eau potable ! »
\item \cle{Délibérer} : s'appuyer sur l'\textbf{image que l'orateur donne de lui-même}, sa crédibilité, sa sagesse, son expérience. C'est le registre de l'\emph{ethos}. Exemple : Baroka, le Lion d'Ilujinle, parle en chef expérimenté que l'on écoute parce qu'il a fait ses preuves.
\end{itemize}

\begin{popfigure}
\begin{tikzpicture}
\begin{axis}[
xbar,
width=\linewidth,
height=6cm,
ytick=data,
symbolic y coords={Ethos (Crédibilité), Pathos (Émotion), Logos (Raison)},
xlabel={Poids dans l'argumentation (estimation en \%)},
title={Les trois piliers de l'argumentation (Aristote)},
xmin=0, xmax=90,
nodes near coords,
every node near coord/.append style={font=\small}
]
\addplot[fill=popblue] coordinates {(70,{Logos (Raison)}) (60,{Pathos (Émotion)}) (50,{Ethos (Crédibilité)})};
\end{axis}
\end{tikzpicture}
\legende{Les trois registres de l'argumentation selon Aristote : le logos (raison) domine dans la dissertation scolaire, mais le pathos (émotion) et l'ethos (crédibilité de l'énonciateur) renforcent l'efficacité du discours. Les pourcentages sont indicatifs : ils montrent une hiérarchie, pas une mesure exacte.}
\end{popfigure}

\begin{attention}[Explication ou argumentation ?]
Ne confonds jamais l'\textbf{explication} et l'\textbf{argumentation}. Expliquer, c'est répondre à la question « \emph{comment} ça marche ? » sans chercher à faire changer d'avis : « La pompe à eau fonctionne grâce à un moteur. » Argumenter, c'est répondre à « \emph{pourquoi} c'est vrai, c'est bien, il faut le faire ? » : « Il faut installer une pompe à Ilujinle car l'eau du puits est insalubre. » La dissertation au Probatoire et au Baccalauréat technique est un exercice d'\textbf{argumentation}, jamais de simple explication.
\end{attention}

\courssub{La structure canonique du texte argumentatif}

Tout texte argumentatif bien construit obéit à une architecture simple, qu'il faut savoir repérer chez les auteurs et reproduire dans ses propres devoirs :

\begin{enumerate}
\item La \cle{thèse} : le point de vue défendu. Elle peut être \textbf{explicite} (clairement annoncée : « Je soutiens que la tradition n'est pas un obstacle au progrès ») ou \textbf{implicite} (suggérée, à deviner, comme dans les ruses de Baroka).
\item Les \cle{arguments} : les raisons avancées pour soutenir la thèse.
\item Les \cle{exemples} et preuves : les faits, anecdotes, citations qui illustrent et vérifient chaque argument.
\item La \cle{conclusion} : le bilan, souvent accompagné d'une \textbf{ouverture} (élargissement vers une question plus générale).
\end{enumerate}

\begin{definition}[L'argument]
Un \cle{argument} est une proposition avancée pour soutenir une thèse. Il est généralement introduit ou relié par un connecteur logique. Exemple : \emph{Baroka est redoutable \textbf{car} il parvient à ses fins sans jamais recourir à la force.} Ici, « car » signale que la seconde proposition justifie la première.
\end{definition}

On distingue classiquement quatre types d'arguments :

\begin{center}
\begin{tabularx}{\linewidth}{|l|X|X|}
\hline
\rowcolor{popblueL} \textbf{Type d'argument} & \textbf{Principe} & \textbf{Exemple (univers de Soyinka)} \\
\hline
Argument \textbf{logique} & Raisonnement, enchaînement de causes et de conséquences & Sans route, pas de commerce ; donc Ilujinle doit s'ouvrir. \\
\hline
Argument \textbf{d'autorité} & Appel à un expert, un aîné, un proverbe, une tradition & « Nos pères disent que le lion vieux connaît tous les sentiers. » \\
\hline
Argument \textbf{par analogie} & Comparaison avec une situation semblable & Un village sans eau est comme un corps sans sang. \\
\hline
Argument \textbf{de l'expérience} & Témoignage vécu, fait observé & « J'ai vu à Lagos des femmes travailler comme des hommes », affirme Lakunle. \\
\hline
\end{tabularx}
\end{center}

\begin{propriete}[L'argument d'autorité dans la culture orale]
L'argument d'autorité (citer un proverbe, un aîné, un sage, un expert) est particulièrement prisé dans la culture orale africaine. Chez Wole Soyinka, dans \emph{Le Lion et la Perle}, Baroka multiplie les proverbes yoruba : ils donnent à ses paroles le poids de la sagesse collective des ancêtres. Dans une dissertation, citer un proverbe ou un auteur reconnu renforce un argument, à condition de ne pas en abuser.
\end{propriete}

\begin{popfigure}
\begin{tikzpicture}[node distance=1.6cm, auto,
every node/.style={font=\small, align=center}]
\node (thesis) [rectangle, draw=popblue, fill=popblueL, rounded corners, align=center]{Thèse : la tradition a encore\\sa force et sa sagesse};
\node (arg1) [rectangle, draw=popgreen, fill=popgreenL, rounded corners, below of=thesis, align=center]{Arg 1 : Baroka triomphe\\par la ruse et l'expérience};
\node (arg2) [rectangle, draw=popgreen, fill=popgreenL, rounded corners, left of=arg1, xshift=-3.2cm, align=center]{Arg 2 : Sidi choisit\\la sécurité du chef};
\node (arg3) [rectangle, draw=popgreen, fill=popgreenL, rounded corners, right of=arg1, xshift=3.2cm, align=center]{Arg 3 : Lakunle échoue\\par rigidité};
\node (ex1) [ellipse, draw=poporange, fill=poporangeL, below of=arg1, align=center]{Ex : le stratagème\\de la fausse impuissance};
\node (ex2) [ellipse, draw=poporange, fill=poporangeL, below of=arg2, align=center]{Ex : Sidi épouse\\Baroka à la fin};
\node (ex3) [ellipse, draw=poporange, fill=poporangeL, below of=arg3, align=center]{Ex : refus de payer\\la dot, promesses vaines};
\node (conc) [rectangle, draw=poppink, fill=poppinkL, rounded corners, below of=ex1, yshift=-0.8cm] {Conclusion : la sagesse du Lion};
\path[->, thick, popdark]
(thesis) edge (arg1)
(thesis) edge (arg2)
(thesis) edge (arg3)
(arg1) edge (ex1)
(arg2) edge (ex2)
(arg3) edge (ex3)
(arg1) edge (conc)
(arg2) edge (conc)
(arg3) edge (conc);
\end{tikzpicture}
\legende{Schéma de la structure d'un développement argumentatif : une thèse soutenue par trois arguments, chacun illustré d'un exemple tiré du \emph{Lion et la Perle}, convergeant vers une conclusion.}
\end{popfigure}

\courssub{Les connecteurs logiques (opérateurs argumentatifs)}

Les \cle{connecteurs logiques} sont les articulations du raisonnement : ils indiquent au lecteur la relation entre les idées. Un paragraphe sans connecteur est une suite de phrases juxtaposées, illisibles et sanctionnée à l'examen.

\begin{center}
\begin{tabularx}{\linewidth}{|l|X|X|}
\hline
\rowcolor{popblueL} \textbf{Relation} & \textbf{Connecteurs} & \textbf{Exemple} \\
\hline
Addition & de plus, en outre, également, aussi, par ailleurs & Lakunle est instruit ; \emph{de plus}, il se croit courageux. \\
\hline
Opposition & cependant, mais, pourtant, en revanche, tandis que & Baroka est âgé ; \emph{cependant}, il reste redoutable. \\
\hline
Cause & car, parce que, puisque, en effet & Sidi refuse Lakunle \emph{parce qu'}il ne veut pas payer la dot. \\
\hline
Conséquence & donc, ainsi, c'est pourquoi, par conséquent, alors & Sa photo paraît dans le magazine ; \emph{donc} Sidi devient fière. \\
\hline
Concession & certes... mais, bien sûr... cependant, il est vrai que... & \emph{Certes}, le progrès est utile, \emph{mais} il ne doit pas détruire nos valeurs. \\
\hline
But & pour, afin de, de sorte que & Baroka ruse \emph{afin de} séduire Sidi. \\
\hline
\end{tabularx}
\end{center}

\begin{exemplebox}[La concession, arme du bon élève]
La structure \emph{certes... mais} est la marque d'un esprit nuancé, très valorisée au Probatoire : « \emph{Certes}, les idées de Lakunle semblent modernes ; \emph{mais} son mépris des traditions le rend ridicule aux yeux des villageois. » On accorde un point à l'adversaire pour mieux faire valoir sa propre position.
\end{exemplebox}

\courssub{Les figures de style au service de l'argumentation}

Les figures de style ne sont pas des ornements : elles rendent l'argumentation plus frappante, plus mémorable, plus émotive.

\begin{itemize}
\item L'\cle{ironie} : dire le contraire de ce qu'on pense pour ridiculiser. Lakunle, avec ses grands mots anglais et son mépris des coutumes, est la cible constante de l'ironie de Soyinka : le « moderniste » apparaît plus naïf que les villageois.
\item L'\cle{hyperbole} : exagération qui amplifie. « Sidi est plus célèbre que la lune ! »
\item La \cle{métonymie} : désigner une chose par une autre qui lui est liée. « Le Lion » pour désigner Baroka, « la Perle » pour Sidi.
\item La \cle{comparaison} (avec \emph{comme}) et la \cle{métaphore} (sans outil de comparaison) : « Baroka est un lion » (métaphore) ; « Sidi brille comme une perle » (comparaison).
\item L'\cle{énumération} : accumulation qui donne une impression d'abondance ou d'étouffement. Lakunle énumère ses projets : routes, écoles, machines, robes modernes...
\item La \cle{question rhétorique} : question qui n'attend pas de réponse, car celle-ci est évidente. « Faut-il vivre comme nos arrière-grands-pères ? » s'écrie Lakunle.
\end{itemize}

\courssub{Les stratégies argumentatives}

\begin{definition}[Argumentation directe et indirecte]
L'argumentation \cle{directe} expose clairement la thèse et la défend ouvertement (essai, éditorial, discours, dissertation). L'argumentation \cle{indirecte} fait passer le message par le détour d'une fiction : fable, conte, dialogue théâtral, ironie, apologue. \emph{Le Lion et la Perle} est une pièce de théâtre : Soyinka y défend une réflexion sur tradition et modernité \emph{sans jamais la formuler directement}.
\end{definition}

La \cle{réfutation} consiste à anticiper les objections de l'adversaire pour les démolir avant qu'il ne les formule : « On objectera que la tradition freine le progrès ; mais Baroka prouve qu'elle peut au contraire l'encadrer et le maîtriser. » C'est la stratégie reine du débat et de la dissertation dialectique (thèse / antithèse / synthèse).

\begin{savaistu}[Le théâtre dans le théâtre]
Dans \emph{Le Lion et la Perle}, la « danse du voyageur perdu », où les villageois rejouent l'arrivée du photographe étranger, est un exemple magnifique d'argumentation indirecte : Baroka y montre sa maîtrise de la modernité sans jamais en avoir l'air. Il \emph{fait voir} au lieu de \emph{dire}. C'est exactement ce que fait un bon auteur de fable : La Fontaine ne dit pas « soyez prudent », il raconte le corbeau et le renard.
\end{savaistu}

\begin{methode}[Analyser un texte argumentatif en cinq étapes]
\begin{enumerate}
\item \textbf{Identifier la thèse} : que veut démontrer l'auteur ? Est-elle explicite ou implicite ?
\item \textbf{Repérer les arguments} : souligner les mots-clés (\emph{car}, \emph{parce que}, \emph{donc}, \emph{en effet}).
\item \textbf{Classer les arguments} : logiques, d'autorité, d'analogie, de l'expérience, affectifs.
\item \textbf{Étudier les figures de style} : ironie, métaphore, question rhétorique... et dire leur effet sur le lecteur.
\item \textbf{Évaluer l'efficacité} : le texte convainc-il ? Pour quel public ? Avec quelles forces et quelles faiblesses ?
\end{enumerate}
\end{methode}

\begin{exoresolu}[Analyse d'un extrait du discours de Lakunle]
\textbf{Énoncé.} Lakunle déclare à Sidi : « Dans un an ou deux, tu auras des machines qui feront la cuisine à ta place. Nous aurons des routes, des voitures, et tu porteras des robes à la mode, comme les femmes de Lagos. Rejette ces coutumes barbares ! » Identifiez la thèse, deux arguments, deux figures de style et la stratégie argumentative.
\tcblower
\textbf{Solution détaillée.}
\begin{enumerate}
\item \emph{Thèse} (explicite) : Ilujinle doit abandonner ses traditions et adopter le mode de vie moderne occidental.
\item \emph{Arguments} : argument par analogie (« comme les femmes de Lagos » : Lagos est le modèle à imiter) ; argument logique implicite (les machines libèrent la femme des tâches pénibles, donc le progrès est un bien).
\item \emph{Figures de style} : l'énumération (« des routes, des voitures, des robes ») qui accumule les promesses et donne le vertige du progrès ; l'hyperbole méprisante (« coutumes barbares ») qui exagère pour condamner.
\item \emph{Stratégie} : argumentation directe (Lakunle annonce ouvertement sa thèse), mais Soyinka, l'auteur, utilise l'\textbf{ironie} : en faisant parler Lakunle avec tant d'excès, il le ridiculise. La stratégie de l'auteur est donc \emph{indirecte} : le spectateur comprend que le vrai progrès ne peut pas naître du mépris des traditions. Conclusion : le texte est doublement argumentatif — Lakunle persuade (mal), Soyinka convainc (bien) par la satire.
\end{enumerate}
\end{exoresolu}

\begin{attention}[Erreurs fréquentes à l'examen]
\begin{itemize}
\item Raconter l'histoire au lieu d'argumenter : rappeler l'intrigue du \emph{Lion et la Perle} n'est pas une argumentation.
\item Citer un exemple sans l'articuler à un argument par un connecteur logique.
\item Confondre argument d'autorité et argument logique : « le chef l'a dit » n'est pas une preuve rationnelle, c'est un appel à l'ethos.
\item Oublier la réfutation : une bonne dissertation anticipe au moins une objection.
\end{itemize}
\end{attention}

\begin{savaistu}[L'argumentation au quotidien au Cameroun]
Les éditoriaux de \emph{Cameroon Tribune}, les discours politiques lors des campagnes, les publicités pour les boissons ou les téléphones mobiles utilisent tous ces stratégies : slogans hyperboliques, témoignages (argument de l'expérience), vedettes (argument d'autorité), émotions (pathos). Savoir les décoder, c'est devenir un citoyen critique, capable de ne pas se laisser manipuler.
\end{savaistu}

\begin{exemplebox}[Enjeu pour l'examen]
Au Probatoire et au Baccalauréat technique, la dissertation est notée sur des critères précis : pertinence du contenu, qualité de l'argumentation, cohérence du plan, correction de la langue. La qualité argumentative (thèse claire, arguments articulés par des connecteurs, exemples précis) représente l'essentiel de la note de contenu. Un paragraphe dépourvu de connecteur logique est systématiquement pénalisé. Maîtriser les outils de cette section, c'est donc sécuriser la moitié de la note de production d'écrits.
\end{exemplebox}

\begin{exercice}[À toi de jouer]
\begin{enumerate}
\item Dans l'éditorial suivant, relevez la thèse et classez les arguments : « Le sport scolaire doit être obligatoire. En effet, il préserve la santé des élèves. De plus, le ministre de l'Éducation l'a rappelé : un esprit sain dans un corps sain. Certes, il prend du temps sur les études, mais un élève fatigué apprend mal. »
\item Construisez, sur le sujet « Faut-il interdire les téléphones portables au lycée ? », une phrase de concession (\emph{certes... mais}) et une phrase de réfutation d'une objection.
\item Relevez dans un extrait du \emph{Lion et la Perle} (scène de la danse du voyageur perdu) deux procédés d'argumentation indirecte et expliquez leur effet.
\end{enumerate}
\end{exercice}

\begin{corrige}[Exercice : éléments de correction]
\begin{enumerate}
\item \emph{Thèse} : le sport scolaire doit être obligatoire (explicite). \emph{Arguments} : logique (« il préserve la santé », introduit par « en effet ») ; d'autorité (la parole du ministre et le proverbe) ; concession-réfutation (« certes... mais ») qui désamorce l'objection du temps perdu.
\item Exemple attendu : « \emph{Certes}, le téléphone permet de joindre ses parents, \emph{mais} il détourne l'attention en classe. » Réfutation : « On dira que le portable sert à faire des recherches ; or l'usage montre qu'il sert surtout à consulter les réseaux sociaux. »
\item Le déguisement et le jeu de rôle (montrer au lieu de dire) ; la mise en scène comique de l'étranger ivre, qui ridiculise la fascination pour l'Occident : le public rit, donc adhère au point de vue de Baroka sans discours explicite.
\end{enumerate}
\end{corrige}

\begin{aretenir}[L'essentiel de la section]
\begin{itemize}
\item Argumenter, c'est défendre une \cle{thèse} par des \cle{arguments} illustrés d'\cle{exemples}, pour convaincre (logos), persuader (pathos) ou crédibiliser (ethos).
\item Les \cle{connecteurs logiques} (addition, opposition, cause, conséquence, concession) sont les articulations obligatoires du raisonnement.
\item Les \cle{figures de style} (ironie, hyperbole, question rhétorique...) renforcent la force du propos.
\item L'argumentation peut être \cle{directe} (éditorial, dissertation) ou \cle{indirecte} (théâtre, fable, ironie) ; la \cle{réfutation} anticipe les objections.
\end{itemize}
\end{aretenir}

Ces outils rhétoriques maîtrisés, nous pouvons maintenant passer à la pratique : la section suivante te montrera comment analyser un sujet de dissertation et formuler une problématique, première étape de la construction de ton propre texte argumentatif.

\coursec{Méthodologie de la Dissertation : Du Sujet à la Problématique}

Dans la section précédente, nous avons appris à reconnaître un texte argumentatif et à décoder ses stratégies : thèse, arguments, exemples, connecteurs logiques. Nous savons désormais \emph{lire} un raisonnement. Il est temps de franchir le pas décisif : \emph{construire} notre propre raisonnement. Or, toute dissertation réussie naît d'un acte fondateur : l'analyse du sujet et la formulation de la problématique. C'est le moment où l'élève cesse d'être un lecteur pour devenir un architecte. Un sujet mal compris, c'est une maison bâtie sur le sable : même bien rédigée, elle s'effondrera à la correction.

\courssub{La typologie des sujets de dissertation}

Avant toute analyse, il faut identifier à quelle \emph{famille} appartient le sujet proposé. Cette reconnaissance conditionne toute la stratégie du devoir.

\begin{definition}[Sujet à thèse imposée et sujet ouvert]
Un \cle{sujet à thèse imposée} (ou sujet à thèse unique) présente une seule position à discuter : « La tradition freine-t-elle le progrès ? » Le candidat doit prendre position par rapport à \emph{une} affirmation. Un \cle{sujet ouvert} (ou sujet à double thèse) met en regard deux positions : « Tradition et modernité : conflit ou complémentarité ? » Le candidat doit confronter les deux avant de trancher. En Première technique, on rencontre souvent des sujets professionnels ouverts : « La sécurité au travail : responsabilité partagée ? »
\end{definition}

On distingue classiquement quatre types de sujets :

\begin{center}
\begin{tabularx}{\linewidth}{|l|X|X|}
\hline
\rowcolor{popblueL} \textbf{Type de sujet} & \textbf{Caractéristique} & \textbf{Exemple} \\
\hline
Sujet à thèse unique & Une seule affirmation à discuter & « La tradition freine-t-elle le progrès ? » \\
\hline
Sujet à double thèse & Deux termes ou deux thèses opposés & « Tradition et modernité : conflit ou complémentarité ? » \\
\hline
Sujet citation & Une phrase d'auteur à commenter et discuter & « "La technique nous délivre de la nature" (Bergson). Discutez. » \\
\hline
Sujet mot-clé & Un ou deux mots à problématiser soi-même & « Le travail. » ou « Progrès. » \\
\hline
\end{tabularx}
\end{center}

\begin{attention}[Le sujet mot-clé, le plus dangereux]
Face à un sujet mot-clé (« Le travail »), le piège est de réciter tout ce qu'on sait sur le mot. C'est le hors-sujet assuré. Il faut \emph{fabriquer} soi-même la tension : le travail est-il libération ou aliénation ? Source de dignité ou d'exploitation ? Sans tension, pas de dissertation.
\end{attention}

\courssub{Étape 1 : L'analyse lexicale et sémantique}

L'intuition est simple : on ne peut pas débattre de ce qu'on n'a pas défini. Deux élèves qui discutent de « tradition » sans s'entendre sur le mot parlent chacun de leur côté, comme deux techniciens qui utiliseraient des plans différents. L'analyse lexicale consiste à dégager les \cle{mots-clés} du sujet et à les définir rigoureusement.

\begin{methode}[Analyser le lexique d'un sujet]
\begin{enumerate}
\item \textbf{Souligner} les mots porteurs de sens (noms, verbes d'action) et écarter les mots-outils.
\item \textbf{Définir chaque mot-clé} avec le dictionnaire, en notant toutes les acceptions (\cle{polysémie}).
\item \textbf{Construire le champ lexical} de chaque terme : synonymes, contraires, mots de la même famille.
\item \textbf{Choisir l'acception pertinente} dans le contexte du sujet.
\end{enumerate}
\end{methode}

Appliquons cette méthode au sujet : « La tradition freine-t-elle le progrès ? »

\begin{itemize}
\item \textbf{Tradition} : du latin \emph{tradere}, « transmettre ». Ensemble de coutumes, de savoir-faire, de valeurs transmis de génération en génération. Champ lexical : coutume, héritage, transmission, ancêtre, rite. Attention à la polysémie : la tradition peut être \emph{vivante} (transmission créatrice) ou \emph{figée} (conservatisme).
\item \textbf{Progrès} : du latin \emph{progredi}, « avancer ». Amélioration, avancée dans un domaine (technique, social, moral). Champ lexical : innovation, modernité, développement, évolution.
\item \textbf{Frein} : terme mécanique (familier aux élèves de série industrielle !) : dispositif qui ralentit ou arrête un mouvement. Par métaphore : obstacle, entrave. Le mot est \emph{négatif} et oriente déjà la réflexion.
\item \textbf{Produire} (dans d'autres sujets) : fabriquer, mais aussi « faire apparaître », « engendrer ». Polysémie à surveiller.
\end{itemize}

\begin{exemplebox}[Champ lexical du mot « frein »]
Frein mécanique $\rightarrow$ ralentir, bloquer, retenir $\rightarrow$ sens figuré : obstacle, entrave, empêchement, retard. Contraires : moteur, levier, accélérateur, stimulant. Ce simple travail fait déjà naître l'idée subversive : et si la tradition était un \emph{levier} plutôt qu'un frein ?
\end{exemplebox}

\courssub{Étape 2 : L'analyse logique et grammaticale}

Une fois les mots définis, il faut examiner comment ils sont \emph{agencés}. La grammaire de la phrase-sujet n'est jamais neutre.

\begin{methode}[Analyser la structure logique du sujet]
\begin{enumerate}
\item Repérer le \textbf{sujet grammatical} : de quoi parle-t-on exactement ? (« La tradition », pas « les traditions africaines » : la portée est générale.)
\item Repérer le \textbf{verbe} et sa valeur : « freine » est un verbe d'action au présent de vérité générale.
\item Repérer la \textbf{modalité} : interrogation directe (« ... ? »), affirmation à discuter, ou injonction (« Discutez », « Commentez », « Vous apprécierez »).
\item Mesurer la \textbf{portée} : l'énoncé est-il général (tout, toujours) ou restreint (certains, parfois, dans tel contexte) ?
\end{enumerate}
\end{methode}

Dans « La tradition freine-t-elle le progrès ? », l'article défini « la » généralise : il s'agit de la tradition en tant que telle, pas d'une coutume particulière. L'interrogation invite à une réponse nuancée : ni oui franc, ni non absolu, mais une réflexion sur les conditions et les limites.

\courssub{Étape 3 : Dégager la tension ou le paradoxe}

C'est l'étape reine, celle que les correcteurs attendent. Un sujet de dissertation n'est jamais une question d'information (« Quelle est la capitale du Cameroun ? ») : c'est une question \emph{débattue}. Pourquoi ? Parce qu'il contient une \cle{tension}, c'est-à-dire deux vérités qui semblent s'opposer.

\begin{exemplebox}[La tension du sujet « La tradition freine-t-elle le progrès ? »]
\textbf{D'un côté} : la tradition est identité, mémoire, cohésion sociale, sagesse accumulée (les savoir-faire artisanaux, la pharmacopée, la solidarité communautaire). \textbf{De l'autre} : la tradition peut être conservatisme, refus de l'innovation, poids des coutumes qui bloquent (refus de certaines techniques agricoles modernes, résistance au changement). \textbf{Le paradoxe} : ce qui nous enracine pourrait nous empêcher d'avancer. Voilà pourquoi le sujet mérite quatre heures de réflexion.
\end{exemplebox}

\begin{attention}[Pas de tension, pas de dissertation]
Si vous ne trouvez pas de débat dans le sujet, c'est que vous n'avez pas fini de l'analyser. Cherchez toujours le « MAIS » : « Il est vrai que... MAIS on peut objecter que... ». Ce « MAIS » est le moteur de votre devoir.
\end{attention}

\courssub{Étape 4 : Formuler la problématique}

La \cle{problématique} est la question directrice qui naît de la tension et qui guidera tout le devoir, de l'introduction à la conclusion. Elle transforme le sujet (donné) en problème personnel (pensé par le candidat).

\begin{definition}[La problématique]
Question unique, ouverte, issue de la tension du sujet, qui annonce l'angle de traitement du devoir. Elle se formule en une phrase interrogative, généralement introduite par « Comment », « Dans quelle mesure » ou « Pourquoi ».
\end{definition}

\begin{methode}[La règle des 3 « P »]
Une bonne problématique est :
\begin{enumerate}
\item \textbf{P}ertinente : directement liée au sujet, sans en ajouter ni en retrancher ;
\item \textbf{P}roblématisée : elle pose un véritable problème, elle ne se contente pas de paraphraser le sujet ;
\item \textbf{P}ortante : elle permet de développer deux ou trois parties argumentées.
\end{enumerate}
\end{methode}

\begin{exemplebox}[De la tension à la problématique]
Sujet : « La tradition freine-t-elle le progrès ? »\\
Tension : la tradition est identité MAIS peut être conservatisme.\\
Problématique : \emph{Dans quelle mesure la tradition, loin d'être un obstacle, peut-elle devenir le levier d'un progrès endogène ?}\\
Cette question est unique, ouverte, et annonce déjà un plan en trois temps (le frein réel, les limites de cette vision, la tradition comme levier).
\end{exemplebox}

\begin{attention}[Les deux erreurs fatales]
\begin{itemize}
\item Problématique \textbf{hors sujet} : parler de l'école générale alors que le sujet porte sur l'école technique ; parler des traditions européennes quand le contexte est africain.
\item Problématique \textbf{fermée} : une question à laquelle on répond par « oui » ou « non » (« La tradition est-elle un frein ? ») ne permet aucun développement. Préférez toujours « Comment », « Dans quelle mesure », « Pourquoi ».
\end{itemize}
\end{attention}

\begin{popfigure}
\begin{center}
\begin{tikzpicture}[node distance=1.5cm, every node/.style={font=\small}]
\node (start) [rectangle, rounded corners, draw, fill=yellow!30] {Sujet Brut};
\node (lex) [rectangle, draw, below of=start] {1. Analyse Lexicale (Mots-clés)};
\node (log) [rectangle, draw, below of=lex] {2. Analyse Logique (Structure)};
\node (tension) [diamond, draw, fill=orange!30, below of=log, aspect=2] {Tension / Paradoxe ?};
\node (pb) [rectangle, rounded corners, draw, fill=green!30, below of=tension] {3. PROBLÉMATIQUE (Question directrice)};
\node (plan) [rectangle, draw, below of=pb] {4. Plan Détaillé};
\draw[->] (start) -- (lex) -- (log) -- (tension);
\draw[->] (tension) -- node[left] {OUI} (pb);
\draw[->] (pb) -- (plan);
\end{tikzpicture}
\end{center}
\legende{La chaîne de traitement du sujet : chaque étape nourrit la suivante. Si le test « tension » échoue, on retourne à l'analyse lexicale.}
\end{popfigure}

\begin{exoresolu}[Analyse complète d'un sujet type Bac technique]
Énoncé : « L'école technique face aux défis de l'insertion professionnelle. » Analysez ce sujet et formulez une problématique.
\tcblower
\textbf{Solution détaillée.}
\textbf{1. Type de sujet} : sujet à thèse implicite, de type professionnel ouvert (fréquent en séries industrielles et tertiaires).
\textbf{2. Analyse lexicale} : \emph{école technique} (établissement formant aux métiers : mécanique, électrotechnique, gestion, comptabilité) ; \emph{défis} (difficultés à surmonter : adéquation formation-emploi, équipements, stages) ; \emph{insertion professionnelle} (entrée effective des lauréats dans l'emploi).
\textbf{3. Analyse logique} : « face à » exprime une confrontation ; la portée est générale (l'école technique en tant que système).
\textbf{4. Tension} : l'école délivre une formation souvent théorique MAIS l'entreprise exige des compétences pratiques immédiates. D'où le paradoxe du diplômé formé... et pourtant sans emploi.
\textbf{5. Problématique} : \emph{Comment l'école technique peut-elle adapter ses curricula pour garantir l'employabilité immédiate de ses lauréats ?}
\end{exoresolu}

\begin{popfigure}
\begin{center}
\begin{tikzpicture}
\node[align=center, text width=10cm, draw=popblue, rounded corners, fill=popblueL!40, inner sep=8pt] {
\textbf{Sujet 2023 (Exemple)} : « L'école technique face aux défis de l'insertion professionnelle. »\\[4pt]
\textbf{Mots-clés} : École technique, Défis, Insertion, Professionnelle.\\
\textbf{Tension} : École = formation théorique VS Entreprise = compétences pratiques.\\
\textbf{Problématique} : Comment l'école technique peut-elle adapter ses curricula pour garantir l'employabilité immédiate de ses lauréats ?
};
\end{tikzpicture}
\end{center}
\legende{Fiche récapitulative d'analyse : un modèle à reproduire au brouillon pour chaque sujet.}
\end{popfigure}

\begin{savaistu}[Le temps, votre allié caché]
Un bon analyseur de sujet gagne 15 à 20 minutes sur l'épreuve de 4 heures. C'est du temps réinvesti dans la rédaction et la relecture, donc des points gagnés sans effort supplémentaire. À l'inverse, l'élève qui rédige trop tôt découvre son hors-sujet à la deuxième partie... trop tard.
\end{savaistu}

\begin{exemplebox}[Ce que vaut l'analyse du sujet dans la note]
Sur 20 points en dissertation, l'analyse du sujet, la problématique et le plan pèsent souvent 5 à 6 points, à travers les critères « Compréhension du sujet » et « Construction du devoir ». Autrement dit : un tiers de la note se joue avant d'écrire la première ligne de l'introduction.
\end{exemplebox}

\courssub{Entraînement : dix sujets types Probatoire et Bac technique}

\begin{exercice}[Analyse et problématique]
Pour chacun des sujets suivants, indiquez le type de sujet, définissez deux mots-clés, dégagez la tension, puis formulez une problématique ouverte.
\begin{enumerate}
\item La tradition freine-t-elle le progrès ? (toutes séries)
\item Tradition et modernité : conflit ou complémentarité ? (toutes séries)
\item La sécurité au travail : une responsabilité partagée ? (séries industrielles)
\item « La technique nous délivre de la nature » (Bergson). Discutez. (séries industrielles)
\item Le numérique en classe : outil d'apprentissage ou source de distraction ? (toutes séries)
\item L'entrepreneuriat des jeunes : solution au chômage au Cameroun ? (séries tertiaires)
\item Le travail manuel a-t-il perdu sa valeur ? (séries industrielles)
\item La lecture est-elle encore utile à l'ère des écrans ? (toutes séries)
\item La qualité avant la quantité : un impératif pour l'entreprise moderne ? (séries tertiaires)
\item Le sport : simple loisir ou école de la vie ? (sciences sociales)
\end{enumerate}
\end{exercice}

\begin{corrige}[Exercice : éléments de correction]
\textbf{Sujet 1} : thèse unique ; tension identité/conservatisme ; problématique : « Dans quelle mesure la tradition peut-elle accompagner le progrès au lieu de l'entraver ? »
\textbf{Sujet 2} : double thèse ; tension rupture/continuité ; problématique : « Comment concilier l'héritage du passé et les exigences du monde moderne ? »
\textbf{Sujet 3} : sujet professionnel ouvert ; tension responsabilité de l'employeur (équipements, formation) / vigilance du salarié (port des protections, respect des consignes) ; problématique : « Comment la prévention des accidents peut-elle reposer à la fois sur l'employeur et sur le travailleur ? »
\textbf{Sujet 4} : sujet citation ; définir « délivrer » et « nature » ; tension libération/nouvelle dépendance aux machines ; problématique : « La technique nous affranchit-elle de la nature ou nous rend-elle dépendants d'elle autrement ? »
\textbf{Sujet 5} : double thèse ; tension concentration/dispersion ; problématique : « Dans quelle mesure le numérique peut-il servir l'apprentissage sans nuire à l'attention des élèves ? »
\textbf{Sujet 6} : thèse unique contextualisée ; tension initiative créatrice d'emplois / risques (financement, formation) ; problématique : « À quelles conditions l'entrepreneuriat des jeunes peut-il réduire durablement le chômage au Cameroun ? »
\textbf{Sujet 7} : thèse unique ; tension dévalorisation sociale / utilité irremplaçable des métiers manuels ; problématique : « Pourquoi le travail manuel, malgré sa dévalorisation, demeure-t-il essentiel à la société ? »
\textbf{Sujet 8} : thèse unique ; tension concurrence des écrans / irremplaçable formation de l'esprit ; problématique : « Que peut encore apporter la lecture que les écrans ne donnent pas ? »
\textbf{Sujet 9} : sujet professionnel ; tension exigence du client / rentabilité à court terme ; problématique : « Comment l'entreprise peut-elle concilier qualité de production et compétitivité ? »
\textbf{Sujet 10} : double thèse ; tension divertissement / formation (discipline, esprit d'équipe, dépassement de soi) ; problématique : « Dans quelle mesure le sport dépasse-t-il le simple loisir pour devenir un instrument d'éducation ? »
\end{corrige}

\begin{aretenir}[L'essentiel de la section]
Analyser un sujet, c'est suivre quatre étapes indissociables : définir les mots-clés (lexique), examiner la structure de l'énoncé (logique), dégager la tension qui rend le sujet débattu, puis formuler une problématique unique, ouverte et porteuse, respectant la règle des 3 « P ». Cette problématique est la boussole du devoir : c'est elle que nous transformerons, dans la section suivante, en plan détaillé.
\end{aretenir}

\coursec{Recherche d'Idées \& Élaboration du Plan Détaillé}

Dans la section précédente, nous avons appris à analyser un sujet de dissertation et à formuler une problématique claire. Mais une problématique, aussi pertinente soit-elle, reste une question sans réponse tant que l'on n'a pas rassemblé, trié et organisé les idées qui permettront d'y répondre. C'est précisément l'objet de cette section : passer de la question au \cle{plan détaillé}, c'est-à-dire à l'architecture complète du devoir, avant même de rédiger la première phrase.

\courssub{Les techniques de recherche d'idées}

\textbf{Intuition d'abord.} Imaginez un technicien qui doit construire un bâtiment : il ne commence jamais par couler le béton. Il rassemble d'abord tous les matériaux disponibles (ciment, fer, sable), puis il les trie, puis il les organise selon un plan d'architecte. La dissertation fonctionne exactement ainsi : avant de rédiger, il faut \emph{produire} des idées en quantité, sans se censurer.

\begin{definition}[Recherche d'idées]
La recherche d'idées est la phase du brouillon où le candidat mobilise toutes ses connaissances (cours, lectures, actualité, expérience personnelle) en rapport avec le sujet, \textbf{sans les juger ni les organiser encore}. Le principe est : \emph{quantité d'abord, qualité ensuite}.
\end{definition}

\textbf{a) Le brainstorming (remue-méninges) et la carte heuristique.}

Le \cle{brainstorming} consiste à noter sur le brouillon, en quelques minutes, tous les mots, exemples, auteurs et faits qui viennent à l'esprit à propos du sujet. Pour structurer ce jaillissement, on utilise la \cle{carte heuristique} (ou carte mentale) : le sujet au centre, et des branches qui partent vers les grandes idées, elles-mêmes subdivisées en exemples.

\textbf{b) La méthode QQOQCP.}

Pour ne rien oublier, on interroge le sujet avec sept questions :

\begin{center}
\begin{tabularx}{\linewidth}{|l|X|}\hline
\rowcolor{popblueL} \textbf{Question} & \textbf{Exemple avec le sujet : « La tradition est-elle un frein au progrès ? »} \\ \hline
\textbf{Qui ?} & Les gardiens des traditions (chefs, sages), les jeunes, l'État. \\ \hline
\textbf{Quoi ?} & Les rites, les langues, les savoir-faire, les techniques modernes. \\ \hline
\textbf{Où ?} & Au Cameroun (villages, villes), en Afrique, ailleurs (Japon). \\ \hline
\textbf{Quand ?} & Hier (époque précoloniale), aujourd'hui, demain. \\ \hline
\textbf{Comment ?} & Par la transmission orale, l'école, les médias. \\ \hline
\textbf{Pourquoi ?} & Pour préserver l'identité ; par peur du changement. \\ \hline
\textbf{Combien ?} & Chiffres : nombre de langues nationales, taux d'accès à l'électricité, etc. \\ \hline
\end{tabularx}
\end{center}

\textbf{c) La méthode du Pour / Contre (la balance).}

Pour tout sujet polémique, on trace deux colonnes sur le brouillon : à gauche les arguments \emph{en faveur} de la thèse implicite du sujet, à droite les arguments \emph{contre}. Cette balance prépare naturellement le plan dialectique (thèse / antithèse).

\textbf{d) Les analogies.}

On enrichit la réflexion par trois déplacements : l'\cle{analogie historique} (qu'en était-il autrefois ?), l'\cle{analogie géographique} (comment un autre pays a-t-il traité la question ? ex. : le Rwanda et les juridictions \emph{gacaca}, le Japon et la modernité enracinée), et l'\cle{analogie de domaine} (que disent la technique, la littérature, l'économie de ce problème ?).

\begin{exemplebox}[Brainstorming guidé sur un sujet technique]
Sujet : « Les énergies renouvelables peuvent-elles résoudre le problème de l'électricité au Cameroun ? » Idées brutes : barrage de Lom Pangar ; soleil abondant dans le Grand Nord ; coût élevé des panneaux solaires ; délestages fréquents ; biomasse (déchets agricoles) ; hydroélectricité = potentiel du fleuve Sanaga ; emplois verts ; dépendance aux importations de panneaux ; Lakunle dans \emph{Le Lion et la Perle} comme figure du « moderne » qui méprise les traditions ; cadre légal du secteur de l'électricité ; électrification rurale.
\end{exemplebox}

\courssub{Tri et hiérarchisation des idées}

Une fois la matière rassemblée, il faut la transformer en structure. Trois opérations successives :

\begin{enumerate}
\item \textbf{Éliminer le hors-sujet.} Toute idée qui ne répond pas à la problématique, même brillante, doit être sacrifiée. Le hors-sujet est la première cause d'échec au Probatoire et au Baccalauréat technique.
\item \textbf{Grouper par axes (thèmes).} On rassemble les idées qui vont ensemble : par exemple « arguments économiques », « arguments sociaux », « arguments techniques ». Chaque groupe deviendra une partie ou une sous-partie.
\item \textbf{Ordonner.} On classe les groupes selon un ordre justifié : du plus simple au plus complexe, du plus connu au moins connu, ordre chronologique, ou ordre logique (constat $\rightarrow$ causes $\rightarrow$ solutions).
\end{enumerate}

\begin{attention}[Le piège du hors-sujet « séduisant »]
Un exemple très documenté mais qui ne répond pas à la problématique affaiblit le devoir. Posez-vous pour chaque idée la question : « Cette idée m'aide-t-elle à répondre à LA question posée ? » Si la réponse est non, rayez-la sans regret.
\end{attention}

\courssub{Les trois grands types de plans au Baccalauréat technique}

\begin{definition}[Les trois plans de la dissertation]
\begin{itemize}
\item Le \cle{plan dialectique} : \textbf{I. Thèse / II. Antithèse / III. Synthèse}. C'est le plan \emph{standard}, adapté aux sujets philosophiques et sociétaux (« Faut-il... ? », « La tradition est-elle... ? »).
\item Le \cle{plan thématique} : on examine successivement des aspects distincts du sujet (aspect économique, social, technique...). Adapté aux sujets descriptifs ou techniques.
\item Le \cle{plan analytique} : \textbf{I. Constat/Causes / II. Conséquences / III. Solutions}. Adapté aux sujets de type « problème -- solutions ».
\end{itemize}
\end{definition}

\begin{propriete}[Choisir le bon plan]
Le plan thématique convient aux sujets techniques du type « Les énergies renouvelables au Cameroun » (I. Le solaire, II. L'hydroélectricité, III. La biomasse). Le plan dialectique convient aux sujets philosophiques ou sociétaux. Le plan analytique s'impose dès que le sujet décrit un problème à résoudre (chômage des jeunes, délestages, insécurité routière).
\end{propriete}

\begin{popfigure}
\begin{tikzpicture}[mindmap, concept color=red!40, level 1/.style={level distance=4.2cm, sibling angle=120}, level 2/.style={level distance=3cm, sibling angle=60}, every node/.style={concept, align=center, font=\small}]
\node[concept, align=center]{Plan dialectique :\\ Tradition \& Progrès}
child[concept color=blue!40] { node[concept, align=center]{I. Thèse :\\ frein} child { node[concept] {A. Conservatisme} } child { node[concept] {B. Rituels figés} } }
child[concept color=orange!40] { node[concept, align=center]{II. Antithèse :\\ moteur} child { node[concept] {A. Valeurs éthiques} } child { node[concept] {B. Savoir-faire endogènes} } }
child[concept color=green!40] { node[concept, align=center]{III. Synthèse :\\ dialectique} child { node[concept] {A. Modernité enracinée} } child { node[concept] {B. Ex : Japon, Rwanda} } };
\end{tikzpicture}
\legende{Carte heuristique d'un plan dialectique : le sujet au centre, les trois parties en branches principales, les sous-parties en branches secondaires.}
\end{popfigure}

\courssub{Construction du plan détaillé}

Le \cle{plan détaillé} est le plan complet rédigé au brouillon, avec toutes les articulations du devoir. Il se construit en quatre niveaux hiérarchiques :

\begin{center}
\begin{tabularx}{\linewidth}{|l|l|X|}\hline
\rowcolor{popblueL} \textbf{Niveau} & \textbf{Notation} & \textbf{Contenu} \\ \hline
Partie & I, II, III & Un grand axe de la réflexion (thèse, antithèse, synthèse). \\ \hline
Sous-partie & A, B & Un aspect de la partie (deux sous-parties par partie). \\ \hline
Argument & 1, 2 & Une idée directrice développée (deux par sous-partie). \\ \hline
Exemple & (ex.) & Une référence précise : fait camerounais, chiffre, loi, auteur, extrait de \emph{Le Lion et la Perle}, actualité. \\ \hline
\end{tabularx}
\end{center}

\begin{methode}[Rédiger chaque argument : le schéma I.E.E.]
Chaque argument (1, 2) du plan détaillé doit suivre le schéma en trois temps :
\begin{enumerate}
\item \textbf{Idée directrice} : la phrase qui affirme l'argument.
\item \textbf{Explication} : deux ou trois phrases qui démontrent, analysent, relient à la problématique.
\item \textbf{Exemple précis} : un chiffre, une loi, un auteur, un fait camerounais, un extrait de \emph{Le Lion et la Perle} (par exemple le conflit entre Baroka, le « Lion » gardien des traditions, et Lakunle, le maître d'école moderniste).
\end{enumerate}
Chaque argument ainsi construit correspondra à \textbf{un alinéa} de la copie.
\end{methode}

\textbf{Structure de l'introduction (au brouillon).} Elle comporte quatre étapes obligatoires : l'\cle{amorce} (une phrase générale qui situe le thème), le \cle{sujet posé} (on reformule le sujet), le \cle{sujet divisé} aboutissant à la \cle{problématique} (la question directrice), et l'\cle{annonce du plan}.

\textbf{Structure de la conclusion.} Elle comporte trois étapes : le \cle{bilan} (résumé des parties), la \cle{réponse explicite à la problématique}, et l'\cle{ouverture} (une question ou une perspective nouvelle, sans hors-sujet).

\begin{popfigure}
\begin{tikzpicture}[node distance=1cm, every node/.style={rectangle, draw=popdark, rounded corners, fill=white, very thick, align=center, font=\small}, fleche/.style={->, very thick, popblue}]
\node[fill=popblueL, align=center] (intro) at (0,3){INTRODUCTION\\ Amorce | Sujet posé | Problématique | Annonce du plan};
\node[fill=poporangeL, align=center] (I) at (-4.5,0.8){I. THÈSE\\ A. Arg 1 (ex : Cameroun)\\ B. Arg 2 (ex : œuvre)};
\node[fill=poporangeL, align=center] (II) at (0,0.8){II. ANTITHÈSE\\ A. Arg 1 (ex : technique)\\ B. Arg 2 (ex : histoire)};
\node[fill=poporangeL, align=center] (III) at (4.5,0.8){III. SYNTHÈSE\\ A. Dépassement\\ B. Modèle endogène};
\node[fill=popgreenL, align=center] (conc) at (0,-1.2){CONCLUSION\\ Bilan | Réponse à la problématique | Ouverture};
\draw[fleche] (intro) -- (I);
\draw[fleche] (intro) -- (II);
\draw[fleche] (intro) -- (III);
\draw[fleche] (I) -- (conc);
\draw[fleche] (II) -- (conc);
\draw[fleche] (III) -- (conc);
\end{tikzpicture}
\legende{Schéma fonctionnel du plan détaillé dialectique : de l'introduction aux trois parties, puis à la conclusion.}
\end{popfigure}

\begin{exemplebox}[Exemple chiffré : la densité minimale d'un plan détaillé]
Un plan détaillé noté au Baccalauréat technique doit comporter au minimum 2 parties majeures (mieux : 3), 2 sous-parties par partie et 2 arguments par sous-partie, soit $3 \times 2 \times 2 = 12$ arguments minimum. Chaque argument deviendra un alinéa de la copie : le devoir comptera donc environ 12 alinéas de développement, plus l'introduction et la conclusion.
\end{exemplebox}

\begin{attention}[L'erreur du « plan en parapluie »]
Erreur fréquente : le plan « I. Les avantages, II. Les inconvénients, III. Mon avis », sans articulation logique. Dans un vrai plan dialectique, la synthèse (III) doit \textbf{dépasser} la simple addition I + II : elle propose un point de vue nouveau qui dépasse l'opposition (par exemple : « une modernité enracinée dans les valeurs endogènes »), et non un simple « les deux ont raison ».
\end{attention}

\begin{savaistu}[Le futur de l'annonce du plan]
L'annonce du plan dans l'introduction (« Nous verrons d'abord..., ensuite..., enfin... ») s'écrit au futur ou à l'infinitif (« Examiner..., analyser..., démontrer... »). C'est un repère majeur pour le correcteur : une annonce de plan bien formulée annonce un devoir maîtrisé et oriente favorablement la notation.
\end{savaistu}

\begin{aretenir}[Boîte à outils : les mots-charnières pour articuler]
\begin{itemize}
\item Pour ouvrir : \emph{En premier lieu, D'abord, Force est de constater que}.
\item Pour ajouter : \emph{Par ailleurs, En outre, De plus, Qui plus est}.
\item Pour opposer : \emph{Cependant, Néanmoins, En revanche, Toutefois}.
\item Pour conclure : \emph{En définitive, Il en résulte que, Ainsi, En somme}.
\end{itemize}
\end{aretenir}

\courssub{Rédaction au brouillon et gestion du temps}

L'introduction et la conclusion se rédigent \textbf{intégralement au brouillon}, car ce sont les deux passages que le correcteur lit avec le plus d'attention. Le développement, lui, reste sous forme de plan détaillé (idées + exemples notés en abrégé).

\begin{methode}[Gérer les 4 heures de l'épreuve]
\begin{enumerate}
\item \textbf{45 minutes de brouillon} : analyse du sujet et problématique (10 min), recherche d'idées (10 min), tri et plan détaillé (15 min), rédaction de l'introduction et de la conclusion au brouillon (10 min).
\item \textbf{2 h 30 de rédaction} au propre, en suivant le plan sans s'en écarter.
\item \textbf{30 minutes de relecture} : orthographe, ponctuation, transitions, vérification que chaque partie répond à la problématique.
\end{enumerate}
\end{methode}

\courssub{Auto-évaluation du plan}

Avant de rédiger, le candidat contrôle son plan avec quatre critères, réunis dans la grille suivante :

\begin{center}
\begin{tabularx}{\linewidth}{|l|X|}\hline
\rowcolor{popblueL} \textbf{Critère} & \textbf{Question de contrôle} \\ \hline
\textbf{Cohérence} & Chaque partie répond-elle à la problématique ? \\ \hline
\textbf{Équilibre} & Les parties sont-elles à peu près égales en importance ? \\ \hline
\textbf{Progression} & L'ordre des parties est-il logique (simple $\rightarrow$ complexe) ? \\ \hline
\textbf{Exemples} & Chaque argument a-t-il un exemple contextualisé (Cameroun, œuvre, actualité) ? \\ \hline
\end{tabularx}
\end{center}

\begin{exoresolu}[Construire le plan détaillé d'un sujet de type « problème -- solutions »]
Énoncé. Sujet : « Le chômage des jeunes au Cameroun : causes, conséquences et solutions. » Élaborez le plan détaillé complet (parties, sous-parties, arguments, exemples).
\tcblower
Solution détaillée. Le sujet appelle un \textbf{plan analytique}.

\textbf{Introduction (brouillon)} : Amorce : la jeunesse est l'avenir d'une nation. Sujet posé : au Cameroun, une large part des jeunes diplômés est sans emploi. Problématique : quelles sont les causes du chômage des jeunes, quelles en sont les conséquences et quelles solutions envisager ? Annonce : nous examinerons d'abord les causes, puis les conséquences, enfin les solutions.

\textbf{I. Les causes du chômage des jeunes.}
A. Causes économiques : 1. faible industrialisation (ex. : faible nombre d'usines hors agroalimentaire) ; 2. raréfaction des recrutements publics.
B. Causes éducatives : 1. inadéquation formation--emploi (ex. : filières générales surchargées, filières techniques délaissées) ; 2. manque de formation professionnelle pratique.

\textbf{II. Les conséquences.}
A. Conséquences sociales : 1. paupérisation et dépendance prolongée des jeunes ; 2. émigration clandestine (ex. : départs vers l'Europe).
B. Conséquences sécuritaires : 1. délinquance urbaine ; 2. vulnérabilité aux groupes radicaux.

\textbf{III. Les solutions.}
A. Solutions étatiques : 1. promotion de l'entrepreneuriat (ex. : programmes d'appui aux jeunes entrepreneurs) ; 2. développement de l'enseignement technique.
B. Solutions individuelles et locales : 1. auto-emploi agricole et artisanal (ex. : savoir-faire endogènes valorisés, à l'image de la vitalité du village d'Ilujinle dans \emph{Le Lion et la Perle}) ; 2. coopératives de jeunes.

\textbf{Conclusion (brouillon)} : Bilan des trois parties ; réponse : le chômage des jeunes, aux causes multiples, appelle des solutions conjuguant l'État et l'initiative individuelle ; ouverture : la jeunesse saura-t-elle transformer ce défi en moteur de développement ?

Vérification par la grille : cohérence (chaque partie répond à un terme du sujet), équilibre (trois parties de poids comparable), progression (causes $\rightarrow$ conséquences $\rightarrow$ solutions), exemples contextualisés (Cameroun, œuvre au programme).
\end{exoresolu}

\begin{exercice}[À vous de jouer]
Sujet : « Les réseaux sociaux : un danger pour la jeunesse camerounaise ? » 1) Produisez une liste d'au moins dix idées par brainstorming et QQOQCP. 2) Choisissez le type de plan adapté et justifiez ce choix. 3) Rédigez le plan détaillé complet (parties, sous-parties, arguments, exemples). 4) Rédigez au brouillon l'annonce du plan et la conclusion.
\end{exercice}

\begin{corrige}[Exercice : éléments de correction]
1) Idées possibles : addiction aux écrans ; cybercriminalité et arnaques ; désinformation ; mais aussi information rapide, apprentissage en ligne, entrepreneuriat numérique, mobilisation citoyenne, coût des données, coupure des études, modèles négatifs.
2) Le sujet est interrogatif et polémique (« un danger ? ») : le \textbf{plan dialectique} s'impose (I. Oui, un danger ; II. Mais aussi une opportunité ; III. Un outil à éduquer).
3) I. A. 1. perte de temps scolaire (ex. : baisse de concentration en classe) ; 2. exposition aux contenus violents. B. 1. arnaques en ligne ; 2. désinformation. II. A. 1. accès au savoir (ex. : cours en ligne) ; 2. information en temps réel. B. 1. entrepreneuriat numérique des jeunes camerounais ; 2. lien social avec la diaspora. III. A. éducation aux médias à l'école ; B. régulation et usage responsable (dépassement : le problème n'est pas l'outil mais l'usage).
4) Annonce : « Nous verrons d'abord les dangers des réseaux sociaux, ensuite leurs apports indéniables, enfin les conditions d'un usage responsable. » Conclusion : bilan + réponse nuancée + ouverture sur la citoyenneté numérique.
\end{corrige}

\coursec{Intégration, Évaluation \& Remédiation}

Dans la section précédente, nous avons appris à rechercher des idées (carte heuristique, brainstorming) et à les organiser en un \cle{plan détaillé} équilibré. Mais un savoir-faire n'est réellement acquis que lorsqu'on sait le mobiliser \emph{seul}, sur un sujet \emph{inédit}, dans un temps limité : c'est exactement ce qu'exigera l'épreuve de dissertation au Probatoire et au Baccalauréat technique. Cette dernière section est donc le moment de la \cle{vérification} : chaque élève produit seul, est évalué, puis \emph{remédie} à ses difficultés précises. C'est aussi le moment de clôturer le projet d'étude de l'œuvre \emph{Le Lion et la Perle} par une soutenance orale.

\courssub{La tâche d'intégration finale (séances 22 à 24)}

\begin{definition}[Tâche d'intégration]
La \cle{tâche d'intégration} est une activité complexe qui demande à l'élève de mobiliser \emph{ensemble}, et en autonomie, tous les savoirs et savoir-faire de l'unité : analyse du sujet, formulation de la problématique, recherche des idées, élaboration du plan détaillé, maîtrise de la langue (ponctuation, phrase simple et composée) et stratégies argumentatives.
\end{definition}

L'intuition est simple : on ne juge pas un mécanicien en lui faisant réciter les noms des outils, mais en lui demandant de réparer un moteur. De même, on ne juge pas un candidat sur sa connaissance de la méthode, mais sur sa capacité à \emph{produire} un travail complet sur un sujet jamais vu.

\begin{popfigure}
\begin{tikzpicture}
\node[align=left, text width=12cm, draw=popgold, rounded corners, fill=popgoldL, inner sep=5pt] {\textbf{Sujet d'intégration type Bac Tech 2027} :\\ « La maîtrise de la ponctuation et de la syntaxe est-elle encore un enjeu majeur pour le technicien supérieur ? »\\ \textbf{Contraintes} : mobiliser les connaissances sur la phrase (simple/composée), la ponctuation (forte/faible), l'argumentation (stratégies) et citer un exemple technique (rapport de stage, notice, code).};
\end{tikzpicture}
\legende{Sujet inédit proposé pour la tâche d'intégration. Un sujet alternatif possible : « L'intelligence artificielle : menace ou opportunité pour la formation technique ? »}
\end{popfigure}

\begin{methode}[Déroulement de la tâche d'intégration en 5 étapes]
\begin{enumerate}
\item \textbf{Analyser le sujet} (15 min) : souligner les mots-clés, définir chaque terme, repérer la question et ses présupposés.
\item \textbf{Formuler la problématique} (10 min) : transformer le sujet en une vraie question de réflexion, ni trop large ni hors sujet.
\item \textbf{Rechercher les idées} (20 min) : produire une \cle{carte heuristique} (sujet au centre, branches d'idées, arguments et exemples).
\item \textbf{Élaborer le plan détaillé} (30 min) : choisir le type de plan adapté (dialectique, thématique ou analytique), rédiger les titres de parties et sous-parties, placer sous chaque sous-partie l'argument et l'exemple.
\item \textbf{Rédiger l'introduction et la conclusion} (30 min) : introduction en 4 temps (amorce, présentation du sujet, problématique, annonce du plan) ; conclusion en 2 temps (bilan, ouverture).
\end{enumerate}
\end{methode}

\begin{exoresolu}[Amorce de traitement du sujet d'intégration]
Énoncé. À partir du sujet « La maîtrise de la ponctuation et de la syntaxe est-elle encore un enjeu majeur pour le technicien supérieur ? », donner : (a) la problématique ; (b) deux axes d'un plan dialectique.
\tcblower
Solution détaillée.
(a) \textbf{Problématique} : « Dans un métier technique dominé par les chiffres et les machines, la maîtrise de la langue écrite reste-t-elle une compétence professionnelle décisive, ou est-elle devenue secondaire ? »
(b) \textbf{Plan dialectique} :
\begin{itemize}
\item \textbf{I. La langue écrite semble secondaire dans les métiers techniques.} A. Le technicien travaille d'abord avec des données chiffrées, des schémas, des machines (exemple : lecture d'un plan de câblage). B. Les outils numériques corrigent automatiquement l'orthographe (exemple : correcteurs intégrés aux logiciels).
\item \textbf{II. Pourtant, la maîtrise de la syntaxe et de la ponctuation reste un enjeu majeur.} A. Une ambiguïté de ponctuation peut provoquer une erreur professionnelle grave (exemple : une consigne de sécurité mal ponctuée dans une notice). B. Le rapport de stage et les écrits professionnels conditionnent l'image et la carrière du technicien (exemple : un rapport clair et bien construit valorise le stagiaire auprès de l'entreprise).
\end{itemize}
\end{exoresolu}

\courssub{L'évaluation formative : la grille critériée}

\begin{definition}[Évaluation formative]
L'\cle{évaluation formative} est une évaluation qui ne sert pas d'abord à sanctionner, mais à \emph{informer} l'élève sur ses points forts et ses lacunes, afin qu'il puisse progresser avant l'évaluation certificative (le Bac).
\end{definition}

L'évaluation se fait à l'aide d'une \cle{grille critériée} : chaque compétence est notée séparément, ce qui rend la correction juste, transparente et identique pour tous.

\begin{center}
\begin{tabularx}{\linewidth}{|l|c|X|}
\hline
\rowcolor{popblueL} \textbf{Critère} & \textbf{Barème} & \textbf{Ce qui est attendu} \\
\hline
Compréhension du sujet & /5 & mots-clés définis, aucun hors sujet \\
\hline
Problématique & /3 & question claire, ni trop large ni absente \\
\hline
Plan détaillé & /6 & type de plan cohérent, parties équilibrées, arguments + exemples \\
\hline
Langue (syntaxe, ponctuation) & /3 & phrases correctes, ponctuation forte et faible maîtrisées \\
\hline
Exemples & /3 & exemples précis, variés, liés aux arguments \\
\hline
\textbf{Total} & \textbf{/20} & \\
\hline
\end{tabularx}
\end{center}

\begin{popfigure}
\begin{tikzpicture}
\begin{axis}[ ybar, bar width=0.6cm, symbolic x coords={Analyse, Problématique, Plan, Intro/Conc, Langue, Exemples}, xtick=data, ymin=0, ymax=5, ylabel={Note /5}, title={Grille d'évaluation intégration (sur 20 converti /5 par critère)}, width=\linewidth, height=7cm ]
\addplot[fill=popblue!50] coordinates {(Analyse, 4) (Problématique, 3) (Plan, 4) (Intro/Conc, 3) (Langue, 2) (Exemples, 3)};
\addplot[fill=poporange!50, pattern=north east lines] coordinates {(Analyse, 5) (Problématique, 5) (Plan, 5) (Intro/Conc, 5) (Langue, 5) (Exemples, 5)};
\legend {Élève type, Objectif max}
\end{axis}
\end{tikzpicture}
\legende{Profil d'un élève type comparé à l'objectif maximal. Ici, la remédiation devra porter en priorité sur la langue (2/5) puis sur la problématique et l'introduction/conclusion (3/5).}
\end{popfigure}

\begin{methode}[La co-évaluation entre pairs]
\begin{enumerate}
\item L'élève correcteur reçoit la copie d'un camarade et la \textbf{grille officielle}.
\item Il surligne en \textbf{VERT} les points forts (bon connecteur logique, bel exemple, problématique claire).
\item Il surligne en \textbf{ROUGE} les faiblesses (hors sujet, faute de syntaxe, argument manquant, exemple absent).
\item Il attribue une note sur /20 critère par critère et \textbf{justifie en deux lignes} sa notation.
\item Le professeur corrige ensuite à son tour et confronte les deux évaluations : l'écart entre la note du pair et la note du professeur est lui-même un objet de discussion en classe.
\end{enumerate}
\end{methode}

\begin{attention}[Pièges de la co-évaluation]
Ne pas noter « à l'amitié » : une note gonflée prive le camarade d'une remédiation utile. Inversement, toute sanction doit être \emph{justifiée} par un élément précis de la copie (citer la ligne ou la phrase fautive). Une remarque sans preuve n'aide pas à progresser.
\end{attention}

\courssub{La remédiation différenciée : les groupes de besoins}

L'observation des grilles montre que les erreurs ne sont pas les mêmes pour tous. Faire refaire le même exercice à toute la classe serait inefficace : on constitue donc des \cle{groupes de besoins}, comme un médecin prescrit un traitement adapté à chaque diagnostic.

\begin{center}
\begin{tabularx}{\linewidth}{|l|X|X|}
\hline
\rowcolor{popblueL} \textbf{Groupe} & \textbf{Difficulté diagnostiquée} & \textbf{Remédiation proposée} \\
\hline
Groupe A & Analyse du sujet défaillante (hors sujet, mots-clés ignorés) & Exercices ciblés : définir les termes de 5 sujets types Bac, repérer les présupposés, reformuler le sujet en une phrase. \\
\hline
Groupe B & Plan déséquilibré ou incohérent & Travail sur les types de plans (dialectique, thématique, analytique) ; squelettes de plans à trous à compléter. \\
\hline
Groupe C & Langue : syntaxe et ponctuation & Exercices de réécriture : transformer des phrases simples en phrases composées, ponctuer un texte dépouillé, corriger des phrases fautives. \\
\hline
\end{tabularx}
\end{center}

\begin{experience}[Atelier de remédiation du Groupe B : le squelette de plan à trous]
\textbf{Matériel} : fiche comportant un plan dont seule l'armature est donnée : « I. \ldots / A. \ldots / 1. \ldots / 2. \ldots / B. \ldots / II. \ldots ».\par
\textbf{Protocole} : (1) l'élève reçoit le sujet et une liste d'arguments et d'exemples en désordre ; (2) il classe chaque élément dans le squelette ; (3) il rédige les intitulés de parties et sous-parties ; (4) il vérifie l'équilibre (deux sous-parties minimum par partie).\par
\textbf{Observation} : la plupart des élèves réussissent le classement mais peinent à distinguer \cle{argument} (l'idée qui défend la thèse) et \cle{exemplebox} (le fait précis qui illustre l'idée).\par
\textbf{Interprétation} : l'argument répond à « pourquoi ? », l'exemple répond à « la preuve ? ». Entraînement : pour chaque argument trouvé, imposer la question « par exemple ? » afin de vérifier qu'une illustration existe.
\end{experience}

\begin{exercice}[Remédiation Groupe C : réécriture]
Ponctue correctement la phrase suivante et indique si elle est simple ou composée : « le technicien relève les mesures il consigne les résultats dans le rapport de stage puis il transmet le document au chef d'atelier ».
\end{exercice}

\begin{corrige}[Exercice de remédiation Groupe C]
Phrase ponctuée : « Le technicien relève les mesures ; il consigne les résultats dans le rapport de stage, puis il transmet le document au chef d'atelier. » Il s'agit d'une \textbf{phrase composée} : trois propositions indépendantes (ou juxtaposées/coordonnées) reliées par la ponctuation forte (point-virgule) et la coordination (« puis »). La majuscule initiale et le point final (ponctuation forte) encadrent la phrase ; la virgule (ponctuation faible) sépare les éléments à l'intérieur.
\end{corrige}

\courssub{Finalisation du projet d'étude : la soutenance orale sur \emph{Le Lion et la Perle}}

Le projet d'étude de l'œuvre de Wole Soyinka, négocié avec la classe au début de l'unité, se clôt par une \cle{soutenance orale} en groupes. Chaque groupe défend un thème choisi : \emph{la femme} (Sidi, Sadiku), \emph{le pouvoir} (Baroka, le Baalè), \emph{le théâtre} (la mise en scène, la danse, le mime), ou \emph{la langue} (anglais, yoruba, registres de langue).

\begin{methode}[Réussir sa soutenance orale]
\begin{enumerate}
\item Choisir une \textbf{problématique littéraire} précise (et non un thème vague).
\item Structurer le propos : introduction (accroche + problématique + annonce du plan), 2 ou 3 parties argumentées, conclusion.
\item Appuyer chaque idée sur un \textbf{passage précis de l'œuvre} (scène, réplique, didascalie).
\item Parler 5 minutes \textbf{sans lire ses notes}, en regardant l'auditoire.
\item Répondre aux questions du jury (professeur et camarades) en justifiant ses réponses.
\end{enumerate}
\end{methode}

\begin{attention}[Le piège du résumé]
Lors de la soutenance orale sur \emph{Le Lion et la Perle}, ne pas faire un résumé de l'intrigue : tout le monde connaît l'histoire de Sidi, Baroka et Lakunle. Il faut répondre à une \textbf{problématique littéraire}, par exemple : « Comment Soyinka théâtralise-t-il le choc des cultures ? » ou « Le personnage de Baroka est-il un tyran ou un stratège ? ». Raconter au lieu d'argumenter fait perdre la moitié des points.
\end{attention}

\begin{savaistu}[L'oral compte pour le contrôle continu]
L'oral de soutenance compte pour la note de \cle{contrôle continu} (CC). Il est évalué sur trois critères : la qualité de l'expression (voix, débit, regard), la solidité de l'argumentation (plan, exemples tirés du texte) et la connaissance de l'œuvre. S'entraîner à parler 5 minutes sans notes, avec une structure claire (introduction, 2 ou 3 parties, conclusion), est la meilleure préparation : chronométrez-vous à la maison !
\end{savaistu}

\courssub{Le bilan méta-cognitif : se connaître pour progresser}

\begin{definition}[Méta-cognition]
La \cle{méta-cognition} est la capacité à réfléchir sur sa propre manière d'apprendre : identifier ce que l'on sait faire, ce qui reste fragile, et quelles stratégies personnelles fonctionnent. C'est un outil décisif pour préparer le Baccalauréat technique.
\end{definition}

Chaque élève remplit sa \textbf{fiche de progression} en trois colonnes, à conserver dans son portfolio jusqu'aux révisions du Bac :

\begin{center}
\begin{tabularx}{\linewidth}{|X|X|X|}
\hline
\rowcolor{popblueL} \textbf{Mes acquis} & \textbf{Mes lacunes} & \textbf{Mes stratégies pour le Bac} \\
\hline
Ex. : je sais formuler une problématique ; je distingue phrase simple et phrase composée. & Ex. : mes plans manquent d'équilibre ; je confonds argument et exemple. & Ex. : refaire un squelette de plan par semaine ; relire ma copie en vérifiant la ponctuation forte. \\
\hline
\end{tabularx}
\end{center}

\begin{exemplebox}[Un chiffre à retenir : le cadrage des 25 séances]
Le module complet représente \textbf{25 séances}. Une répartition type : étude de l'œuvre \emph{Le Lion et la Perle} (6 séances), langue et communication (5 séances), méthodologie de la dissertation (10 séances), intégration et évaluation (4 séances). Respecter ce cadrage est crucial : chaque séance de retard sur l'œuvre est une séance en moins pour s'entraîner à la dissertation, épreuve reine du Bac technique. Un module bien tenu, c'est un programme terminé avant les révisions.
\end{exemplebox}

\begin{exoresolu}[Auto-évaluation finale]
Énoncé. Un élève obtient à la grille : Compréhension 4/5, Problématique 2/3, Plan 3/6, Langue 2/3, Exemples 2/3. (a) Calculez sa note. (b) Identifiez son groupe de besoins prioritaire. (c) Proposez-lui une stratégie concrète.
\tcblower
Solution détaillée.
(a) Note : $4 + 2 + 3 + 2 + 2 = 13/20$.
(b) Le critère le plus faible proportionnellement est le \textbf{plan} (3/6, soit 50 \%) : l'élève relève du \textbf{Groupe B} (plan déséquilibré).
(c) Stratégie : s'entraîner avec des squelettes de plans à trous, vérifier systématiquement que chaque partie contient au moins deux sous-parties, et appliquer le réflexe « argument $\rightarrow$ par exemple ? » pour garantir la présence d'une illustration sous chaque idée.
\end{exoresolu}

\begin{aretenir}[L'essentiel de la section]
\begin{itemize}
\item La \textbf{tâche d'intégration} mobilise en autonomie toute la chaîne : analyse du sujet $\rightarrow$ problématique $\rightarrow$ carte heuristique $\rightarrow$ plan détaillé $\rightarrow$ introduction et conclusion.
\item L'\textbf{évaluation formative} repose sur une grille critériée sur /20 et sur la co-évaluation entre pairs (vert = points forts, rouge = faiblesses, note justifiée).
\item La \textbf{remédiation} est différenciée : Groupe A (analyse du sujet), Groupe B (plan), Groupe C (langue et ponctuation).
\item La \textbf{soutenance orale} sur \emph{Le Lion et la Perle} répond à une problématique littéraire, jamais à un résumé ; elle compte pour le contrôle continu.
\item Le \textbf{bilan méta-cognitif} (acquis / lacunes / stratégies) transforme chaque élève en acteur de sa préparation au Baccalauréat technique.
\end{itemize}
\end{aretenir}

\newpage

\coursec{Activité d'intégration}

\courssub{Objectifs de l'activité}
Cette activité d'intégration vise à mobiliser l'ensemble des savoir-faire acquis au cours de la séquence « Dissertation : du sujet au plan détaillé ». Elle permet à l'élève de :
\begin{itemize}
    \item Analyser un sujet inédit en identifiant les termes clés, les relations logiques et les attendus du correcteur.
    \item Formuler une problématique pertinente, ouverte et débattue, ancrée dans le contexte camerounais.
    \item Mener une recherche d'idées structurée (carte heuristique) en s'appuyant sur des connaissances techniques, des exemples locaux et des références littéraires (\emph{Le Lion et la Perle} de Wole Soyinka).
    \item Élaborer un plan détaillé complet (introduction, trois parties équilibrées, conclusion) avec arguments, contre-arguments et illustrations précises.
    \item Rédiger l'introduction et la conclusion en respectant la ponctuation forte, la syntaxe (phrase simple / phrase composée) et les connecteurs argumentatifs.
    \item S'auto-évaluer à l'aide d'une grille critériée et participer à un débriefing collectif pour identifier les pièges récurrents.
\end{itemize}
L'objectif final est de placer l'élève en situation d'examen (Probatoire technique) afin de consolider son autonomie et sa rigueur méthodologique.

\begin{aretenir}[Rappel des savoirs mobilisés]
Cette activité suppose acquis : l'analyse d'un sujet de dissertation et la formulation de la problématique ; les techniques de recherche des idées ; l'élaboration du plan détaillé ; les caractéristiques et stratégies du texte argumentatif ; la ponctuation (notamment la ponctuation forte) ; la structure de la phrase (phrase simple et phrase composée) ; les situations et conditions de la communication.
\end{aretenir}

\courssub{Situation}
\begin{quote}
\textbf{Contexte :} Vous êtes élève en Première technique (spécialité Maintenance industrielle) au Lycée Technique de Douala. Dans le cadre de la préparation au Probatoire, votre professeur organise un \emph{Bac Blanc} de dissertation. Le sujet tiré au sort est :
\end{quote}
\begin{center}
\textbf{« La maintenance préventive : coût ou investissement pour l'entreprise camerounaise ? »}
\end{center}
\begin{quote}
\textbf{Données chiffrées fournies avec le sujet (étude fictive d'entraînement, inspirée des publications de la CCIMA) :}
\begin{itemize}
    \item Coût moyen d'une panne non planifiée : \num{2500000} FCFA par incident.
    \item Budget annuel moyen alloué à la maintenance préventive : \num{12000000} FCFA pour une PME de 50 salariés.
    \item Taux de réduction des pannes après mise en place d'un plan préventif : 68\,\%.
    \item Délai moyen de retour sur investissement (ROI) observé : 14 mois.
\end{itemize}
\textbf{Contraintes de l'épreuve :} Durée 1\,h\,30. Brouillon autorisé. Remise de cinq livrables (voir consignes). Évaluation sur 20 points selon la grille critériée fournie.
\end{quote}

\begin{attention}[Lecture critique des données]
Avant d'utiliser un chiffre dans une dissertation, l'élève doit vérifier sa cohérence. Ici, si un plan préventif réduit les pannes de 68\,\%, l'économie moyenne par incident évité est de :
\[ \num{2500000} \times \frac{68}{100} = \num{1700000} \text{ FCFA}. \]
Ce calcul simple, mené au brouillon, fournit un argument chiffré solide pour la deuxième partie du développement.
\end{attention}

\courssub{Consignes}
\begin{enumerate}
    \item \textbf{Fiche d'analyse lexicale et logique} : Relevez les mots-clés, définissez les termes techniques (maintenance préventive, coût, investissement, ROI), identifiez la structure du sujet (thème, question, portée) et dégagez les présupposés.
    \item \textbf{Problématique} : Formulez une unique question problématisée qui met en tension les deux pôles (coût / investissement) et qui appelle une réponse nuancée.
    \item \textbf{Carte heuristique (mindmap)} : Sur une feuille A4, représentez visuellement les idées principales, les arguments, les contre-arguments et les exemples (camerounais, techniques, littéraires). Utilisez des couleurs pour distinguer les trois parties futures.
    \item \textbf{Plan détaillé complet} : Rédigez l'introduction (accroche, définition, problématique, annonce du plan), trois parties équilibrées (I, II, III) chacune avec deux sous-parties (A, B), arguments, exemples chiffrés et références, puis la conclusion (bilan, ouverture).
    \item \textbf{Introduction et conclusion rédigées} : Soignez la ponctuation forte (deux-points, point-virgule, points de suspension), variez les structures de phrases (simples, composées), employez les connecteurs argumentatifs (\emph{en outre, toutefois, par conséquent, en définitive}).
    \item \textbf{Correction croisée} : Échangez votre copie avec un binôme, notez sur la grille critériée (analyse, problématique, idées, plan, langue) et ajoutez un commentaire constructif.
    \item \textbf{Débriefing collectif} : Participez à la synthèse orale animée par le professeur : pièges relevés, bonnes pratiques, remédiation.
\end{enumerate}

\courssub{Ressources à mobiliser}
\begin{itemize}
    \item \textbf{Méthode d'analyse du sujet} : \emph{Repérage lexical $\rightarrow$ Définition $\rightarrow$ Structure logique $\rightarrow$ Présupposés $\rightarrow$ Problématique}. (Voir fiche méthode « Analyse d'un sujet de dissertation ».)
    \item \textbf{Techniques de recherche d'idées} : brainstorming, carte heuristique, tableau « Pour / Contre », analogie avec \emph{Le Lion et la Perle} (tradition et modernité, coût du changement).
    \item \textbf{Structure du plan détaillé} : introduction (accroche, mise en contexte, problématique, annonce du plan), développement en trois parties (thèse, antithèse, synthèse), conclusion (bilan + ouverture).
    \item \textbf{Règles de ponctuation forte} : deux-points pour annoncer ou expliquer, point-virgule pour lier des propositions longues, points de suspension pour l'ouverture.
    \item \textbf{Connecteurs argumentatifs} : \emph{en premier lieu, de surcroît, néanmoins, par ailleurs, en conséquence, pour conclure}.
    \item \textbf{Grille critériée d'évaluation} (20 pts) : Analyse (4), Problématique (3), Idées / Mindmap (4), Plan détaillé (5), Introduction + Conclusion (4), Langue / Ponctuation (3). Total : $4+3+4+5+4+3 = 23$ ramené à 20 par harmonisation, ou pondération ajustée par le professeur.
\end{itemize}

\begin{methode}[De la carte heuristique au plan détaillé]
\begin{enumerate}
    \item Placer le sujet au centre de la feuille et dégager trois branches principales correspondant aux trois parties envisagées.
    \item Sur chaque branche, accrocher deux sous-branches (A et B) : ce seront les sous-parties.
    \item Ajouter à chaque sous-branche un argument, un exemple chiffré ou local, et si possible un contre-argument à réfuter.
    \item Numéroter les branches dans l'ordre logique (thèse, antithèse, synthèse) : la carte devient le plan.
    \item Rédiger l'introduction et la conclusion à partir de la problématique placée au centre.
\end{enumerate}
\end{methode}

\courssub{Démarche et corrigé commenté}
\begin{exoresolu}{Corrigé type — Sujet : « La maintenance préventive : coût ou investissement pour l'entreprise camerounaise ? »}
\textbf{1. Fiche d'analyse lexicale et logique}
\begin{itemize}
    \item \textbf{Mots-clés} : maintenance préventive, coût, investissement, entreprise camerounaise.
    \item \textbf{Définitions} : la \cle{maintenance préventive} est l'ensemble des actions planifiées (inspection, lubrification, remplacement programmé) destinées à éviter les pannes ; le \cle{coût} est une sortie financière immédiate ; l'\cle{investissement} est une dépense générant un retour futur (gain de productivité, réduction des arrêts) ; l'\cle{entreprise camerounaise} désigne ici les PME/PMI opérant dans un contexte d'infrastructures énergétiques instables et de main-d'œuvre qualifiée limitée.
    \item \textbf{Structure} : sujet à double pôle (coût / investissement) + portée géographique (Cameroun). Question implicite : « La maintenance préventive est-elle une charge ou un levier de performance ? »
    \item \textbf{Présupposés} : l'entreprise a le choix entre maintenance curative (réactive) et préventive ; les données chiffrées sont disponibles ; le décideur cherche la rentabilité.
\end{itemize}
\textbf{2. Problématique}
\emph{« Dans quelle mesure la maintenance préventive, malgré son coût initial élevé, constitue-t-elle un investissement stratégique pour la compétitivité et la pérennité des entreprises camerounaises ? »}
\textbf{3. Carte heuristique (représentation)}
\begin{center}
\begin{tikzpicture}[mindmap, grow cyclic, every node/.style=concept, concept color=popblue!60,
    level 1/.append style={level distance=4.5cm,sibling angle=120},
    level 2/.append style={level distance=3cm,sibling angle=60}]
\node {Maintenance préventive}
    child[concept color=popgreen!60] {node {I. Coût apparent}
        child {node {A. Budget direct (12 M FCFA/an)}}
        child {node {B. Formation et outillage}}}
    child[concept color=poporange!60] {node {II. Gains mesurables}
        child {node {A. Réduction des pannes : 68\,\%}}
        child {node {B. ROI : 14 mois ; productivité en hausse}}}
    child[concept color=poppurple!60] {node {III. Facteurs contextuels}
        child {node {A. Instabilité énergétique (coupures)}}
        child {node {B. Cadre réglementaire et formation}}};
\end{tikzpicture}
\end{center}
\textbf{4. Plan détaillé}
\begin{itemize}
    \item \textbf{Introduction} : accroche (citation attribuée à la tradition orale et mise en relation avec \emph{Le Lion et la Perle} : le refus du changement coûte plus cher que le changement lui-même), définition des termes, problématique, annonce du plan (I. Coût à court terme ; II. Retour sur investissement ; III. Conditions de réussite au Cameroun).
    \item \textbf{I. La maintenance préventive : un coût immédiat assumé}
        \begin{itemize}
            \item A. Budget direct et charges annexes (ex. : 12 millions de FCFA par an pour une PME de 50 salariés ; formation des techniciens, achat de capteurs).
            \item B. Perception comptable : charge d'exploitation, impact sur la trésorerie à court terme.
        \end{itemize}
    \item \textbf{II. Un investissement rentable à moyen terme}
        \begin{itemize}
            \item A. Réduction drastique des pannes (68\,\% selon les données fournies), soit une économie moyenne de \num{1700000} FCFA par incident évité.
            \item B. ROI à 14 mois, gain de productivité, image de fiabilité vis-à-vis des clients (ex. : une usine agroalimentaire de la zone industrielle de Douala-Bonabéri).
        \end{itemize}
    \item \textbf{III. Les conditions de succès dans le contexte camerounais}
        \begin{itemize}
            \item A. Adaptation aux aléas énergétiques (groupes électrogènes, onduleurs) et à la pénurie de pièces détachées (stock de sécurité).
            \item B. Cadre incitatif : avantages fiscaux (amortissement), certification qualité, formation continue en partenariat avec les établissements d'enseignement technique.
        \end{itemize}
    \item \textbf{Conclusion} : bilan (coût initial réel mais investissement stratégique), ouverture vers la maintenance prédictive (capteurs connectés) et la responsabilité sociétale de l'entreprise.
\end{itemize}
\textbf{5. Introduction rédigée (extrait)}
\emph{« Dans \emph{Le Lion et la Perle}, Wole Soyinka montre que refuser le changement, c'est souvent payer le prix le plus lourd. Cette leçon résonne avec acuité pour l'industriel camerounais confronté au dilemme de la maintenance préventive : charge immédiate ou levier de performance ? Si le budget annuel de 12 millions de FCFA pèse sur la trésorerie, les données fournies révèlent une réduction de 68\,\% des pannes et un retour sur investissement à 14 mois. Dès lors, dans quelle mesure la maintenance préventive, malgré son coût initial élevé, constitue-t-elle un investissement stratégique pour la compétitivité et la pérennité des entreprises camerounaises ? Nous analyserons d'abord le coût apparent (I), puis les gains mesurables (II), avant d'examiner les conditions de réussite dans le contexte national (III). »}
\textbf{6. Conclusion rédigée (extrait)}
\emph{« En définitive, la maintenance préventive n'est pas une simple ligne budgétaire ; elle est le socle d'une résilience industrielle. Les chiffres parlent : 68\,\% de pannes en moins, un retour sur investissement à 14 mois, une productivité accrue. Toutefois, sa réussite exige une adaptation aux réalités camerounaises — instabilité énergétique, pénurie de pièces détachées — et un cadre incitatif fort. L'avenir appartient à la maintenance prédictive, où les capteurs connectés transformeront durablement le coût en avantage concurrentiel... »}

\tcblower
\textbf{Commentaires du corrigé}
\begin{itemize}
    \item L'analyse lexicale isole les termes techniques et replace le sujet dans son environnement économique camerounais.
    \item La problématique met en tension « coût initial » et « investissement stratégique » tout en intégrant la dimension géographique ; elle est unique, ouverte et appelle une réponse nuancée.
    \item La carte heuristique structure visuellement les trois parties et facilite la rédaction du plan : chaque branche devient une partie, chaque sous-branche une sous-partie.
    \item Le plan détaillé respecte l'équilibre thèse / antithèse / synthèse ; chaque sous-partie est étayée par des données chiffrées et un exemple local.
    \item Le calcul $68\,\% \times \num{2500000} = \num{1700000}$ FCFA transforme une donnée brute en argument : c'est la démarche « rédiger et justifier » attendue au Probatoire.
    \item L'introduction et la conclusion illustrent l'usage de la ponctuation forte (deux-points, point-virgule, tirets, points de suspension) et des connecteurs argumentatifs (\emph{si... dès lors, toutefois, en définitive}).
    \item La grille critériée permet une évaluation transparente : chaque indicateur est pondéré, l'élève identifie ses points forts et ses axes de progrès.
\end{itemize}
\end{exoresolu}

\begin{attention}[Pièges fréquents relevés en correction croisée]
\begin{itemize}
    \item Problématique descriptive (« Qu'est-ce que la maintenance préventive ? ») au lieu d'une question problématisée qui met en tension deux pôles.
    \item Plan déséquilibré : une partie très développée, les autres expédiées ; ou trois parties sans lien logique visible.
    \item Absence d'exemples chiffrés : les données de l'énoncé doivent être exploitées et, si possible, combinées par un calcul simple.
    \item Ponctuation forte mal employée : deux-points sans annonce réelle, point-virgule entre deux propositions sans rapport.
    \item Hors-sujet technique : la dissertation de français n'est pas un exposé de spécialité ; les connaissances techniques servent l'argumentation, elles ne s'y substituent pas.
\end{itemize}
\end{attention}

\courssub{Bilan de l'activité}
Cette simulation de Bac Blanc a permis de mettre en évidence plusieurs acquis majeurs :
\begin{itemize}
    \item \textbf{Autonomie méthodologique} : les élèves ont appliqué, sans guidage pas à pas, la chaîne complète « analyse $\rightarrow$ problématique $\rightarrow$ idées $\rightarrow$ plan $\rightarrow$ rédaction ».
    \item \textbf{Transfert des connaissances techniques} : l'exploitation des données chiffrées et l'intégration d'exemples concrets (zone industrielle de Douala, coupures d'électricité) montrent la capacité à mobiliser le savoir technique dans un exercice littéraire.
    \item \textbf{Maîtrise de la langue} : l'attention portée à la ponctuation forte, à la variation syntaxique (phrase simple, phrase composée) et aux connecteurs a produit des introductions et des conclusions de qualité attendue au Probatoire.
    \item \textbf{Esprit critique et remédiation} : la correction croisée a révélé des pièges récurrents (problématique trop descriptive, plan déséquilibré, absence d'exemples chiffrés) et a favorisé l'auto-correction.
    \item \textbf{Ancrage culturel} : la référence à \emph{Le Lion et la Perle} a enrichi l'argumentation en liant tradition et modernité, illustrant la transversalité du programme.
\end{itemize}
Les difficultés résiduelles (gestion du temps, formulation d'une problématique unique, articulation thèse / antithèse) feront l'objet de séances de remédiation ciblées : exercices d'analyse de sujets, ateliers de carte heuristique chronométrés, entraînement à l'emploi des connecteurs. L'activité valide ainsi l'atteinte des compétences « Produire un plan détaillé de dissertation » et « Rédiger et justifier une démarche, une réponse ou une conclusion » attendues pour le Probatoire technique.

\newpage

\coursec{Méthodes et exercices résolus}

\coursec{Leçon N02 : Produire le plan détaillé d'une dissertation}

\courssub{Analyser un sujet de dissertation}

\begin{definition}[Sujet de dissertation]
Énoncé qui invite à une réflexion organisée, argumentée et illustrée. Il peut prendre la forme d'une question (fermée ou ouverte), d'une affirmation à discuter, ou d'une citation à commenter. En Première technique, il porte souvent sur des thèmes de société, des œuvres au programme (comme \emph{Le Lion et la Perle}) ou des notions de civilisation.
\end{definition}

\begin{methode}[Analyser un sujet de dissertation : 5 étapes indispensables]
\begin{enumerate}
\item \textbf{Lire et relire} le sujet (au moins trois fois) ; souligner les \cle{mots-clés} (notions centrales, verbes d'action, noms propres, œuvres citées, connecteurs logiques).
\item \textbf{Identifier le type de sujet} :
  \begin{itemize}
  \item Sujet de \textbf{réflexion} (question générale : \og La tradition est-elle un obstacle au progrès ? \fg).
  \item Sujet de \textbf{commentaire de citation} (\og "L'école est l'arme la plus puissante...", Mandela. Commentez. \fg).
  \item Sujet \textbf{lié à l'œuvre} (\og Dans \emph{Le Lion et la Perle}, la tradition triomphe-t-elle de la modernité ? \fg).
  \end{itemize}
\item \textbf{Définir chaque mot-clé} avec ses propres mots (dictionnaire de la langue, contexte culturel camerounais/africain) pour baliser le champ de la réflexion et éviter le hors-sujet.
\item \textbf{Dégager l'enjeu (la tension)} : repérer l'opposition, le paradoxe ou le problème caché (ex. : tradition vs modernité, individuel vs collectif, conservation vs changement).
\item \textbf{Formuler la problématique} : transformer l'enjeu en question(s) ouverte(s) qui guideront tout le devoir.
\end{enumerate}
\textbf{Réflexe Bac} : Sujet bien analysé = Mots-clés définis + Enjeu dégagé + Problématique interrogative et ouverte.
\end{methode}

\begin{exoresolu}[Analyse complète d'un sujet type Probatoire]
Énoncé. Sujet : \emph{« La coutume est-elle un frein à l'émancipation de la femme africaine ? »} Analysez ce sujet et formulez la problématique.
\tcblower
Solution détaillée.

\textbf{Étape 1 : Mots-clés.} \cle{Coutume}, \cle{frein}, \cle{émancipation}, \cle{femme africaine}.

\textbf{Étape 2 : Type de sujet.} Sujet de réflexion sur un thème de société, forme interrogative fermée (appelant oui/non), mais exigeant une réponse nuancée.

\textbf{Étape 3 : Définitions contextuelles.}
\begin{itemize}
\item \emph{Coutume} : ensemble de règles coutumières, rites, interdits régissant la vie sociale (mariage, héritage, place dans la cité) transmis oralement.
\item \emph{Frein} : obstacle, ralentissement, blocage d'un processus d'évolution.
\item \emph{Émancipation} : accession à l'autonomie (juridique, économique, intellectuelle, politique), libération des tutelles.
\item \emph{Femme africaine} : sujet historique et social, prise ici dans la diversité de ses statuts (rurale/urbaine, lettrée/non lettrée).
\end{itemize}

\textbf{Étape 4 : Enjeu / Tension.} Opposition entre la \textbf{légitimité coutumière} (cohésion du groupe, identité, protection) et l'\textbf{aspiration à l'autonomie} (droits humains, scolarisation, activité génératrice de revenus). Le sujet interroge la compatibilité entre fidélité aux ancêtres et conquête de soi.

\textbf{Étape 5 : Problématique.} \og Dans quelle mesure les normes coutumières constituent-elles un obstacle à l'autonomie de la femme africaine, ou peuvent-elles, réinterprétées, devenir un levier pour son émancipation ? \fg

\textbf{Vérification.} Reprend les 4 mots-clés ; ouvre deux pistes (obstacle / levier) ; annonce un plan dialectique. Valide.
\end{exoresolu}

\courssub{Formuler la problématique}

\begin{propriete}[Critères d'une bonne problématique]
Une problématique recevable au Probatoire / Baccalauréat technique doit :
\begin{enumerate}
\item Être une \textbf{question ouverte} (commencer par \og Dans quelle mesure \fg, \og Comment \fg, \og Pourquoi \fg, \og De quelle manière \fg).
\item \textbf{Reprendre les mots-clés} du sujet (ou leurs synonymes précis).
\item \textbf{Exprimer une tension} (opposition, paradoxe, double sens).
\item \textbf{Annoncer implicitement le plan} (souvent en deux ou trois temps : thèse / antithèse / synthèse).
\item Être \textbf{autonome} : compréhensible sans relire le sujet.
\end{enumerate}
\end{propriete}

\begin{methode}[De l'enjeu à la problématique : 4 gestes]
\begin{enumerate}
\item Repérer la \cle{tension} (deux pôles qui s'affrontent : ex. fidélité / rupture).
\item Refuser la question fermée (\og Faut-il... ? \fg $\rightarrow$ \og Dans quelle mesure doit-on... ? \fg).
\item Rédiger une \textbf{question principale} large, puis \textbf{deux questions secondaires} qui découpent le débat (souvent : causes/manifestations // limites/conditions // perspectives).
\item Tester : si je réponds aux questions secondaires, est-ce que je réponds à la principale ? Est-ce que je traite le sujet ?
\end{enumerate}
\textbf{Formules types} : \og Dans quelle mesure [concept A] et [concept B] sont-ils compatibles ? \fg ; \og Comment concilier [valeur 1] et [exigence 2] ? \fg ; \og Quelles sont les conditions pour que [phénomène] ne soit pas [risque] ? \fg
\end{methode}

\begin{exoresolu}[Problématique et questions secondaires]
Énoncé. Sujet : \emph{« L'école moderne menace-t-elle l'identité culturelle camerounaise ? »} Proposez une problématique principale et deux questions secondaires.
\tcblower
Solution détaillée.

\textbf{Tension.} D'un côté, l'école véhicule des valeurs exogènes (langues officielles, savoirs universels, individualisme) ; de l'autre, elle est le lieu de transmission du patrimoine (histoire, langues nationales, civisme).

\textbf{Problématique principale.} \og Dans quelle mesure l'école moderne, par ses contenus et ses méthodes, fragilise-t-elle ou au contraire renforce-t-elle l'identité culturelle camerounaise ? \fg

\textbf{Questions secondaires (annonce du plan).}
\begin{itemize}
\item Q1 (Thèse) : \og En quoi l'école moderne, par l'importation de modèles étrangers, peut-elle déstructurer les repères culturels traditionnels ? \fg
\item Q2 (Antithèse/Synthèse) : \og Comment l'école, en intégrant les langues et cultures nationales, peut-elle devenir le creuset d'une identité camerounaise renouvelée ? \fg
\end{itemize}

\textbf{Vérification.} Mots-clés présents (\emph{école}, \emph{identité}, \emph{camerounaise}) ; tension explicite ; plan dialectique annoncé.
\end{exoresolu}

\courssub{Rechercher les idées et les arguments}

\begin{methode}[Recherche et classement des idées : la méthode des 3 sources]
\begin{enumerate}
\item \textbf{Brainstorming (brouillon) :} 5 à 10 minutes. Noter \emph{tout} sans censure : mots, noms, proverbes, souvenirs, titres d'œuvres, faits d'actualité.
\item \textbf{Puiser dans les 3 réservoirs obligatoires} :
  \begin{itemize}
  \item \textbf{Œuvres au programme} : \emph{Le Lion et la Perle} (Sidi, Baroka, Lakunle, Sadiku), autres textes étudiés.
  \item \textbf{Connaissances personnelles \& actualité camerounaise} : vie familiale, village, ville, tontine, chefferies, médias, lois, faits divers.
  \item \textbf{Culture générale / Lectures} : romans africains (Betty, Oyono, Sembène), contes, proverbes, histoire, géographie.
  \end{itemize}
\item \textbf{Classer selon la problématique} : deux colonnes (Thèse / Antithèse) ou trois (Constat / Causes / Solutions) selon le plan choisi.
\item \textbf{Valider chaque argument} : 1 idée force + 1 connecteur logique + 1 exemple précis (nom, chiffre, scène, proverbe). \textbf{Éliminer} les doublons et le hors-sujet. Garder \textbf{2 à 3 arguments solides par partie}.
\end{enumerate}
\textbf{Règle d'or} : Argument sans exemple = affirmation gratuite = 0 point au Probatoire.
\end{methode}

\begin{exoresolu}[Recherche et classement pour un plan dialectique]
Énoncé. Sujet : \emph{« Le progrès technique améliore-t-il les conditions de vie au village ? »} Classez les idées pour un plan dialectique.
\tcblower
Solution détaillée.

\textbf{Brainstorming sélectionné.} Pompage solaire, téléphone portable, moto-taxi, pesticides, exode rural, perte des savoirs agricoles, école au village, centre de santé, individualisme, tontine modernisée, radio communautaire.

\begin{center}
\begin{tabularx}{\linewidth}{|>{\columncolor{popblueL}}X|>{\columncolor{poporangeL}}X|}
\hline
\textbf{Thèse : Le progrès améliore la vie} & \textbf{Antithèse : Le progrès fragilise les équilibres} \\
\hline
\textbf{Accès à l'eau/santé} : forage solaire $\rightarrow$ baisse maladies hydriques. Ex. : village de Koutaba. & \textbf{Dépendance technique} : panne pompe = retour corvée d'eau ; pièces détachées introuvables. \\
\hline
\textbf{Désenclavement} : moto-taxi + téléphone $\rightarrow$ vente produits champs, lien famille ville. Ex. : coopérative café Ngoundéré. & \textbf{Exode des jeunes} : attirance ville, main-d'œuvre agricole vieillissante. Ex. : villages fantômes du Nord. \\
\hline
\textbf{Éducation/Santé} : école/radio au village $\rightarrow$ alphabétisation, prévention. & \textbf{Pollution / Santé} : pesticides mal dosés $\rightarrow$ sols morts, cancers. Ex. : cotonnière Nord. \\
\hline
\end{tabularx}
\end{center}

\textbf{Sélection finale (2 arguments/partie).}
\begin{itemize}
\item I. A. Santé/Eau (forage) ; B. Économie (moto/téléphone).
\item II. A. Exode rural ; B. Environnement (pesticides).
\end{itemize}
Transition : \og Si le progrès apporte des services vitaux, ses effets pervers menacent la survie même du village. Comment le maîtriser ? \fg
\end{exoresolu}

\courssub{Élaborer le plan détaillé}

\begin{definition}[Plan détaillé]
Squelette complet de la dissertation rédigé au brouillon. Il comporte : Introduction (amorce, définitions, problématique, annonce), Parties (I, II, III), Sous-parties (A, B), Arguments (1, 2) avec exemples, Transitions, Conclusion (bilan, réponse, ouverture). Il se rédige \textbf{au brouillon} avant la rédaction finale.
\end{definition}

\begin{propriete}[Les trois plans types au Probatoire Technique]
\begin{center}
\begin{tabularx}{\linewidth}{|l|X|X|}\hline
\rowcolor{popblueL} \textbf{Type de plan} & \textbf{Structure} & \textbf{Quand l'utiliser ?} \\ \hline
\textbf{Dialectique} (Thèse / Antithèse / Synthèse) & I. Pour / II. Contre / III. Dépassement & Sujet polémique, question fermée, opposition nette (ex. \og La tradition freine-t-elle le progrès ? \fg). \\ \hline
\textbf{Analytique} (Constat / Causes / Conséquences ou Solutions) & I. Constat / II. Causes / III. Solutions & Sujet "problème-solution", analyse d'un phénomène (ex. \og Causes et conséquences de l'exode rural \fg). \\ \hline
\textbf{Thématique} (Aspects / Dimensions) & I. Aspect 1 / II. Aspect 2 / III. Aspect 3 & Sujet large, description/analyse d'une notion (ex. \og Les formes de la solidarité en milieu bantou \fg). \\ \hline
\end{tabularx}
\end{center}
\textbf{Choix stratégique} : 80\% des sujets de Probatoire appellent le plan \textbf{dialectique}.
\end{propriete}

\begin{methode}[Construire le plan détaillé pas à pas]
\begin{enumerate}
\item \textbf{Choisir le type de plan} adapté à la problématique (majoritairement dialectique).
\item \textbf{Rédiger l'Introduction (schéma)} : Amorce (citation, proverbe, fait d'actualité) $\rightarrow$ Définitions $\rightarrow$ Problématique $\rightarrow$ Annonce du plan (\og Nous examinerons d'abord... puis... enfin... \fg).
\item \textbf{Détailler chaque partie} (I, II, III) :
  \begin{itemize}
  \item \textbf{Idée directrice (A)} : phrase nominale ou verbale résumant la thèse de la sous-partie.
  \item \textbf{Arguments (1, 2)} : phrases complètes, connectées (\emph{en effet, par ailleurs, de surcroît}).
  \item \textbf{Exemples précis} : entre parenthèses ou après deux-points (Personnage d'œuvre, lieu camerounais, proverbe, chiffre, loi).
  \end{itemize}
\item \textbf{Rédiger les transitions} (1 phrase : bilan de la partie finie + annonce de la suivante).
\item \textbf{Rédiger la Conclusion (schéma)} : Bilan synthétique $\rightarrow$ Réponse nette à la problématique $\rightarrow$ Ouverture (autre œuvre, actualité, question prospective).
\item \textbf{Numérotation rigoureuse} : I, II, III / A, B / 1, 2.
\end{enumerate}
\end{methode}

\begin{exoresolu}[Plan détaillé complet : Sujet dialectique]
Énoncé. Sujet : \emph{« La tradition est-elle un obstacle au progrès ? »} Élaborez le plan détaillé.
\tcblower
Solution détaillée.

\textbf{INTRODUCTION}
\begin{itemize}
\item \textbf{Amorce} : Proverbe bamiléké : \og On ne peut pas marcher en avant en regardant seulement en arrière. \fg
\item \textbf{Définitions} : Tradition = héritage coutumier ; Progrès = amélioration conditions de vie (technique, social).
\item \textbf{Problématique} : La tradition, héritage du passé, empêche-t-elle réellement le progrès ou peut-elle l'accompagner et le guider ?
\item \textbf{Annonce} : Nous verrons que la tradition peut freiner le changement (I), avant de montrer qu'elle peut aussi en être le socle moral (II), pour conclure sur la nécessité d'une tradition vivante (III).
\end{itemize}

\textbf{I. La tradition peut freiner le progrès (Thèse)}
\begin{itemize}
\item \textbf{A. Le poids des interdits bloque l'innovation technique.}
  \begin{enumerate}
  \item Tabous fonciers / refus machines agricoles. Ex. : refus du labour attelé dans certains chefferies Grassfields (années 60).
  \item Crainte de la modernité médicale (accouchement à l'hôpital vs sage-femme traditionnelle).
  \end{enumerate}
\item \textbf{B. La gérontocratie limite l'initiative des jeunes.}
  \begin{enumerate}
  \item Pouvoir des anciens / silence des cadets. Ex. : Baroka dans \emph{Le Lion et la Perle} use de la ruse traditionnelle pour empêcher le chemin de fer.
  \item Héritage coutumier excluant les filles (terre, titres).
  \end{enumerate}
\end{itemize}
\textbf{Transition} : \og Toutefois, réduire la tradition à un frein serait ignorer sa fonction de régulation sociale. \fg

\textbf{II. La tradition peut guider et humaniser le progrès (Antithèse)}
\begin{itemize}
\item \textbf{A. Les valeurs traditionnelles offrent un cadre éthique au changement.}
  \begin{enumerate}
  \item Solidarité / Tontine $\rightarrow$ financement micro-projets modernes. Ex. : Groupements d'initiative commune (GIC) au Cameroun.
  \item Respect aînés $\rightarrow$ transmission savoir-faire adaptés (agroforesterie).
  \end{enumerate}
\item \textbf{B. L'enracinement culturel conditionne la pertinence du développement.}
  \begin{enumerate}
  \item Éducation bilingue / programmes intégrant cultures locales. Ex. : Manuels \emph{Elan-Afrique} (français + langues nationales).
  \item Justice coutumière + droit moderne $\rightarrow$ paix sociale (ex. \emph{Ndop} chez les Bamiléké).
  \end{enumerate}
\end{itemize}
\textbf{Transition} : \og La tradition n'est donc pas figée ; elle se réinvente au contact du progrès. \fg

\textbf{III. Vers une tradition dynamique : le dialogue des temps (Synthèse)}
\begin{itemize}
\item \textbf{A. La réinterprétation sélective : garder le sens, changer la forme.}
  \begin{enumerate}
  \item Modernisation rites initiatiques (durée, hygiène) sans perdre le sens symbolique.
  \item Chefferies 2.0 : chefs connectés, médiateurs développement local.
  \end{enumerate}
\item \textbf{B. La responsabilité de la jeunesse : acteurs de la synthèse.}
  \begin{enumerate}
  \item Jeunes formés (école moderne) + initiés (coutume) = ingénieurs culturels.
  \item Ex. : Start-ups camerounaises valorisant pharmacopée traditionnelle (ex. \emph{Phytothérapie}).
  \end{enumerate}
\end{itemize}

\textbf{CONCLUSION}
\begin{itemize}
\item \textbf{Bilan} : La tradition n'est obstacle que si elle se fige ; elle est ressource si elle s'adapte.
\item \textbf{Réponse} : Le progrès a besoin de la tradition pour ne pas être déshumanisé ; la tradition a besoin du progrès pour ne pas s'étioler.
\item \textbf{Ouverture} : \og Et si le vrai progrès, c'était l'art de faire dialoguer les temps ? \fg (Perspective : rôle de l'école technique dans cette médiation).
\end{itemize}
\end{exoresolu}

\courssub{Maîtriser la langue : Ponctuation, Syntaxe \& Argumentation}

\begin{definition}[Ponctuation forte et structure de la phrase]
La \cle{ponctuation forte} (point \textbf{.}, point d'interrogation \textbf{?}, point d'exclamation \textbf{!}, points de suspension \textbf{...}) délimite les \textbf{phrases} (unités syntaxiques autonomes, une idée complète, un verbe conjugué).
\begin{itemize}
\item \textbf{Phrase simple} : une seule proposition (un verbe conjugué). Ex. : \og La tradition guide le progrès. \fg
\item \textbf{Phrase composée} : plusieurs propositions liées par \textbf{juxtaposition} (ponctuation faible : , ; :), \textbf{coordination} (mais, or, et, donc, or, ni, car) ou \textbf{subordination} (conjonctions : que, quand, si, parce que ; relatifs : qui, que, dont).
\end{itemize}
\end{definition}

\begin{methode}[Ponctuer et structurer son plan (et sa dissertation)]
\begin{enumerate}
\item \textbf{Séparer les phrases} par un point (majusucre suivante). Ne pas confondre virgule et point.
\item \textbf{Utiliser le point-virgule (;)} pour lier deux propositions proches de sens sans conjonction (ex. : \og La tradition unit ; le progrès divise. \fg).
\item \textbf{Utiliser les deux-points (:)} pour annoncer explication, énumération ou citation.
\item \textbf{Encadrer les citations} par des guillemets français \og \fg.
\item \textbf{Relire à l'oral} : chaque point = chute de voix ; chaque virgule = pause courte. Vérifier : 1 verbe conjugué = 1 phrase (sauf subordination).
\end{enumerate}
\end{methode}

\begin{exoresolu}[Correction syntaxique et ponctuation]
Énoncé. Corrigez ce brouillon : \og la tradition transmet les valeurs elle est sacrée cependant certains rites sont dangereux il faut les réformer comment faire \fg
\tcblower
Solution détaillée.

\textbf{Analyse :} 4 propositions principales + 1 question. Verbes : \emph{transmet, est, sont, faut, faire}.

\textbf{Correction :}
\og La tradition transmet les valeurs ; elle est sacrée. Cependant, certains rites sont dangereux. Il faut donc les réformer : comment faire ? \fg

\textbf{Justifications :}
\begin{itemize}
\item Point-virgule entre deux propositions courtes, liées, sans conjonction.
\item Point après \og sacrée \fg (fin idée 1). Majuscule à \og Cependant \fg (adverbe de liaison en tête de phrase + virgule).
\item Deux-points avant la question rhétorique qui conclut l'argumentation.
\item Point d'interrogation final obligatoire.
\end{itemize}
\end{exoresolu}

\begin{definition}[Texte argumentatif : structure et stratégies]
Un texte argumentatif vise à \textbf{convaincre} (raison) ou \textbf{persuader} (émotion) un \textbf{destinataire} (ici : le correcteur du Probatoire). Sa structure de base par paragraphe :
\textbf{Idée directrice (Thèse partielle)} $\rightarrow$ \textbf{Argument(s)} (raisonnement) $\rightarrow$ \textbf{Exemple(s) / Preuve(s)} $\rightarrow$ \textbf{Mini-conclusion / Transition}.
\end{definition}

\begin{propriete}[Stratégies argumentatives attendues]
\begin{center}
\begin{tabularx}{\linewidth}{|l|X|l|}\hline
\rowcolor{popblueL} \textbf{Stratégie} & \textbf{Procédés / Marqueurs} & \textbf{Exemple dissertation} \\ \hline
\textbf{Argument d'autorité} & Citer auteur, proverbe, loi, expert. \emph{Selon..., Comme l'écrit..., Le proverbe dit...} & \og Comme l'affirme Cheikh Anta Diop, "l'histoire est le témoin des temps". \fg \\ \hline
\textbf{Argument par l'exemple} & Fait précis, chiffre, scène, personnage. \emph{Par exemple, ainsi, citons le cas de...} & \og Ainsi, la tontine villageoise finance des tracteurs. \fg \\ \hline
\textbf{Argument logique} & Causalité, conséquence, analogie. \emph{Parce que, donc, par conséquent, dès lors...} & \og Puisque l'école instruit, elle libère. \fg \\ \hline
\textbf{Concession / Réfutation} & Admettre l'argument adverse pour mieux le tourner. \emph{Certes... mais / Cependant / Néanmoins / Il est vrai que... toutefois...} & \og Certes, la tradition peut freiner ; toutefois elle évite le chaos. \fg \\ \hline
\end{tabularx}
\end{center}
\end{propriete}

\begin{methode}[Rédiger un paragraphe argumentatif type "Bac"]
\begin{enumerate}
\item Ouvrir par un \textbf{connecteur d'enchaînement} (\emph{D'abord, En premier lieu, Par ailleurs, En outre}).
\item Énoncer l'\textbf{idée directrice} (phrase nominale ou verbale claire).
\item Développer : \textbf{Argument} (connecteur \emph{car, puisque, en effet}) + \textbf{Exemple précis} (connecteur \emph{ainsi, par exemple, c'est le cas de}).
\item Fermer par une \textbf{phrase de conclusion partielle} (connecteur \emph{donc, par conséquent, en définitive}) qui relie à la problématique.
\item \textbf{Relire} : ponctuation, accords, vocabulaire soutenu, ton objectif.
\end{enumerate}
\end{methode}

\begin{exoresolu}[Rédaction d'un paragraphe de la partie II (Antithèse)]
Énoncé. Rédigez le paragraphe correspondant à l'argument : \og La solidarité traditionnelle finance le développement moderne \fg (Plan : Tradition / Progrès).
\tcblower
Solution détaillée.

\textbf{Paragraphe rédigé :}
\og \textbf{En outre}, la tradition, loin d'être un boulet, fournit au progrès des leviers financiers endogènes. \textbf{En effet}, la solidarité mécanique qui unit les membres d'une lignée ou d'un village se mue aujourd'hui en capital d'investissement grâce à la tontine. \textbf{Ainsi}, dans de nombreuses localités du Cameroun (Bafoussam, Foumban, Yaoundé), les tontines hebdomadaires ou mensuelles permettent l'acquisition de motos-taxis, de pompes à eau ou le financement des frais de scolarité universitaire des enfants du groupe. \textbf{Par conséquent}, cet épargne collective, héritée de la coutume, devient le moteur d'une modernité maîtrisée et partagée. \fg

\textbf{Analyse des marqueurs :}
\begin{itemize}
\item \emph{En outre} (enchainement parties) ; \emph{En effet} (explication) ; \emph{Ainsi} (illustration) ; \emph{Par conséquent} (conclusion logique).
\item Exemple géolocalisé (Cameroun), concret (moto, pompe, scolarité).
\item Phrases correctement ponctuées, vocabulaire précis (\emph{endogène, solidarité mécanique, capital d'investissement}).
\end{itemize}
\end{exoresolu}

\courssub{Communication : Situations et Conditions de l'énonciation}

\begin{aretenir}[Situation de communication de la dissertation]
\begin{center}
\begin{tabularx}{\linewidth}{|X|X|}\hline
\rowcolor{popblueL} \textbf{Composante} & \textbf{Réalité pour le candidat au Probatoire} \\ \hline
\textbf{Émetteur} & Candidat (élève de 1re Technique), locuteur responsable, objectif, cultivé. \\ \hline
\textbf{Récepteur} & Correcteur (professeur), attend : rigueur méthodique, culture générale, langue maîtrisée. \\ \hline
\textbf{Message} & Dissertation structurée (Introduction, Développement, Conclusion), argumentée, illustrée. \\ \hline
\textbf{Code} & Français standard soutenu (syntaxe correcte, vocabulaire précis, pas de familier). \\ \hline
\textbf{Contexte} & Épreuve certificative, temps limité (3h-4h), anonymat (numéro), enjeu diplôme. \\ \hline
\textbf{Contact} & Copie papier (lisibilité, propreté, paragraphes visibles, marges). \\ \hline
\end{tabularx}
\end{center}
\textbf{Conséquence} : Adaptez le ton (ni familier, ni agressif), soignez la présentation (saut de ligne pour chaque partie/sous-partie), gérez le temps (brouillon 45 min / rédaction 2h15 / relecture 15 min).
\end{aretenir}

\courssub{Intégration, Évaluation \& Remédiation}

\begin{exercice}[Entraînement type Probatoire - Sujet 1]
\textbf{Sujet :} \emph{« Dans quelle mesure l'exode rural menace-t-il la cohésion sociale au Cameroun ? »}
\begin{enumerate}
\item Analysez le sujet (mots-clés, type, définitions, enjeu).
\item Formulez la problématique principale et deux questions secondaires.
\item Faites un brainstorming classé (Thèse / Antithèse) avec 2 arguments + 1 exemple précis par argument.
\item Élaborez le plan détaillé complet (Introduction, I, II, Transition, Conclusion avec ouverture).
\end{enumerate}
\end{exercice}

\begin{exercice}[Entraînement type Probatoire - Sujet 2]
\textbf{Sujet :} \emph{« Le dialogue des cultures est-il une condition de la paix en Afrique ? »}
\begin{enumerate}
\item Identifiez la tension et proposez une problématique ouverte.
\item Choisissez le type de plan justifié (dialectique / analytique / thématique).
\item Rédigez l'introduction complète (amorce, définitions, problématique, annonce).
\item Rédigez le premier paragraphe de la partie I (Thèse) en respectant la structure argumentative et la ponctuation forte.
\end{enumerate}
\end{exercice}

\begin{exercice}[Remédiation : Correction de copies "pièges"]
Corrigez les erreurs majeures dans ces extraits de copies d'élèves :
\begin{enumerate}
\item \og \textbf{Sujet :} La femme au village. \textbf{Introduction :} La femme est importante. Elle travaille. Elle élève les enfants. \textbf{Problématique :} La femme est-elle importante ? \textbf{Plan :} I. Elle travaille. II. Elle élève. \fg
\item \og \textbf{Développement :} La tradition c'est bien. Par exemple mon grand-père. Il dit que c'est bien. Donc c'est bien. \fg
\item \og \textbf{Conclusion :} En conclusion, j'ai dit que la tradition c'est bien mais aussi pas bien. Merci. \fg
\end{enumerate}
\textit{Indications :} Hors-sujet (sujet mal analysé), absence de définitions, problématique fermée/répétitive, plan thématique au lieu de dialectique, arguments sans exemples valides, style oral, conclusion vide.
\end{exercice}

\begin{corrige}[Exercice Remédiation - Points clés]
\begin{enumerate}
\item \textbf{Erreurs :} Sujet trop large (pas de thèse), définitions absentes, problématique fermée (oui/non) et répétitive, plan descriptif (pas argumentatif), pas d'exemples littéraires/sociologiques.
\item \textbf{Correction type :} Sujet recentré \og Place et rôle de la femme rurale face à la modernité \fg. Mots-clés : femme, rural, modernité. Problématique : \og Dans quelle mesure la modernité transforme-t-elle le statut de la femme rurale ? \fg Plan dialectique : I. Pesanteurs traditionnelles / II. Leviers d'émancipation / III. Nouveaux équilibres.
\item \textbf{Erreurs :} Argument d'autorité familial non recevable (subjectif), raisonnement circulaire (\emph{c'est bien car c'est bien}), absence de connecteurs, ponctuation défaillante.
\item \textbf{Correction :} Argument logique (cohésion sociale) + Exemple précis (tontine, résolution conflits par chef) + Connecteurs (\emph{en effet, ainsi}).
\item \textbf{Erreurs :} Bilan absent, réponse à la problématique absente, ouverture absente, ton familier (\emph{Merci}).
\item \textbf{Correction :} Bilan synthétique des 2 parties $\rightarrow$ Réponse nuancée $\rightarrow$ Ouverture sur rôle de l'école technique.
\end{enumerate}
\end{corrige}

\begin{attention}[Les 5 erreurs fatales au Probatoire / Baccalauréat technique]
\begin{enumerate}
\item \textbf{Hors-sujet radical} : Ne pas définir les mots-clés $\rightarrow$ dissertation sur un thème voisin. \textbf{Parade :} 10 min d'analyse obligatoire au brouillon.
\item \textbf{Problématique "coquille vide"} : Recopier le sujet avec un point d'interrogation (\og La tradition freine-t-elle le progrès ? \fg). \textbf{Parade :} Formuler une tension (\og Dans quelle mesure... \fg).
\item \textbf{Plan "squelette vide"} : Parties sans arguments, arguments sans exemples précis (œuvre, lieu, chiffre, proverbe). \textbf{Parade :} Exigence 2 arguments + 1 exemple par sous-partie (A, B).
\item \textbf{Déséquilibre flagrant} : Partie I = 1 page, Partie II = 3 lignes. \textbf{Parade :} Même nombre d'arguments (2 ou 3) par grande partie.
\item \textbf{Négligence linguistique} : Phrases sans verbe, majuscules oubliées, points remplacés par des virgules, fautes d'accords systématiques. \textbf{Parade :} 15 min de relecture ciblée (ponctuation forte, accords sujets/verbes, homophones).
\end{enumerate}
\end{attention}

\begin{savaistu}[Le Lion et la Perle : un laboratoire pour la dissertation]
Wole Soyinka (Nigeria, Prix Nobel 1986) met en scène dans \emph{Le Lion et la Perle} (1963) le choc frontal entre \textbf{Baroka} (le Bale, tradition rusée, pouvoir coutumier) et \textbf{Lakunle} (instituteur, modernité mimétique, maladroite). \textbf{Sidi} (la Perle) est l'enjeu : elle hésite entre les deux mondes.
\begin{itemize}
\item \textbf{Thèse (Tradition frein) :} Lakunle veut supprimer la dot (\og coutume barbare \fg), Baroka use de la ruse traditionnelle pour séduire Sidi et bloquer le "progrès" (chemin de fer).
\item \textbf{Antithèse (Tradition socle) :} La ruse de Baroka (sagesse des anciens) triomphe de la naïveté moderne de Lakunle. Sidi choisit la sécurité coutumière.
\item \textbf{Synthèse :} La pièce ne tranche pas ; elle montre l'adaptation de la tradition (Baroka moderne dans son palais : timbre, machine à coudre) face à une modernité déracinée.
\end{itemize}
\textbf{Ressource inépuisable} pour illustrer : tradition/modernité, condition féminine, pouvoir, ruse/force, langue/identité.
\end{savaistu}

\newpage

\coursec{Exercices}

\courssub{Je vérifie mes connaissances}

\begin{exercice}[QCM : les étapes de la dissertation]
Pour chaque question, choisis la bonne réponse (une seule possible).
\begin{enumerate}
\item La première étape du travail sur un sujet de dissertation est :
\begin{enumerate}
\item la recherche des idées ; \item l'analyse du sujet ; \item la rédaction de l'introduction ; \item l'élaboration du plan détaillé.
\end{enumerate}
\item La problématique est :
\begin{enumerate}
\item un résumé du sujet ; \item la question centrale qui naît de l'analyse du sujet ; \item la conclusion du devoir ; \item une citation d'auteur.
\end{enumerate}
\item Un plan dialectique se présente sous la forme :
\begin{enumerate}
\item thèse / antithèse / synthèse ; \item cause / conséquence ; \item description / narration ; \item problème / solution.
\end{enumerate}
\item Dans le plan détaillé, chaque partie doit contenir au minimum :
\begin{enumerate}
\item un seul argument ; \item deux arguments illustrés d'exemples ; \item cinq exemples ; \item une citation obligatoire.
\end{enumerate}
\end{enumerate}
\end{exercice}

\begin{exercice}[Vrai ou faux : le texte argumentatif]
Indique si chaque affirmation est vraie ou fausse, puis justifie ta réponse en une phrase.
\begin{enumerate}
\item Un texte argumentatif vise uniquement à distraire le lecteur.
\item « Cependant », « néanmoins » et « toutefois » sont des connecteurs exprimant l'opposition.
\item Dans une dissertation, les exemples peuvent être inventés de toutes pièces, sans lien avec la réalité.
\item La phrase simple ne contient qu'un seul verbe conjugué.
\end{enumerate}
\end{exercice}

\begin{exercice}[Définitions]
Définis en deux ou trois phrases chacun des termes suivants, tels qu'ils s'emploient dans la méthode de la dissertation :
\begin{enumerate}
\item la problématique ;
\item l'argument ;
\item le connecteur logique ;
\item le plan détaillé.
\end{enumerate}
\end{exercice}

\begin{exercice}[Ponctuation et phrase]
\begin{enumerate}
\item Réécris la phrase suivante en plaçant correctement la ponctuation : « le développement durable est il une contrainte pour l'industrialisation du cameroun »
\item Indique si chacune des phrases suivantes est une phrase simple ou une phrase complexe, en justifiant par le nombre de verbes conjugués :
\begin{enumerate}
\item « Baroka ruse avec Sidi. »
\item « Lakunle, qui rêve de modernité, refuse de payer le prix de la dot. »
\end{enumerate}
\end{enumerate}
\end{exercice}

\courssub{Je m'entraîne}

\begin{exercice}[Analyse d'un sujet et problématique]
On donne le sujet suivant : « Le téléphone portable : un progrès pour la jeunesse camerounaise ? »
\begin{enumerate}
\item Dégage le thème du sujet.
\item Délimite le sujet (domaine, espace, catégorie de personnes concernées).
\item Relève les termes importants et définis-les brièvement.
\item Formule la problématique sous la forme d'une question unique et claire.
\end{enumerate}
\end{exercice}

\begin{exercice}[Réécriture syntaxique]
Voici une suite de phrases simples extraites d'un rapport de stage technique : « L'entreprise dispose de machines modernes. Les ouvriers ne sont pas formés. Les machines restent souvent inutilisées. La production baisse. Le chef d'atelier propose une formation. »
\begin{enumerate}
\item Transforme cette suite en un paragraphe fluide de deux ou trois phrases complexes, en utilisant la subordination (qui, parce que, bien que, si) et la coordination (mais, donc, or).
\item Souligne les connecteurs logiques que tu as employés et précise la relation qu'ils expriment (cause, conséquence, opposition).
\end{enumerate}
\end{exercice}

\begin{exercice}[Analyse de connecteurs]
Dans un éditorial de \emph{Cameroon Tribune} consacré à la formation professionnelle, on relève les connecteurs suivants : « d'abord », « ensuite », « en effet », « cependant », « néanmoins », « toutefois », « donc », « par conséquent », « car », « enfin », « de plus », « c'est pourquoi », « en outre », « ainsi », « pourtant ».
\begin{enumerate}
\item Classe ces quinze connecteurs selon la relation logique qu'ils expriment (addition, cause, conséquence, opposition, ordre/énumération).
\item Explique la nuance de sens apportée par « cependant », « néanmoins » et « toutefois ». Peut-on toujours les intervertir ? Justifie ta réponse par un exemple de phrase.
\end{enumerate}
\end{exercice}

\begin{exercice}[Recherche d'idées et plan]
On donne le sujet : « Faut-il préférer la formation technique à la formation générale ? »
\begin{enumerate}
\item Note au brouillon au moins six idées (arguments ou exemples) que ce sujet t'inspire.
\item Regroupe ces idées en deux ou trois grandes parties.
\item Indique le type de plan le plus adapté (thématique, dialectique ou analytique) et justifie ton choix.
\item Rédige le plan détaillé : chaque partie avec son titre, ses arguments et les exemples prévus.
\end{enumerate}
\end{exercice}

\courssub{Je me prépare au Bac}

\begin{exercice}[Sujet type Probatoire : le développement durable]
Sujet : « Le développement durable, une contrainte pour l'industrialisation du Cameroun ? »
\begin{enumerate}
\item Analyse le sujet : dégage le thème, délimite-le, définis les termes importants (« développement durable », « industrialisation », « contrainte »).
\item Formule la problématique.
\item Recherche les idées : relève la tension entre les exigences économiques (emplois, croissance, usines) et les exigences écologiques (protection des forêts, de l'eau, des sols). Tu pourras mobiliser des exemples locaux : le barrage de Lom Pangar, les plantations de la SOCAPALM, l'exploitation forestière dans la région de l'Est.
\item Élabore un plan dialectique détaillé (thèse / antithèse / synthèse) : pour chaque partie, donne le titre, deux arguments et un exemple précis.
\item Rédige l'introduction complète du devoir (amorce, présentation du sujet, problématique, annonce du plan).
\end{enumerate}
\end{exercice}

\begin{exercice}[Cas pratique : \emph{Le Lion et la Perle} de Wole Soyinka]
Question : « Baroka incarne-t-il une modernité africaine authentique ? »
\begin{enumerate}
\item Rappelle qui est Baroka et quel rôle il joue dans le village d'Ilujinle.
\item Rédige un paragraphe argumentatif complet (environ quinze lignes) comprenant : une thèse clairement énoncée, deux arguments développés et des exemples textuels précis (l'usage de l'ironie et des proverbes par Baroka, la scène de la poste et du téléphone, son rapport à la machine et aux timbres).
\item Veille à employer au moins trois connecteurs logiques différents et à soigner la ponctuation.
\item Propose ensuite le plan détaillé d'une dissertation entière sur cette question, en annonçant le type de plan choisi et en le justifiant.
\end{enumerate}
\end{exercice}

\begin{exercice}[Sujet de Bac blanc (durée conseillée : 4 heures)]
Sujet : « L'innovation technologique menace-t-elle les savoir-faire traditionnels ? »
Consignes :
\begin{enumerate}
\item Analyse le sujet et formule la problématique : tu opposeras notamment la logique de la « destruction créatrice » (le nouveau remplace l'ancien) à celle de la préservation du patrimoine culturel et technique.
\item Recherche les idées et construis un plan thématique détaillé en trois parties, par exemple : l'agriculture (machines et tracteurs face à la main-d'œuvre et aux techniques culturales locales), l'artisanat (production industrielle face au tissage, à la poterie, à la forge traditionnelle), la formation (études et diplômes face à l'apprentissage traditionnel auprès d'un maître artisan).
\item Pour chaque partie, développe au moins deux arguments illustrés d'exemples concrets tirés de la réalité camerounaise ou africaine.
\item Rédige le devoir complet : introduction (avec problématique et annonce du plan), développement en trois parties avec transitions, conclusion (bilan et ouverture).
\item Soigne la langue : phrases simples et phrases complexes bien construites, ponctuation correcte, connecteurs logiques variés.
\end{enumerate}
\end{exercice}

\newpage

\coursec{Corrigés des exercices}

\coursec{Leçon 02 : Produire le plan détaillé d'une dissertation}

\courssub{Exercice 1 — QCM : les étapes de la dissertation}
\begin{corrige}[Exercice 1 — QCM : les étapes de la dissertation]
\begin{enumerate}
\item \textbf{Réponse b : l'analyse du sujet.} Avant toute recherche d'idées, il faut lire attentivement le sujet, dégager le thème, délimiter le sujet et définir les termes importants. Rédiger l'introduction ou le plan sans cette étape expose au hors-sujet.
\item \textbf{Réponse b : la question centrale qui naît de l'analyse du sujet.} La problématique n'est ni un résumé ni une citation : c'est l'interrogation directrice que le devoir entier devra résoudre.
\item \textbf{Réponse a : thèse / antithèse / synthèse.} Le plan dialectique confronte deux positions opposées avant de les dépasser dans une synthèse.
\item \textbf{Réponse b : deux arguments illustrés d'exemples.} Un seul argument ne suffit pas à étayer une partie ; chaque argument doit être suivi d'un exemple précis qui le rend convaincant.
\end{enumerate}
\textbf{Point de vigilance :} à l'examen, l'ordre des étapes est impératif : analyse du sujet $\rightarrow$ problématique $\rightarrow$ recherche des idées $\rightarrow$ plan détaillé $\rightarrow$ rédaction.
\end{corrige}

\courssub{Exercice 2 — Vrai ou faux : le texte argumentatif}
\begin{corrige}[Exercice 2 — Vrai ou faux : le texte argumentatif]
\begin{enumerate}
\item \textbf{Faux.} Un texte argumentatif vise à \cle{convaincre} (par la raison) ou à \cle{persuader} (par l'émotion) le lecteur en défendant une thèse à l'aide d'arguments et d'exemples ; distraire est la fonction du texte narratif ou ludique.
\item \textbf{Vrai.} « Cependant », « néanmoins » et « toutefois » sont bien des connecteurs d'\cle{opposition} (ou de concession) : ils introduisent une idée qui s'oppose à la précédente.
\item \textbf{Faux.} Les exemples doivent être \cle{plausibles} et liés à la réalité (faits observés, œuvres lues, événements historiques ou locaux) ; un exemple invraisemblable affaiblit l'argumentation au lieu de la soutenir.
\item \textbf{Vrai.} La \cle{phrase simple} ne contient qu'une seule proposition, donc un seul verbe conjugué ; dès qu'il y a deux verbes conjugués, on a une \cle{phrase complexe} (qui regroupe les phrases composées — coordination, juxtaposition — et les phrases à subordination).
\end{enumerate}
\end{corrige}

\courssub{Exercice 3 — Définitions}
\begin{corrige}[Exercice 3 — Définitions]
\begin{enumerate}
\item \textbf{La problématique} est la question centrale qui naît de l'analyse du sujet. Elle exprime le problème intellectuel que le devoir doit résoudre et oriente toute l'argumentation. Elle se formule en une phrase interrogative claire et unique.
\item \textbf{L'argument} est une raison avancée pour défendre ou combattre une thèse. Il doit être clair, pertinent et suivi d'un exemple qui l'illustre et le rend convaincant.
\item \textbf{Le connecteur logique} est un mot ou une locution (« d'abord », « cependant », « donc », « en effet ») qui relie les idées entre elles. Il explicite la relation logique : addition, opposition, cause, conséquence ou énumération.
\item \textbf{Le plan détaillé} est l'organisation écrite du devoir avant la rédaction : il présente les grandes parties avec leurs titres, et sous chaque partie les arguments, les exemples et les transitions prévus.
\end{enumerate}
\end{corrige}

\courssub{Exercice 4 — Ponctuation et phrase}
\begin{corrige}[Exercice 4 — Ponctuation et phrase]
\begin{enumerate}
\item Phrase correctement ponctuée : « Le développement durable est-il une contrainte pour l'industrialisation du Cameroun ? » \\
Justification : majuscule en début de phrase ; inversion du sujet avec trait d'union (\emph{est-il}) car c'est une interrogation directe ; point d'interrogation final ; majuscule au nom propre « Cameroun ».
\item \begin{enumerate}
\item « Baroka ruse avec Sidi. » : \textbf{phrase simple}, car elle ne contient qu'un seul verbe conjugué (\emph{ruse}), donc une seule proposition.
\item « Lakunle, qui rêve de modernité, refuse de payer le prix de la dot. » : \textbf{phrase complexe}, car elle contient deux verbes conjugués (\emph{rêve}, \emph{refuse}) répartis en deux propositions : une principale (« Lakunle refuse de payer le prix de la dot ») et une subordonnée relative (« qui rêve de modernité »). Attention : « payer » est à l'infinitif, il ne compte pas comme verbe conjugué.
\end{enumerate}
\end{enumerate}
\textbf{Point de vigilance :} ne compter que les verbes \emph{conjugués} ; les infinitifs et participes ne créent pas de proposition.
\end{corrige}

\courssub{Exercice 5 — Analyse d'un sujet et problématique}
\begin{corrige}[Exercice 5 — Analyse d'un sujet et problématique]
\begin{enumerate}
\item \textbf{Thème :} les nouvelles technologies de communication et leurs effets sur la jeunesse.
\item \textbf{Délimitation :} domaine technologique et social ; espace : le Cameroun ; catégorie concernée : la jeunesse (adolescents et jeunes adultes). Le sujet n'interroge donc ni l'usage professionnel des adultes ni la situation d'autres pays.
\item \textbf{Termes importants :}
\begin{itemize}
\item \emph{téléphone portable} : appareil de communication mobile, donnant aussi accès à Internet, aux réseaux sociaux, au paiement mobile ;
\item \emph{progrès} : ce qui marque une amélioration, une avancée bénéfique ; le mot est mis en question par le point d'interrogation du sujet ;
\item \emph{jeunesse camerounaise} : les jeunes générations du Cameroun, élèves, étudiants, apprentis, jeunes travailleurs.
\end{itemize}
\item \textbf{Problématique :} « Le téléphone portable constitue-t-il réellement un progrès pour la jeunesse camerounaise, ou ses dangers l'emportent-ils sur ses bienfaits ? »
\end{enumerate}
\textbf{Point de vigilance :} la problématique doit reprendre la tension du sujet (progrès \emph{ou} danger) sans y répondre d'avance.
\end{corrige}

\courssub{Exercice 6 — Réécriture syntaxique}
\begin{corrige}[Exercice 6 — Réécriture syntaxique]
\begin{enumerate}
\item \textbf{Paragraphe possible :} « Bien que l'entreprise dispose de machines modernes, celles-ci restent souvent inutilisées \textbf{parce que} les ouvriers ne sont pas formés. La production baisse \textbf{donc}, \textbf{c'est pourquoi} le chef d'atelier propose une formation. » (Autre formulation acceptée : « L'entreprise dispose de machines modernes, \textbf{mais} elles restent souvent inutilisées \textbf{car} les ouvriers ne sont pas formés ; la production baisse \textbf{par conséquent}, et le chef d'atelier propose une formation. »)
\item \textbf{Connecteurs et relations exprimées :}
\begin{center}
\begin{tabularx}{\linewidth}{|l|X|}\hline
\rowcolor{popblueL} \textbf{Connecteur} & \textbf{Relation logique} \\ \hline
bien que & concession / opposition (concession) \\ \hline
parce que (car) & cause \\ \hline
mais & opposition \\ \hline
donc / par conséquent & conséquence \\ \hline
c'est pourquoi & conséquence \\ \hline
\end{tabularx}
\end{center}
\end{enumerate}
\textbf{Point de vigilance :} la réécriture doit supprimer les répétitions (« les machines » repris par « celles-ci » ou « elles ») et conserver le sens du texte de départ.
\end{corrige}

\courssub{Exercice 7 — Analyse de connecteurs}
\begin{corrige}[Exercice 7 — Analyse de connecteurs]
\begin{enumerate}
\item \textbf{Classement des quinze connecteurs :}
\begin{center}
\begin{tabularx}{\linewidth}{|l|X|}\hline
\rowcolor{popblueL} \textbf{Relation logique} & \textbf{Connecteurs} \\ \hline
Addition & de plus, en outre \\ \hline
Cause & car \\ \hline
Conséquence & donc, par conséquent, c'est pourquoi, ainsi \\ \hline
Opposition & cependant, néanmoins, toutefois, pourtant \\ \hline
Ordre / énumération & d'abord, ensuite, enfin \\ \hline
Explication / confirmation & en effet \\ \hline
\end{tabularx}
\end{center}
Remarque : « en effet » n'entre dans aucune des cinq catégories proposées au sens strict ; il exprime l'\emph{explication} ou la \emph{confirmation} d'une idée précédente. « Ainsi » peut aussi, selon le contexte, marquer la manière ou la comparaison.
\item \textbf{Nuance entre « cependant », « néanmoins » et « toutefois » :} les trois marquent l'opposition, mais « néanmoins » et « toutefois » introduisent une \cle{concession} plus forte : la seconde idée l'emporte malgré la première (« Il pleuvait ; néanmoins, le match a eu lieu »). « Cependant » signale une simple opposition ou restriction. On ne peut donc pas toujours les intervertir : dans « Il est pauvre, \emph{cependant} il est généreux », le contraste est direct ; dans « La tâche est dure ; \emph{néanmoins}, il faut l'accomplir », la concession s'impose et « cependant » affaiblirait le sens.
\end{enumerate}
\end{corrige}

\courssub{Exercice 8 — Recherche d'idées et plan}
\begin{corrige}[Exercice 8 — Recherche d'idées et plan]
\begin{enumerate}
\item \textbf{Idées au brouillon (exemples) :}
\begin{itemize}
\item la formation technique donne rapidement un métier concret (mécanicien, électricien, menuisier) ;
\item elle répond aux besoins du marché de l'emploi camerounais (BTP, agro-industrie) ;
\item elle favorise l'auto-emploi et l'entrepreneuriat des jeunes ;
\item la formation générale développe la culture, l'esprit critique et l'expression ;
\item elle ouvre la voie aux longues études universitaires (médecine, droit, enseignement) ;
\item un technicien sans culture générale peut être limité dans son évolution ; un bachelier général sans débouché peut rester chômeur.
\end{itemize}
\item \textbf{Regroupement :} partie 1 : les atouts de la formation technique ; partie 2 : les atouts de la formation générale ; partie 3 : la complémentarité nécessaire des deux.
\item \textbf{Type de plan :} le plan \cle{dialectique} (thèse / antithèse / synthèse) est le plus adapté, car le sujet oppose explicitement deux formations (« faut-il préférer... ») et invite à dépasser l'opposition.
\item \textbf{Plan détaillé :}
\begin{itemize}
\item \textbf{I. La formation technique, une voie d'insertion rapide.} Arg. 1 : elle transmet un savoir-faire directement exploitable (ex. : un titulaire d'un CAP en électricité trouve du travail à Douala). Arg. 2 : elle répond aux besoins de l'économie nationale (ex. : pénurie de techniciens dans le BTP et l'agro-industrie).
\item \textbf{II. Mais la formation générale reste irremplaçable.} Arg. 1 : elle forme l'esprit critique et la capacité d'adaptation (ex. : le lycéen généraliste peut se réorienter plus facilement). Arg. 2 : elle est indispensable aux professions exigeant de longues études (ex. : médecins, ingénieurs, magistrats).
\item \textbf{III. Synthèse : deux formations complémentaires.} Arg. 1 : le pays a besoin des deux profils (ex. : un chantier réunit ingénieur et techniciens). Arg. 2 : l'idéal est une formation technique enrichie de culture générale (ex. : les passerelles entre filières techniques et universitaires).
\end{itemize}
\end{enumerate}
\end{corrige}

\courssub{Exercice 9 — Sujet type Probatoire : le développement durable}
\begin{corrige}[Exercice 9 — Sujet type Probatoire : le développement durable]
\textbf{Barème indicatif (sur 20 \%) :} analyse du sujet et problématique : 4 pts ; recherche des idées : 4 pts ; plan détaillé : 6 pts ; introduction rédigée : 4 pts ; langue et présentation : 2 pts.
\begin{enumerate}
\item \textbf{Analyse du sujet.} Thème : l'environnement et l'économie. Délimitation : espace, le Cameroun ; domaine, l'industrie et l'écologie. Définitions : le \emph{développement durable} est un développement qui satisfait les besoins du présent sans compromettre ceux des générations futures ; l'\emph{industrialisation} est la transformation d'une économie par la création d'usines et la production de masse ; la \emph{contrainte} est ce qui limite, freine ou empêche une action.
\item \textbf{Problématique :} « Le respect de l'environnement freine-t-il nécessairement l'industrialisation du Cameroun, ou peut-il au contraire l'accompagner ? »
\item \textbf{Recherche des idées :} exigences économiques (emplois industriels, croissance, transformation locale des matières premières) contre exigences écologiques (protection des forêts de l'Est, de l'eau, des sols). Exemples : le barrage de Lom Pangar (énergie mais impacts sur l'environnement), les plantations de la SOCAPALM (emplois mais déforestation et conflits fonciers), l'exploitation forestière dans la région de l'Est (devises mais disparition du couvert végétal).
\item \textbf{Plan dialectique détaillé :}
\begin{itemize}
\item \textbf{I. Thèse : le développement durable apparaît comme une contrainte.} Arg. 1 : les normes écologiques renchérissent les coûts de production (ex. : traitement des rejets d'usine). Arg. 2 : la protection des forêts limite l'exploitation des ressources (ex. : quotas d'abattage dans la région de l'Est).
\item \textbf{II. Antithèse : mais l'industrialisation sauvage est plus coûteuse encore.} Arg. 1 : la dégradation des sols et de l'eau détruit l'agriculture, base de l'économie (ex. : contestations autour des plantations de la SOCAPALM). Arg. 2 : une énergie propre est une condition de l'industrie (ex. : le barrage de Lom Pangar alimente le réseau électrique national).
\item \textbf{III. Synthèse : vers une industrialisation durable.} Arg. 1 : l'écologie peut créer des emplois nouveaux (ex. : recyclage, énergies solaires). Arg. 2 : transformer localement en respectant les normes valorise les produits camerounais (ex. : certification du bois pour l'exportation).
\end{itemize}
\item \textbf{Introduction possible :} « Le Cameroun aspire à devenir un pays émergent, et cette ambition passe par l'industrie. Or, chaque usine implantée, chaque forêt exploitée pose la question du prix écologique de la croissance : c'est tout le débat sur le développement durable. Le sujet invite donc à se demander si ce développement constitue une contrainte pour l'industrialisation du pays. En d'autres termes, le respect de l'environnement freine-t-il nécessairement l'industrie camerounaise ? Nous verrons d'abord que le développement durable peut apparaître comme une contrainte, puis qu'une industrialisation irrespectueuse de la nature se révèle plus coûteuse encore, avant de montrer qu'une industrialisation durable est possible. »
\end{enumerate}
\textbf{Points de vigilance :} la synthèse ne doit pas être un simple résumé des deux premières parties mais un dépassement ; l'introduction doit contenir les quatre éléments exigés (amorce, présentation du sujet, problématique, annonce du plan) ; chaque argument doit être suivi d'un exemple précis.
\end{corrige}

\courssub{Exercice 10 — Cas pratique : \emph{Le Lion et la Perle} de Wole Soyinka}
\begin{corrige}[Exercice 10 — Cas pratique : \emph{Le Lion et la Perle} de Wole Soyinka]
\textbf{Barème indicatif (sur 20 \%) :} présentation de Baroka : 2 pts ; paragraphe argumentatif (thèse, arguments, exemples, connecteurs, ponctuation) : 10 pts ; plan détaillé justifié : 6 pts ; langue : 2 pts.
\begin{enumerate}
\item \textbf{Baroka} est le chef (le \emph{Bale}) du village d'Ilujinle, surnommé « le Lion ». Polygame rusé et attaché aux traditions yoruba, il s'oppose à Lakunle, l'instituteur « moderniste », dans la conquête de Sidi, « la Perle ». Il incarne la ruse et la force de la tradition face à une modernité superficielle.
\item \textbf{Paragraphe argumentatif possible :} « Baroka incarne, à sa manière, une modernité africaine authentique, car il sait faire dialoguer tradition et progrès au lieu de les opposer. D'abord, le chef d'Ilujinle n'est pas un ennemi du progrès : il comprend l'usage des objets modernes, comme le montre la scène de la poste et du téléphone, où il révèle qu'il maîtrise mieux que Lakunle les codes du monde moderne. En effet, là où l'instituteur récite une modernité apprise dans les livres, Baroka l'utilise avec intelligence, par exemple en faisant imprimer l'image de Sidi sur les timbres, détournant ainsi la machine au service de sa stratégie. Ensuite, son attachement aux proverbes et à l'ironie n'est pas un signe d'archaïsme : c'est une arme rhétorique qui lui permet de démasquer les prétentions de son rival. Cependant, cette modernité reste ambiguë, car elle sert aussi la ruse et la domination du vieux chef. Baroka prouve donc que la vraie modernité africaine ne consiste pas à copier l'Occident, mais à intégrer le progrès à la sagesse traditionnelle. » (Connecteurs employés : \emph{d'abord}, \emph{en effet}, \emph{ensuite}, \emph{cependant}, \emph{donc}.)
\item \textbf{Plan détaillé — plan dialectique} (justifié par la forme interrogative du sujet, qui appelle une confrontation puis un dépassement) :
\begin{itemize}
\item \textbf{I. Baroka semble incarner la tradition contre la modernité.} Arg. 1 : il défend les coutumes du village (ex. : la polygamie, le paiement de la dot). Arg. 2 : il a freiné certains projets modernes (ex. : son opposition au passage de la voie ferrée à Ilujinle).
\item \textbf{II. Mais il comprend et utilise la modernité mieux que Lakunle.} Arg. 1 : il maîtrise les objets modernes (ex. : la scène de la poste et du téléphone). Arg. 2 : il détourne la technique à son profit (ex. : les timbres à l'effigie de Sidi, la machine).
\item \textbf{III. Synthèse : une modernité africaine authentique, fondée sur l'adaptation.} Arg. 1 : la modernité véritable intègre le progrès à la culture locale (ex. : l'usage des proverbes comme arme d'argumentation). Arg. 2 : Soyinka critique la modernité d'imitation incarnée par Lakunle (ex. : le ridicule de l'instituteur qui refuse la dot sans comprendre le village).
\end{itemize}
\end{enumerate}
\textbf{Points de vigilance :} les exemples doivent renvoyer à des scènes précises de la pièce ; le paragraphe doit s'ouvrir sur la thèse et se fermer sur une conclusion partielle ; soigner la ponctuation des deux-points et des points-virgules.
\end{corrige}

\courssub{Exercice 11 — Sujet de Bac blanc : innovation technologique et savoir-faire traditionnels}
\begin{corrige}[Exercice 11 — Sujet de Bac blanc : innovation technologique et savoir-faire traditionnels]
\textbf{Barème indicatif (sur 20 \%) :} introduction (problématique + annonce du plan) : 4 pts ; développement (3 parties, arguments et exemples) : 10 pts ; conclusion (bilan + ouverture) : 3 pts ; langue, ponctuation, connecteurs : 3 pts.
\begin{enumerate}
\item \textbf{Analyse et problématique.} Thème : technique et culture. L'\emph{innovation technologique} désigne l'apparition de machines et de procédés nouveaux ; les \emph{savoir-faire traditionnels} sont les techniques transmises de génération en génération (tissage, poterie, forge, agriculture). Problématique : « L'innovation technologique condamne-t-elle nécessairement les savoir-faire traditionnels à disparaître, ou peut-elle les transformer et les valoriser ? »
\item \textbf{Plan thématique détaillé en trois parties :}
\begin{itemize}
\item \textbf{I. Dans l'agriculture, la machine bouscule les techniques locales.} Arg. 1 : le tracteur augmente les rendements mais réduit l'emploi de la main-d'œuvre (ex. : mécanisation des plantations dans la plaine des Mbo). Arg. 2 : les techniques culturales locales (rotation, jachère, variétés locales) risquent d'être abandonnées alors qu'elles protègent les sols (ex. : cultures vivrières traditionnelles de l'Ouest).
\item \textbf{II. Dans l'artisanat, l'industrie concurrence le fait main.} Arg. 1 : la production industrielle, moins chère, menace le tissage, la poterie et la forge traditionnelle (ex. : tissus industriels face au ndop et au kaba tissé main). Arg. 2 : mais le fait main garde une valeur culturelle et économique unique (ex. : artisanat d'art de Foumban recherché par les touristes).
\item \textbf{III. Dans la formation, le diplôme concurrence l'apprentissage auprès du maître artisan.} Arg. 1 : les études et les diplômes sont devenus la voie valorisée, délaissant l'apprentissage traditionnel (ex. : jeunes préférant les filières scolaires aux ateliers familiaux). Arg. 2 : pourtant l'apprentissage auprès d'un maître transmet des gestes que l'école n'enseigne pas (ex. : forgerons et sculpteurs formés en atelier dans la région de l'Extrême-Nord).
\end{itemize}
\item \textbf{Exemples concrets :} voir ci-dessus ; chaque argument est illustré d'une réalité camerounaise ou africaine précise.
\item \textbf{Éléments de rédaction.} \emph{Introduction :} amorce sur la place de la technique dans la vie quotidienne ; présentation du sujet ; problématique ; annonce des trois parties (agriculture, artisanat, formation). \emph{Développement :} chaque partie s'ouvre sur une phrase d'accroche, développe deux arguments avec exemples, et se relie à la suivante par une transition (ex. : « Ce que l'on observe dans les champs se vérifie aussi dans les ateliers »). \emph{Conclusion :} bilan — l'innovation menace les savoir-faire lorsqu'elle les ignore, mais elle peut les valoriser lorsqu'elle s'appuie sur eux ; ouverture possible : « Ne faut-il pas former des techniciens capables de moderniser les gestes traditionnels plutôt que de les effacer ? »
\item \textbf{Langue :} alterner phrases simples (pour la clarté) et phrases complexes (subordonnées avec \emph{qui}, \emph{parce que}, \emph{bien que}) ; varier les connecteurs (\emph{d'abord}, \emph{en effet}, \emph{cependant}, \emph{par conséquent}, \emph{enfin}) ; vérifier la ponctuation des deux-points, des points-virgules et des interrogations.
\end{enumerate}
\textbf{Points de vigilance :} respecter le plan thématique annoncé (ne pas glisser vers un plan dialectique non annoncé) ; chaque partie doit comporter au moins deux arguments illustrés ; la conclusion ne doit introduire aucun argument nouveau, seulement un bilan et une ouverture ; gérer le temps : environ 30 min pour l'analyse et le plan, 2 h 30 pour la rédaction, 30 min pour la relecture.
\end{corrige}

\newpage

\coursec{Fiche bilan}

\begin{popfigure}
\centering
\begin{tikzpicture}[
    mindmap,
    concept color=popblue,
    level 1 concept/.append style={level distance=4.5cm, sibling angle=360/6},
    level 2 concept/.append style={level distance=3cm},
    every node/.style={font=\small, text=popdark},
    grow cyclic,
    text width=2.8cm,
    align=center
]
\node[concept, text=white, font=\bfseries\small, align=center]{Leçon 02\\Dissertation \& Argumentation}
    child[concept color=poporange] { node[concept, text=white, font=\bfseries\tiny, align=center]{\textit{Le Lion et la Perle}\\Paratextes \& Lecture ouvroir} }
    child[concept color=popgreen] { node[concept, text=white, font=\bfseries\tiny, align=center]{Communication\\Situations \& Ponctuation} }
    child[concept color=poppurple] { node[concept, text=white, font=\bfseries\tiny, align=center]{Syntaxe\\Phrase simple / composée} }
    child[concept color=poppink] { node[concept, text=white, font=\bfseries\tiny, align=center]{Texte argumentatif\\Caractères \& Stratégies} }
    child[concept color=popgold] { node[concept, text=white, font=\bfseries\tiny, align=center]{Dissertation\\Analyse sujet $\to$ Problématique} }
    child[concept color=popblue] { node[concept, text=white, font=\bfseries\tiny, align=center]{Plan détaillé\\Recherche idées $\to$ Structure} };
\end{tikzpicture}
\legende{Carte mentale de la leçon 02 : articulation œuvre, langue, rhétorique et méthode dissertation.}
\end{popfigure}

\begin{bilan}

\courssub{Définitions clés}

\begin{itemize}
    \item \textbf{Paratexte} (Genette) : Ensemble des éléments entourant le texte (titre, préface, couverture, notes) qui guident la lecture.
    \item \textbf{Lecture ouvroir} : Première lecture active, annotée, pour saisir l'intrigue, les personnages, l'espace/temps et les impressions subjectives.
    \item \textbf{Situation de communication} : Configuration \cle{Émetteur} $\to$ \cle{Message} $\to$ \cle{Récepteur} via un \cle{Canal}, un \cle{Code}, dans un \cle{Contexte} (Jakobson).
    \item \textbf{Ponctuation faible} (, ; : ...) : Organise la phrase ; \textbf{forte} (. ? ! ...) : Clôture l'énoncé et marque l'intonation.
    \item \textbf{Phrase simple} : Une seule proposition (un verbe conjugué). \textbf{Phrase composée} : Au moins deux propositions liées par \cle{coordination}, \cle{subordination} ou \cle{juxtaposition}.
    \item \textbf{Texte argumentatif} : Vise à \cle{convaincre} (raison) ou \cle{persuader} (émotion). Structure : \cle{Thèse} -- \cle{Arguments} -- \cle{Exemples} -- \cle{Connecteurs logiques}.
    \item \textbf{Stratégies argumentatives} : Concession, réfutation, analogie, argument d'autorité, \textit{ad hominem}, exemple, chiffres/statistiques.
    \item \textbf{Analyse de sujet} : 1) Dépouiller les mots-clés, 2) Définir les termes, 3) Délimiter le champ, 4) Dégager la tension/paradoxe, 5) Formuler la \cle{problématique} (question ouverte, débattue, issue du sujet).
    \item \textbf{Plan détaillé} : Squelette complet de la dissertation. \cle{Introduction} (Amorce, Sujet posé, Problématique, Annonce du plan) -- \cle{Développement} (2 ou 3 parties \textbf{I, II, III} ; chacune 2--3 sous-parties \textbf{A, B, C} avec arguments + exemples précis) -- \cle{Conclusion} (Bilan, Ouverture).
\end{itemize}

\courssub{Méthodes express (pas à pas)}

\begin{methode}[Analyser un sujet de dissertation (5 étapes)]
    \begin{enumerate}
        \item \textbf{Surlineur} : Mots-clés, termes techniques, connecteurs, énoncé (citation, question, affirmation).
        \item \textbf{Dictionnaire} : Définir chaque mot-clé (sens littéral, connoté, opposés).
        \item \textbf{Frontières} : Définir l'espace (œuvre, époque, genre), le temps, le point de vue imposé ou libre.
        \item \textbf{Tension} : Chercher l'opposition, le paradoxe, la nuance (ex: « \textit{Le théâtre est-il seulement un divertissement ?} » $\to$ tension divertissement / engagement).
        \item \textbf{Problématique} : Rédiger \textbf{une question unique}, ouverte, qui appelle une réponse nuancée (ni oui, ni non). Éviter : « Qu'est-ce que... ? », « Pourquoi... ? » Préférer : « Dans quelle mesure... ? », « Comment... ? », « Faut-il... ? ».
    \end{enumerate}
\end{methode}

\begin{methode}[Rechercher et classer les idées (Brainstorming $\to$ Arborescence)]
    \begin{enumerate}
        \item \textbf{Explosion} : Noter \textbf{tout} ce qui vient (mots, citations, persos, scènes, concepts) sans trier (10--15 min).
        \item \textbf{Triage} : Grouper par affinités thématiques (couleurs, symboles). Éliminer le hors-sujet.
        \item \textbf{Hiérarchiser} : Identifier 2--3 \textbf{grands axes} (devenir parties I, II, III). Vérifier l'équilibre et la progression logique (dialectique : Thèse/Antithèse/Synthèse ; ou Thématique : Cause/Conséquence ; ou Analytique : Définition/Analyse/Évaluation).
        \item \textbf{Détailler} : Pour chaque sous-partie (A, B, C), écrire : \textbf{Idée directrice} + \textbf{Argument} + \textbf{Exemple précis (citation, scène, fait)} + \textbf{Connecteur}.
    \end{enumerate}
\end{methode}

\begin{methode}[Rédiger le plan détaillé (Modèle standard)]
    \begin{enumerate}
        \item \textbf{Introduction} : [Amorce] $\to$ [Sujet posé reformulé] $\to$ [Problématique] $\to$ [Annonce plan : « Nous verrons d'abord que... puis que... enfin... »].
        \item \textbf{Développement} :
        \begin{itemize}
            \item \textbf{I. Premier axe} (Idée force)
            \item A. Argument 1 + Exemple précis
            \item B. Argument 2 + Exemple précis
            \item (C. Argument 3 + Exemple précis)
            \item \textbf{II. Deuxième axe} (Idée force, souvent nuance/opposition)
            \item A. ...
            \item B. ...
            \item \textbf{III. Troisième axe} (Dépassement, synthèse, ouverture critique)
            \item A. ...
            \item B. ...
        \end{itemize}
        \item \textbf{Conclusion} : [Bilan synthétique répondant à la problématique] $\to$ [Ouverture (autre œuvre, actualité, question nouvelle)].
    \end{enumerate}
\end{methode}

\courssub{Rappels syntaxiques \& rhétoriques (Tableau récapitulatif)}

\begin{center}
\begin{tabularx}{\linewidth}{|l|X|X|}
\hline
\rowcolor{popblueL}
\textbf{Notion} & \textbf{Point clé / Règle} & \textbf{Exemple / Repère} \\
\hline
Phrase simple & 1 verbe conjugué = 1 proposition & \textit{« Lakunle enseigne. »} \\
\hline
Phrase composée & $\ge$ 2 propositions liées & \textit{Coordination : « Lakunle enseigne \textbf{mais} Baroka règne. »} \\
 & \textbf{Coordination} : \textit{mais, ou, et, donc, or, ni, car} & \textit{Subordination : « \textbf{Quand} Lakunle parle, Sidi écoute. »} \\
 & \textbf{Subordination} : subordonnée (relative, conjonctive, infinitive) & \textit{Juxtaposition : « Lakunle parle ; Sidi rit. »} \\
 & \textbf{Juxtaposition} : ; , : sans mot-outil & \\
\hline
Ponctuation forte & . ? ! ... : Fin de phrase, ton & \textit{« Le lion rugit ! »} \\
\hline
Ponctuation faible & , ; : ( ) " " -- / & \textit{« Le lion, roi des animaux, rugit. »} \\
\hline
Texte argumentatif & Thèse + Arguments + Exemples + Connecteurs & Connecteurs : \textit{donc, car, en effet, cependant, par conséquent} \\
\hline
Stratégie : Concession & Reconnaître une part de vérité adverse pour mieux la réfuter & \textit{« Certes, le divertissement existe, \textbf{mais} le théâtre instruit. »} \\
\hline
Stratégie : Réfutation & Démontrer la fausseté d'un contre-argument & \textit{« On dit que... Or, les faits prouvent que... »} \\
\hline
\end{tabularx}
\end{center}

\end{bilan}

\begin{aretenir}[Checklist avant le Bac]
    $\square$ Je sais identifier et commenter les paratextes de \textit{Le Lion et la Perle} (titre, auteur, genre, contexte).\par
    $\square$ Je maîtrise les 6 facteurs de la communication (Jakobson) et leurs fonctions associées.\par
    $\square$ Je distingue ponctuation forte/faible et j'explique leurs effets de sens (rythme, ton, syntaxe).\par
    $\square$ J'analyse la structure de la phrase : simple vs composée ; je nomme le type de liaison (coord, subord, juxtapos).\par
    $\square$ Je reconnais un texte argumentatif : je repère la thèse, les arguments, les exemples, les connecteurs logiques.\par
    $\square$ J'identifie et j'explique au moins 3 stratégies argumentatives (concession, réfutation, analogie, autorité...).\par
    $\square$ J'applique la méthode en 5 étapes pour analyser un sujet de dissertation et formuler une problématique pertinente.\par
    $\square$ Je sais faire un brainstorming efficace et classer mes idées en 2 ou 3 grands axes équilibrés.\par
    $\square$ Je construis un plan détaillé complet (Intro, Développement I/II/III avec A/B/C arguments+exemples, Conclusion).\par
    $\square$ Je vérifie la cohérence du plan : progression logique, réponse à la problématique, annonces respectées.\par
    $\square$ Je relis mon plan pour corriger la syntaxe, la ponctuation et l'orthographe (soin de la langue).\par
    $\square$ Je suis prêt(e) pour l'activité d'intégration : produire un plan détaillé en autonomie dans le temps imparti.
\end{aretenir}

\findecours

\end{document}
